\documentclass[12pt,a4paper]{report}
\usepackage[utf8]{inputenc}
\usepackage[T1]{fontenc}
\usepackage[french]{babel}
\usepackage{amsmath,amssymb,amsfonts}
\usepackage{amsthm}
\usepackage{geometry}
\geometry{a4paper, margin=1in}
\usepackage{tikz}
\usepackage{listings}
\usepackage{xcolor}
\usepackage{longtable}
\usepackage{booktabs}
\usepackage{calc}

\newcommand{\passthrough}[1]{#1}
\providecommand{\tightlist}{%
  \setlength{\itemsep}{0pt}\setlength{\parskip}{0pt}}

\lstset{
    language=Python,
    basicstyle=\ttfamily\small,
    keywordstyle=\color{blue},
    stringstyle=\color{red},
    commentstyle=\color{green!60!black},
    numbers=left,
    numberstyle=\tiny,
    stepnumber=1,
    numbersep=5pt,
    backgroundcolor=\color{white},
    showspaces=false,
    showstringspaces=false,
    showtabs=false,
    frame=single,
    tabsize=4,
    captionpos=b,
    breaklines=true,
    breakatwhitespace=false,
}

\title{Applications, injections, surjections, bijections et composition de fonctions}
\author{Charles EDOU NZE}
\date{}

\begin{document}
\maketitle
\tableofcontents

\subsection{1. Présentation du concept
clé}\label{pruxe9sentation-du-concept-cluxe9}

\begin{itemize}
\tightlist
\item
  \textbf{La Métaphore :} Imaginez un archer qui tire des flèches sur
  des cibles.

  \begin{itemize}
  \tightlist
  \item
    L'\textbf{Application}, c'est la règle du jeu : chaque flèche
    \emph{doit} être tirée et atteindre \emph{une seule} cible (on ne
    peut pas rater la cible, ni toucher deux cibles avec une seule
    flèche).
  \item
    L'\textbf{Injection}, c'est quand chaque cible ne reçoit \emph{au
    plus qu'une} flèche (pas de jaloux, personne ne partage sa cible).
  \item
    La \textbf{Surjection}, c'est quand \emph{toutes} les cibles
    reçoivent au moins une flèche (personne n'est oublié).
  \item
    La \textbf{Bijection}, c'est le cas parfait : chaque flèche a sa
    cible unique, et chaque cible a sa flèche unique. C'est comme un
    dictionnaire parfait entre deux mondes.
  \end{itemize}
\item
  \textbf{Le ``Pourquoi on a inventé ça'' :} Les mathématiques
  consistent souvent à transformer des objets en d'autres objets. Les
  fonctions sont les ``machines à transformer''. Comprendre si une
  machine est ``réversible'' (bijection) ou si elle ``perd de
  l'information'' (non injective) est crucial pour savoir si on peut
  revenir en arrière.
\item
  \textbf{Visualisation :} Imaginez deux ensembles de points reliés par
  des fils. Si les fils ne s'emmêlent jamais sur le même point
  d'arrivée, c'est injectif. Si chaque point d'arrivée a au moins un fil
  qui lui arrive dessus, c'est surjectif.
\end{itemize}

\subsection{2. Formalisation}\label{formalisation}

\subsubsection{A. Définitions
Formelles}\label{a.-duxe9finitions-formelles}

Soient \(E\) (ensemble de départ non vide) et \(F\) (ensemble d'arrivée
non vide) deux ensembles structurés. 1. \textbf{Application
(\(f : E \to F\)) :} Relation qui à chaque élément \(x \in E\) associe
un unique élément \(y \in F\), noté \(f(x)\). 2. \textbf{Injection :}
\(f\) est injective si
\(\forall x, x' \in E, f(x) = f(x') \Rightarrow x = x'\). 3.
\textbf{Surjection :} \(f\) est surjective si
\(\forall y \in F, \exists x \in E, y = f(x)\). 4. \textbf{Bijection :}
\(f\) est bijective si elle est à la fois injective et surjective. Cela
équivaut à dire que \(\forall y \in F, \exists ! x \in E, y = f(x)\). 5.
\textbf{Composition (\(g \circ f\)) :} Soient \(f : E \to F\) et
\(g : F \to G\). L'application \(g \circ f : E \to G\) est définie par
\((g \circ f)(x) = g(f(x))\).

\subsubsection{B. Théorèmes, Propositions \&
Lemmes}\label{b.-thuxe9oruxe8mes-propositions-lemmes}

\begin{quote}
\textbf{Proposition (Composition et Bijection) :} Soient \(f : E \to F\)
et \(g : F \to G\) deux applications bijectives. Alors
\(g \circ f : E \to G\) est bijective et sa réciproque est donnée par
\((g \circ f)^{-1} = f^{-1} \circ g^{-1}\).
\end{quote}

\begin{quote}
\textbf{Caractérisation de l'injectivité par la composition :} Soit un
ensemble de départ non vide \(E \neq \emptyset\). Une application
\(f : E \to F\) est injective si et seulement s'il existe une
application \(g : F \to E\) telle que \(g \circ f = \text{Id}_E\)
(rétraction à gauche).
\end{quote}

\subsection{3. Démonstrations}\label{duxe9monstrations}

\subsubsection{\texorpdfstring{Démonstration du Théorème Pivot :
Injection de la composée
\(g \circ f\)}{Démonstration du Théorème Pivot : Injection de la composée g \textbackslash circ f}}\label{duxe9monstration-du-thuxe9oruxe8me-pivot-injection-de-la-composuxe9e-g-circ-f}

Soient \(f : E \to F\) et \(g : F \to G\). Montrons que si \(f\) et
\(g\) sont injectives, alors \(g \circ f\) est injective.

\begin{enumerate}
\def\labelenumi{\arabic{enumi}.}
\item
  \textbf{Initialisation / Cadre :} Soient \(x, x' \in E\). Supposons
  que \((g \circ f)(x) = (g \circ f)(x')\). Notre but est de montrer que
  \(x = x'\).
\item
  \textbf{Étape 1 : Utilisation de la définition de la composée}
  \((g \circ f)(x) = (g \circ f)(x') \implies g(f(x)) = g(f(x'))\).
\item
  \textbf{Étape 2 : Utilisation de l'injectivité de \(g\)} Comme \(g\)
  est injective, pour tous \(y, y' \in F\),
  \(g(y) = g(y') \implies y = y'\). Ici, posons \(y = f(x)\) et
  \(y' = f(x')\). Puisque \(g(f(x)) = g(f(x'))\), alors par injectivité
  de \(g\), on a : \[f(x) = f(x')\]
\item
  \textbf{Étape 3 : Utilisation de l'injectivité de \(f\)} Comme \(f\)
  est injective, pour tous \(z, z' \in E\),
  \(f(z) = f(z') \implies z = z'\). Comme nous avons établi que
  \(f(x) = f(x')\), alors par injectivité de \(f\), on a : \[x = x'\]
\item
  \textbf{Conclusion :} Nous avons montré que
  \((g \circ f)(x) = (g \circ f)(x') \Rightarrow x = x'\). L'application
  \(g \circ f\) est donc injective. La démonstration est achevée.
\end{enumerate}

\subsection{4. Exercices d'Application}\label{exercices-dapplication}

\subsubsection{Exercice 1 : Application Directe (Fonction
réelle)}\label{exercice-1-application-directe-fonction-ruxe9elle}

\textbf{Énoncé :} Soit \(f : \mathbb{R} \to \mathbb{R}\) définie par
\(f(x) = 2x + 3\). Démontrer que \(f\) est bijective et déterminer sa
réciproque \(f^{-1}\). \textbf{Correction Détaillée :} * \emph{Analyse
de l'énoncé :} On doit montrer que pour tout \(y \in \mathbb{R}\),
l'équation \(f(x) = y\) possède une unique solution
\(x \in \mathbb{R}\). * \emph{Résolution pas-à-pas :} 1. Soit
\(y \in \mathbb{R}\). Résolvons \(2x + 3 = y\). 2. \(2x = y - 3\). 3.
\(x = \frac{y - 3}{2}\). 4. Comme \(y \in \mathbb{R}\), alors
\(\frac{y - 3}{2}\) est un nombre réel bien défini. 5. Il existe donc
une unique solution \(x = \frac{1}{2}y - \frac{3}{2}\) pour chaque
\(y\). 6. \(f\) est donc bijective. 7. Sa réciproque est
\(f^{-1} : \mathbb{R} \to \mathbb{R}\) définie par
\(f^{-1}(y) = \frac{1}{2}y - \frac{3}{2}\). * \emph{Conclusion :} La
fonction est une bijection de la droite réelle sur elle-même.

\subsubsection{Exercice 2 : Niveau Avancé (Composition et
Surjection)}\label{exercice-2-niveau-avancuxe9-composition-et-surjection}

\textbf{Énoncé :} Soient \(f : E \to F\) et \(g : F \to G\). Démontrer
que si \(g \circ f\) est surjective, alors \(g\) est surjective.
\textbf{Correction Détaillée :} * \emph{Analyse de l'énoncé :} On
utilise la définition de la surjectivité : atteindre tous les points de
l'ensemble d'arrivée. * \emph{Résolution pas-à-pas :} 1. Soit
\(z \in G\). Notre but est de trouver un antécédent \(y \in F\) par
\(g\) tel que \(g(y) = z\). 2. Comme \(g \circ f : E \to G\) est
surjective, par définition, il existe \(x \in E\) tel que
\((g \circ f)(x) = z\). 3. En utilisant la définition de la composition
: \(g(f(x)) = z\). 4. Posons \(y = f(x)\). Comme \(x \in E\) et
\(f : E \to F\), alors \(y\) est un élément de \(F\). 5. Nous avons
trouvé un élément \(y \in F\) tel que \(g(y) = z\). 6. Cela est vrai
pour tout \(z \in G\). * \emph{Conclusion :} L'application \(g\) est
donc surjective. (Note : On ne peut rien conclure sur la surjectivité de
\(f\) sans information supplémentaire).

\subsection{5. Application en Intelligence
Artificielle}\label{application-en-intelligence-artificielle}

\begin{itemize}
\tightlist
\item
  \textbf{Le Pont Théorique :} Les réseaux de neurones sont des
  \textbf{compositions massives} de fonctions élémentaires (couches).
\item
  \textbf{Exemple Concret :} Dans les \textbf{Normalizing Flows}
  (modèles génératifs d'IA), on cherche à apprendre une suite de
  transformations \(f_1, f_2, ..., f_n\) qui sont toutes des
  \textbf{bijections}. Pourquoi ? Parce qu'une bijection permet de
  transformer une distribution de probabilité simple (comme une
  Gaussienne) en une distribution complexe (une image), tout en étant
  capable de faire le calcul inverse (calculer la probabilité d'une
  image) de manière exacte. Si une couche n'est pas bijective, on perd
  de l'information et le modèle devient ``aveugle'' à certaines données.
\end{itemize}

\subsection{6. Liens Sémantiques}\label{liens-suxe9mantiques}

\begin{itemize}
\tightlist
\item
  \textbf{Concepts Précédents requis :} Jalon-4
\item
  \textbf{Concepts Futurs dépendants :} Jalon-6.md\textbar Jalon 6,
  Jalon 52 (Applications continues entre espaces topologiques et
  définition fine des homéomorphismes.), Jalon 111 (Applications
  différentiables entre variétés)
\end{itemize}
\chapter{Exercices d'application}
\section{Exercice 1}
\subsection{Exercice 1/10}\label{exercice-110}

\textbf{Niveau de difficulté :} \(\star \text{ sur } 5\)

\textbf{Sujet :} Reconnaître une application, injection, surjection sur
des ensembles finis (diagrammes sagittaux traduits en ensembles).

\begin{center}\rule{0.5\linewidth}{0.5pt}\end{center}

\subsubsection{Énoncé}\label{uxe9noncuxe9}

Soient les ensembles finis \(E = \{1, 2, 3\}\) et
\(F = \{a, b, c, d\}\). On définit la relation \(f\) de \(E\) vers \(F\)
par l'ensemble de couples : \[f = \{(1, a), (2, b), (3, a)\}\]

Déterminez, en justifiant rigoureusement chaque réponse : 1. Si \(f\)
est une application de \(E\) dans \(F\). 2. Si \(f\) est une injection
de \(E\) dans \(F\). 3. Si \(f\) est une surjection de \(E\) dans \(F\).

\begin{center}\rule{0.5\linewidth}{0.5pt}\end{center}

\subsubsection{Analyse de l'énoncé}\label{analyse-de-luxe9noncuxe9}

Cet exercice introductif vise à solidifier la compréhension des
définitions fondamentales d'une application, d'une injection et d'une
surjection sur des ensembles finis. La relation \(f\) est donnée
explicitement sous forme d'un ensemble de couples, ce qui est une
représentation directe d'un ``diagramme sagittal'' où les flèches sont
les couples \((x, f(x))\).

Nous allons rappeler les définitions clés :

\begin{itemize}
\item
  \textbf{Application (ou fonction) :} Une relation \(f\) de \(E\) vers
  \(F\) est une application si et seulement si tout élément \(x \in E\)
  possède une unique image \(y \in F\). Formellement :

  \begin{enumerate}
  \def\labelenumi{\arabic{enumi}.}
  \tightlist
  \item
    Pour tout \(x \in E\), il existe \(y \in F\) tel que
    \((x, y) \in f\). (Existence d'une image)
  \item
    Pour tout \(x \in E\), si \((x, y_1) \in f\) et \((x, y_2) \in f\),
    alors \(y_1 = y_2\). (Unicité de l'image) En d'autres termes, chaque
    élément de l'ensemble de départ \(E\) doit être le premier élément
    d'exactement un couple dans \(f\).
  \end{enumerate}
\item
  \textbf{Injection :} Une application \(f: E \to F\) est injective si
  et seulement si deux éléments distincts de \(E\) ont toujours des
  images distinctes dans \(F\). Formellement : Pour tout
  \(x_1, x_2 \in E\), si \(x_1 \neq x_2\), alors \(f(x_1) \neq f(x_2)\).
  Une formulation équivalente, souvent plus pratique pour la
  démonstration, est la contraposée : Pour tout \(x_1, x_2 \in E\), la
  condition formelle exige que si \(f(x_1) = f(x_2)\), alors
  nécessairement \(x_1 = x_2\), ce qui démontre l'absence d'information
  perdue. En d'autres termes, chaque élément de l'ensemble d'arrivée
  \(F\) possède au plus une pré-image dans \(E\).
\item
  \textbf{Surjection :} Une application \(f: E \to F\) est surjective si
  et seulement si tout élément \(y \in F\) possède au moins une
  pré-image \(x \in E\). Formellement : Pour tout \(y \in F\), il existe
  \(x \in E\) tel que \(f(x) = y\). Ceci est équivalent à dire que
  l'image de l'application, notée \(f(E) = \{f(x) \mid x \in E\}\), est
  égale à l'ensemble d'arrivée \(F\). En d'autres termes, chaque élément
  de l'ensemble d'arrivée \(F\) doit être le second élément d'au moins
  un couple dans \(f\).
\end{itemize}

Pour des ensembles finis, ces propriétés peuvent être vérifiées par
inspection directe de tous les éléments et de leurs images/pré-images.

\begin{center}\rule{0.5\linewidth}{0.5pt}\end{center}

\subsubsection{Correction exhaustive
pas-à-pas}\label{correction-exhaustive-pas-uxe0-pas}

Soient les ensembles \(E = \{1, 2, 3\}\) et \(F = \{a, b, c, d\}\). La
relation \(f\) est donnée par \(f = \{(1, a), (2, b), (3, a)\}\).

\paragraph{\texorpdfstring{1. \(f\) est-elle une application de \(E\)
dans \(F\)
?}{1. f est-elle une application de E dans F ?}}\label{f-est-elle-une-application-de-e-dans-f}

Pour que \(f\) soit une application, chaque élément de \(E\) doit avoir
une et une seule image dans \(F\). Nous allons vérifier cette condition
pour chaque élément de \(E\).

\begin{itemize}
\item
  \textbf{Pour l'élément \(1 \in E\) :} Nous cherchons les couples dans
  \(f\) dont le premier élément est \(1\). Nous trouvons le couple
  \((1, a)\). Il n'y a pas d'autre couple dans \(f\) dont le premier
  élément est \(1\). Donc, l'élément \(1\) a une unique image, qui est
  \(a\).
\item
  \textbf{Pour l'élément \(2 \in E\) :} Nous cherchons les couples dans
  \(f\) dont le premier élément est \(2\). Nous trouvons le couple
  \((2, b)\). Il n'y a pas d'autre couple dans \(f\) dont le premier
  élément est \(2\). Donc, l'élément \(2\) a une unique image, qui est
  \(b\).
\item
  \textbf{Pour l'élément \(3 \in E\) :} Nous cherchons les couples dans
  \(f\) dont le premier élément est \(3\). Nous trouvons le couple
  \((3, a)\). Il n'y a pas d'autre couple dans \(f\) dont le premier
  élément est \(3\). Donc, l'élément \(3\) a une unique image, qui est
  \(a\).
\end{itemize}

Puisque chaque élément de \(E\) (à savoir \(1, 2, 3\)) possède une et
une seule image dans \(F\), la relation \(f\) satisfait la définition
d'une application.

\textbf{Conclusion pour la question 1 :} Oui, \(f\) est une application
de \(E\) dans \(F\).

\paragraph{\texorpdfstring{2. \(f\) est-elle une injection de \(E\) dans
\(F\)
?}{2. f est-elle une injection de E dans F ?}}\label{f-est-elle-une-injection-de-e-dans-f}

Pour que \(f\) soit une injection, des éléments distincts de \(E\)
doivent avoir des images distinctes dans \(F\). Autrement dit, si
\(f(x_1) = f(x_2)\), alors \(x_1\) doit être égal à \(x_2\).

Nous allons examiner les images des éléments de \(E\) : * \(f(1) = a\) *
\(f(2) = b\) * \(f(3) = a\)

Nous observons que \(f(1) = a\) et \(f(3) = a\). Nous avons donc
\(f(1) = f(3)\). Cependant, les éléments de départ sont \(1\) et \(3\),
et \(1 \neq 3\). Puisque nous avons trouvé deux éléments distincts de
\(E\) (à savoir \(1\) et \(3\)) qui ont la même image dans \(F\) (à
savoir \(a\)), la condition d'injectivité n'est pas remplie.

\textbf{Conclusion pour la question 2 :} Non, \(f\) n'est pas une
injection de \(E\) dans \(F\).

\paragraph{\texorpdfstring{3. \(f\) est-elle une surjection de \(E\)
dans \(F\)
?}{3. f est-elle une surjection de E dans F ?}}\label{f-est-elle-une-surjection-de-e-dans-f}

Pour que \(f\) soit une surjection, chaque élément de \(F\) doit avoir
au moins une pré-image dans \(E\). Autrement dit, l'image de \(f\),
notée \(f(E)\), doit être égale à l'ensemble d'arrivée \(F\).

Calculons l'image de \(f\), \(f(E)\) : \(f(E) = \{f(x) \mid x \in E\}\)
\(f(E) = \{f(1), f(2), f(3)\}\) En utilisant les images que nous avons
identifiées : \(f(E) = \{a, b, a\}\) En éliminant les doublons pour
obtenir un ensemble : \(f(E) = \{a, b\}\)

Maintenant, comparons \(f(E)\) avec l'ensemble d'arrivée
\(F = \{a, b, c, d\}\). Nous avons \(f(E) = \{a, b\}\) et
\(F = \{a, b, c, d\}\). Pour que \(f\) soit surjective, il faut que
\(f(E) = F\). Cependant, nous constatons que \(c \in F\) mais
\(c \notin f(E)\). De même, \(d \in F\) mais \(d \notin f(E)\). Puisque
les éléments \(c\) et \(d\) de \(F\) n'ont aucune pré-image dans \(E\),
la condition de surjectivité n'est pas remplie.

\textbf{Conclusion pour la question 3 :} Non, \(f\) n'est pas une
surjection de \(E\) dans \(F\).

\begin{center}\rule{0.5\linewidth}{0.5pt}\end{center}

\subsubsection{Liens avec l'Intelligence
Artificielle}\label{liens-avec-lintelligence-artificielle}

Les concepts d'applications, d'injections et de surjections sont
fondamentaux en mathématiques discrètes et ont des répercussions
directes dans de nombreux domaines de l'Intelligence Artificielle, même
si leur application n'est pas toujours explicite au premier abord.

\begin{enumerate}
\def\labelenumi{\arabic{enumi}.}
\tightlist
\item
  \textbf{Représentation des Données et des Relations :}

  \begin{itemize}
  \tightlist
  \item
    En IA, les données sont souvent représentées sous forme de
    structures discrètes (graphes, bases de données relationnelles,
    ensembles de tuples). Une relation comme celle de l'exercice peut
    être vue comme une table dans une base de données où \(E\) est un
    ensemble de clés et \(F\) un ensemble de valeurs.
  \item
    Vérifier si une relation est une application, c'est s'assurer de
    l'intégrité des données : chaque ``clé'' (élément de \(E\)) doit
    correspondre à une ``valeur'' unique (élément de \(F\)). C'est une
    contrainte de base pour de nombreux systèmes de gestion de données
    utilisés en IA.
  \end{itemize}
\item
  \textbf{Fonctions de Hachage et Tables de Hachage :}

  \begin{itemize}
  \tightlist
  \item
    Les fonctions de hachage, essentielles pour l'efficacité des
    structures de données comme les tables de hachage (utilisées pour
    l'accès rapide aux données en IA), sont des applications.
    Idéalement, une fonction de hachage serait injective (pas de
    collisions), mais en pratique, ce n'est pas le cas pour des
    ensembles de départ plus grands que l'ensemble d'arrivée. L'étude
    des collisions (non-injectivité) est cruciale pour la performance
    des algorithmes.
  \end{itemize}
\item
  \textbf{Apprentissage Automatique (Machine Learning) :}

  \begin{itemize}
  \tightlist
  \item
    \textbf{Classification :} Un classifieur en apprentissage
    automatique est une application \(h: X \to Y\), où \(X\) est
    l'espace des caractéristiques d'entrée et \(Y\) est l'ensemble des
    étiquettes de classe. Pour un classifieur, il est essentiel que
    chaque entrée ait une unique prédiction (c'est une application). La
    surjectivité est souvent souhaitable (le modèle doit pouvoir prédire
    toutes les classes possibles), tandis que l'injectivité est rarement
    une propriété recherchée ou réalisable (des entrées différentes
    peuvent et doivent souvent être classées de la même manière).
  \item
    \textbf{Intégration de Caractéristiques (Feature Engineering) :} Les
    transformations de caractéristiques (par exemple, la normalisation,
    la vectorisation de texte) sont des applications. Comprendre si ces
    transformations sont injectives (préservent l'information unique de
    chaque donnée) ou non (introduisent des collisions ou une perte
    d'information) est crucial pour la conception de modèles
    performants.
  \item
    \textbf{Réseaux de Neurones :} Chaque couche d'un réseau de neurones
    applique une fonction (combinaison linéaire suivie d'une fonction
    d'activation) à son entrée. L'ensemble du réseau est une composition
    de telles applications. L'analyse de la surjectivité des couches de
    sortie est liée à la capacité du réseau à générer une diversité de
    sorties.
  \end{itemize}
\item
  \textbf{Logique et Raisonnement Automatisé :}

  \begin{itemize}
  \tightlist
  \item
    Dans la logique formelle et les systèmes de raisonnement automatisé,
    les relations et les fonctions sont utilisées pour modéliser des
    connaissances et des inférences. La vérification des propriétés de
    ces relations est une étape fondamentale pour garantir la cohérence
    et la validité des systèmes.
  \end{itemize}
\end{enumerate}

En somme, bien que l'exercice porte sur des ensembles finis très
simples, les principes sous-jacents de l'existence et de l'unicité des
images (application), de la distinction des pré-images (injection) et de
la couverture de l'espace d'arrivée (surjection) sont des concepts
mathématiques omniprésents qui structurent la conception, l'analyse et
l'optimisation de nombreux algorithmes et systèmes d'Intelligence
Artificielle.
\newpage
\section{Exercice 2}
\subsection{Exercice 2/10 : Injectivité et Surjectivité de Fonctions
Réelles
Simples}\label{exercice-210-injectivituxe9-et-surjectivituxe9-de-fonctions-ruxe9elles-simples}

\textbf{Niveau de difficulté :} \(\star \text{ (1 étoile sur 5)}\)

\subsubsection{Énoncé}\label{uxe9noncuxe9}

Soient les fonctions suivantes :

\begin{enumerate}
\def\labelenumi{\arabic{enumi}.}
\tightlist
\item
  \(f_1: \mathbb{R} \to \mathbb{R}\) définie par \(f_1(x) = ax+b\), où
  \(a, b \in \mathbb{R}\) et \(a \neq 0\).
\item
  \(f_2: \mathbb{R} \to \mathbb{R}\) définie par \(f_2(x) = x^2\).
\item
  \(f_3: \mathbb{R} \to \mathbb{R}\) définie par \(f_3(x) = e^x\).
\end{enumerate}

Pour chacune de ces fonctions, déterminez si elle est injective,
surjective, les deux (bijective), ou aucune des deux. Justifiez
rigoureusement vos réponses.

\subsubsection{Analyse de l'énoncé}\label{analyse-de-luxe9noncuxe9}

Cet exercice vise à consolider la compréhension des définitions
d'injectivité et de surjectivité pour des fonctions réelles
élémentaires. Il s'agit d'appliquer directement les définitions
formelles :

\begin{itemize}
\tightlist
\item
  \textbf{Injectivité :} Une fonction \(f: E \to F\) est injective si et
  seulement si pour tout \(x_1, x_2 \in E\), si \(f(x_1) = f(x_2)\),
  alors \(x_1 = x_2\). Autrement dit, chaque élément de l'ensemble
  d'arrivée est l'image d'au plus un élément de l'ensemble de départ.
\item
  \textbf{Surjectivité :} Une fonction \(f: E \to F\) est surjective si
  et seulement si pour tout \(y \in F\), il existe au moins un
  \(x \in E\) tel que \(f(x) = y\). Autrement dit, chaque élément de
  l'ensemble d'arrivée est l'image d'au moins un élément de l'ensemble
  de départ. L'ensemble image \(f(E)\) est égal à l'ensemble d'arrivée
  \(F\).
\end{itemize}

Pour prouver l'injectivité, on part de l'hypothèse \(f(x_1) = f(x_2)\)
pour des éléments arbitraires \(x_1, x_2\) du domaine \(E\), et on
cherche à déduire \(x_1 = x_2\). Pour prouver la non-injectivité, il
suffit de trouver un contre-exemple : deux éléments distincts
\(x_1 \neq x_2\) dans \(E\) qui ont la même image \(f(x_1) = f(x_2)\).

Pour prouver la surjectivité, on prend un \(y\) arbitraire dans
l'ensemble d'arrivée \(F\) et on cherche à résoudre l'équation
\(f(x) = y\) pour \(x \in E\). Si une solution \(x\) existe et
appartient à \(E\) pour tout \(y \in F\), la fonction est surjective.
Pour prouver la non-surjectivité, on cherche un contre-exemple : un
élément \(y \in F\) pour lequel l'équation \(f(x) = y\) n'a aucune
solution dans \(E\).

La difficulté est faible car les fonctions sont bien connues et les
manipulations algébriques sont simples. L'accent est mis sur la rigueur
de la démonstration et la bonne application des définitions.

\subsubsection{Correction exhaustive
pas-à-pas}\label{correction-exhaustive-pas-uxe0-pas}

Appliquons les définitions à chaque fonction.

\paragraph{\texorpdfstring{1. Fonction
\(f_1: \mathbb{R} \to \mathbb{R}\) définie par \(f_1(x) = ax+b\), avec
\(a, b \in \mathbb{R}\) et
\(a \neq 0\).}{1. Fonction f\_1: \textbackslash mathbb\{R\} \textbackslash to \textbackslash mathbb\{R\} définie par f\_1(x) = ax+b, avec a, b \textbackslash in \textbackslash mathbb\{R\} et a \textbackslash neq 0.}}\label{fonction-f_1-mathbbr-to-mathbbr-duxe9finie-par-f_1x-axb-avec-a-b-in-mathbbr-et-a-neq-0.}

\subparagraph{\texorpdfstring{Injectivité de
\(f_1\)}{Injectivité de f\_1}}\label{injectivituxe9-de-f_1}

Pour prouver l'injectivité, nous devons montrer que pour tout
\(x_1, x_2 \in \mathbb{R}\), si \(f_1(x_1) = f_1(x_2)\), alors
\(x_1 = x_2\).

Soient \(x_1, x_2 \in \mathbb{R}\) tels que \(f_1(x_1) = f_1(x_2)\). Par
définition de \(f_1\), l'égalité des images s'écrit :
\[ax_1 + b = ax_2 + b\] Pour isoler les termes en \(x\), nous
soustrayons \(b\) des deux côtés de l'équation :
\[ax_1 + b - b = ax_2 + b - b\] \[ax_1 = ax_2\] Puisque l'énoncé
spécifie que \(a \neq 0\), nous pouvons diviser les deux côtés de
l'équation par \(a\) : \[\frac{ax_1}{a} = \frac{ax_2}{a}\] \[x_1 = x_2\]
Nous avons démontré que si \(f_1(x_1) = f_1(x_2)\), alors \(x_1 = x_2\).
Par conséquent, la fonction \(f_1\) est injective.

\subparagraph{\texorpdfstring{Surjectivité de
\(f_1\)}{Surjectivité de f\_1}}\label{surjectivituxe9-de-f_1}

Pour prouver la surjectivité, nous devons montrer que pour tout
\(y \in \mathbb{R}\) (l'ensemble d'arrivée), il existe au moins un
\(x \in \mathbb{R}\) (l'ensemble de départ) tel que \(f_1(x) = y\).

Soit \(y \in \mathbb{R}\) un élément arbitraire de l'ensemble d'arrivée.
Nous cherchons à trouver un \(x \in \mathbb{R}\) tel que \(f_1(x) = y\).
Par définition de \(f_1\), cela revient à résoudre l'équation suivante
pour \(x\) : \[ax + b = y\] Pour isoler \(x\), nous soustrayons \(b\)
des deux côtés de l'équation : \[ax + b - b = y - b\] \[ax = y - b\]
Puisque \(a \neq 0\), nous pouvons diviser les deux côtés par \(a\) :
\[x = \frac{y - b}{a}\] Étant donné que \(y \in \mathbb{R}\),
\(b \in \mathbb{R}\) et \(a \in \mathbb{R}\) avec \(a \neq 0\), la
valeur \(x = \frac{y - b}{a}\) est un nombre réel bien défini. Ainsi,
pour tout \(y \in \mathbb{R}\), nous avons trouvé un
\(x \in \mathbb{R}\) tel que \(f_1(x) = y\). Par conséquent, la fonction
\(f_1\) est surjective.

\subparagraph{\texorpdfstring{Conclusion pour
\(f_1\)}{Conclusion pour f\_1}}\label{conclusion-pour-f_1}

Puisque \(f_1\) est à la fois injective et surjective, \(f_1\) est une
fonction bijective.

\paragraph{\texorpdfstring{2. Fonction
\(f_2: \mathbb{R} \to \mathbb{R}\) définie par
\(f_2(x) = x^2\).}{2. Fonction f\_2: \textbackslash mathbb\{R\} \textbackslash to \textbackslash mathbb\{R\} définie par f\_2(x) = x\^{}2.}}\label{fonction-f_2-mathbbr-to-mathbbr-duxe9finie-par-f_2x-x2.}

\subparagraph{\texorpdfstring{Injectivité de
\(f_2\)}{Injectivité de f\_2}}\label{injectivituxe9-de-f_2}

Pour prouver la non-injectivité, nous devons trouver deux éléments
distincts \(x_1, x_2 \in \mathbb{R}\) tels que \(f_2(x_1) = f_2(x_2)\).

Considérons les valeurs \(x_1 = 2\) et \(x_2 = -2\). Ces deux éléments
appartiennent à l'ensemble de départ \(\mathbb{R}\) et sont distincts :
\(x_1 \neq x_2\) car \(2 \neq -2\). Calculons leurs images sous la
fonction \(f_2\): \(f_2(x_1) = f_2(2) = (2)^2 = 4\).
\(f_2(x_2) = f_2(-2) = (-2)^2 = 4\). Nous observons que
\(f_2(2) = f_2(-2)\) alors que \(2 \neq -2\). Par conséquent, la
fonction \(f_2\) n'est pas injective.

\subparagraph{\texorpdfstring{Surjectivité de
\(f_2\)}{Surjectivité de f\_2}}\label{surjectivituxe9-de-f_2}

Pour prouver la non-surjectivité, nous devons trouver un élément
\(y \in \mathbb{R}\) (l'ensemble d'arrivée) pour lequel il n'existe
aucun \(x \in \mathbb{R}\) (l'ensemble de départ) tel que
\(f_2(x) = y\).

Soit \(y \in \mathbb{R}\) un élément arbitraire. Nous cherchons
\(x \in \mathbb{R}\) tel que \(f_2(x) = y\). Par définition de \(f_2\),
cela revient à résoudre l'équation : \[x^2 = y\] Nous savons que le
carré de tout nombre réel est toujours un nombre réel positif ou nul.
C'est-à-dire, pour tout \(x \in \mathbb{R}\), \(x^2 \ge 0\). Considérons
un \(y\) strictement négatif, par exemple \(y = -1\). L'équation devient
\(x^2 = -1\). Il n'existe aucun nombre réel \(x\) dont le carré est égal
à \(-1\). Par conséquent, pour \(y = -1 \in \mathbb{R}\), il n'existe
aucun \(x \in \mathbb{R}\) tel que \(f_2(x) = -1\). Par conséquent, la
fonction \(f_2\) n'est pas surjective.

\subparagraph{\texorpdfstring{Conclusion pour
\(f_2\)}{Conclusion pour f\_2}}\label{conclusion-pour-f_2}

Puisque \(f_2\) n'est ni injective ni surjective, \(f_2\) n'est pas une
bijection.

\paragraph{\texorpdfstring{3. Fonction
\(f_3: \mathbb{R} \to \mathbb{R}\) définie par
\(f_3(x) = e^x\).}{3. Fonction f\_3: \textbackslash mathbb\{R\} \textbackslash to \textbackslash mathbb\{R\} définie par f\_3(x) = e\^{}x.}}\label{fonction-f_3-mathbbr-to-mathbbr-duxe9finie-par-f_3x-ex.}

\subparagraph{\texorpdfstring{Injectivité de
\(f_3\)}{Injectivité de f\_3}}\label{injectivituxe9-de-f_3}

Pour prouver l'injectivité, nous devons montrer que pour tout
\(x_1, x_2 \in \mathbb{R}\), si \(f_3(x_1) = f_3(x_2)\), alors
\(x_1 = x_2\).

Soient \(x_1, x_2 \in \mathbb{R}\) tels que \(f_3(x_1) = f_3(x_2)\). Par
définition de \(f_3\), l'égalité des images s'écrit :
\[e^{x_1} = e^{x_2}\] Pour résoudre cette équation, nous pouvons
appliquer la fonction logarithme népérien (ln) aux deux côtés. La
fonction ln est bien définie pour les nombres strictement positifs, et
la fonction exponentielle \(e^x\) est toujours strictement positive pour
tout \(x \in \mathbb{R}\). \[\ln(e^{x_1}) = \ln(e^{x_2})\] En utilisant
la propriété fondamentale des logarithmes \(\ln(e^u) = u\) pour tout
\(u \in \mathbb{R}\) : \[x_1 = x_2\] Nous avons démontré que si
\(f_3(x_1) = f_3(x_2)\), alors \(x_1 = x_2\). Par conséquent, la
fonction \(f_3\) est injective.

\subparagraph{\texorpdfstring{Surjectivité de
\(f_3\)}{Surjectivité de f\_3}}\label{surjectivituxe9-de-f_3}

Pour prouver la non-surjectivité, nous devons trouver un élément
\(y \in \mathbb{R}\) (l'ensemble d'arrivée) pour lequel il n'existe
aucun \(x \in \mathbb{R}\) (l'ensemble de départ) tel que
\(f_3(x) = y\).

Soit \(y \in \mathbb{R}\) un élément arbitraire. Nous cherchons
\(x \in \mathbb{R}\) tel que \(f_3(x) = y\). Par définition de \(f_3\),
cela revient à résoudre l'équation : \[e^x = y\] Nous savons que la
fonction exponentielle \(e^x\) est toujours strictement positive pour
tout \(x \in \mathbb{R}\). C'est-à-dire, pour tout \(x \in \mathbb{R}\),
\(e^x > 0\). Considérons un \(y\) négatif ou nul, par exemple \(y = 0\).
L'équation devient \(e^x = 0\). Il n'existe aucun nombre réel \(x\) tel
que \(e^x = 0\). De même, si nous prenons \(y = -5\), l'équation devient
\(e^x = -5\). Il n'existe aucun nombre réel \(x\) tel que \(e^x = -5\),
car \(e^x\) doit toujours être positif. Par conséquent, pour tout
\(y \in \mathbb{R}\) tel que \(y \le 0\), il n'existe aucun
\(x \in \mathbb{R}\) tel que \(f_3(x) = y\). Par exemple, pour
\(y = -10 \in \mathbb{R}\), il n'existe aucun \(x \in \mathbb{R}\) tel
que \(f_3(x) = -10\). Par conséquent, la fonction \(f_3\) n'est pas
surjective.

\subparagraph{\texorpdfstring{Conclusion pour
\(f_3\)}{Conclusion pour f\_3}}\label{conclusion-pour-f_3}

Puisque \(f_3\) est injective mais n'est pas surjective, \(f_3\) n'est
pas une bijection.

\subsubsection{Liens avec l'Intelligence
Artificielle}\label{liens-avec-lintelligence-artificielle}

Les concepts d'injectivité et de surjectivité, bien que fondamentaux en
mathématiques pures, trouvent des échos significatifs dans plusieurs
domaines de l'Intelligence Artificielle, notamment en traitement de
données, en apprentissage automatique et en cryptographie.

\begin{enumerate}
\def\labelenumi{\arabic{enumi}.}
\item
  \textbf{Représentations et Encodages :} En IA, les données brutes sont
  souvent transformées en représentations numériques (encodages) pour
  être traitées par des algorithmes. Une fonction d'encodage idéale est
  injective. Si un encodage n'est pas injectif (comme \(f_2(x)=x^2\) où
  \(2\) et \(-2\) sont mappés à \(4\)), cela signifie que différentes
  entrées peuvent produire la même représentation. Cela peut entraîner
  une perte d'information ou une ambiguïté, rendant impossible de
  distinguer les entrées originales. Par exemple, dans l'encodage de
  mots (word embeddings), on souhaite que des mots sémantiquement
  distincts aient des représentations distinctes.
\item
  \textbf{Fonctions de Hachage :} Les fonctions de hachage sont
  utilisées pour mapper des données de taille arbitraire à des valeurs
  de taille fixe. Elles sont cruciales pour les tables de hachage, la
  vérification d'intégrité et la cryptographie. Idéalement, une fonction
  de hachage devrait être ``quasi-injective'' pour minimiser les
  collisions (où différentes entrées produisent la même sortie). Bien
  qu'il soit impossible d'avoir une fonction de hachage parfaitement
  injective pour des domaines potentiellement infinis vers des
  codomaines finis, l'objectif est de s'en approcher le plus possible
  pour garantir l'unicité des identifiants ou la robustesse des
  signatures numériques.
\item
  \textbf{Compression de Données :} Les algorithmes de compression de
  données visent à réduire la taille des données. Une compression ``sans
  perte'' est une fonction qui est bijective de l'ensemble des données
  originales vers l'ensemble des données compressées. Cela garantit que
  les données originales peuvent être entièrement reconstruites à partir
  des données compressées, ce qui implique que la fonction de
  compression doit être injective (pas de perte d'information) et
  surjective (toutes les données compressées peuvent être
  décompressées).
\item
  \textbf{Modèles Génératifs (GANs, VAEs) :} Dans les modèles
  génératifs, on cherche à générer de nouvelles données (images, textes,
  etc.) à partir d'un espace latent de variables aléatoires. La
  ``qualité'' de la génération peut être liée à la capacité de la
  fonction de génération (le générateur) à couvrir l'espace des données
  réelles (surjectivité) et à produire des échantillons variés et non
  redondants (injectivité dans l'espace latent). Un générateur qui n'est
  pas suffisamment surjectif pourrait souffrir de ``mode collapse'', où
  il ne génère qu'une petite variété d'échantillons.
\item
  \textbf{Fonctions d'Activation dans les Réseaux de Neurones :} Bien
  que l'injectivité et la surjectivité ne soient pas des critères
  directs pour les fonctions d'activation (comme ReLU, Sigmoid, Tanh),
  la compréhension de leur comportement sur leur domaine et codomaine
  est essentielle. Par exemple, la fonction sigmoïde mappe
  \(\mathbb{R}\) à \((0,1)\), elle est injective mais pas surjective sur
  \(\mathbb{R}\). Cela signifie qu'elle compresse l'information dans un
  intervalle borné, ce qui peut être utile pour la normalisation mais
  peut aussi entraîner une perte de gradient pour des entrées très
  grandes ou très petites.
\end{enumerate}

En somme, la compréhension des propriétés d'injectivité et de
surjectivité permet d'évaluer la fidélité, la réversibilité et la
couverture des transformations de données, des concepts cruciaux pour la
conception et l'analyse des algorithmes d'Intelligence Artificielle.
\newpage
\section{Exercice 3}
\section{Exercice 3/10 - Jalon 5 : Composition et propriétés des
fonctions}\label{exercice-310---jalon-5-composition-et-propriuxe9tuxe9s-des-fonctions}

\textbf{Niveau de difficulté :} \(\star\star\)

\begin{center}\rule{0.5\linewidth}{0.5pt}\end{center}

\subsubsection{Énoncé}\label{uxe9noncuxe9}

Soient les fonctions \(f\) et \(g\) définies comme suit : *
\(f: \mathbb{R} \to \mathbb{R}\), donnée par \(f(x) = x-1\). *
\(g: \mathbb{R} \to \mathbb{R}\), donnée par \(g(x) = x^2\).

On considère la fonction composée \(h = g \circ f\).

\begin{enumerate}
\def\labelenumi{\arabic{enumi}.}
\tightlist
\item
  Déterminer l'expression explicite de la fonction \(h(x)\). Préciser
  rigoureusement son ensemble de départ et son ensemble d'arrivée.
\item
  Étudier l'injectivité de la fonction \(h\). Justifier votre réponse de
  manière exhaustive.
\item
  Étudier la surjectivité de la fonction \(h\). Justifier votre réponse
  de manière exhaustive.
\end{enumerate}

\begin{center}\rule{0.5\linewidth}{0.5pt}\end{center}

\subsubsection{Analyse de l'énoncé}\label{analyse-de-luxe9noncuxe9}

Cet exercice vise à évaluer la compréhension des concepts fondamentaux
de composition de fonctions, d'injectivité et de surjectivité.

\begin{enumerate}
\def\labelenumi{\arabic{enumi}.}
\tightlist
\item
  \textbf{Composition de fonctions (\(h = g \circ f\)) :} Il s'agit
  d'appliquer la fonction \(f\) en premier, puis la fonction \(g\) au
  résultat de \(f\). La définition formelle \((g \circ f)(x) = g(f(x))\)
  sera utilisée. Il est crucial de vérifier que les ensembles de
  définition et d'arrivée sont compatibles pour que la composition soit
  bien définie. Ici, l'ensemble d'arrivée de \(f\) est \(\mathbb{R}\),
  qui est aussi l'ensemble de départ de \(g\), donc la composition est
  bien définie sur \(\mathbb{R}\). L'ensemble de départ de \(h\) sera
  celui de \(f\), et l'ensemble d'arrivée de \(h\) sera celui de \(g\).
\item
  \textbf{Injectivité :} Une fonction \(\phi: E \to F\) est injective si
  et seulement si pour tout \(x_1, x_2 \in E\), l'égalité
  \(\phi(x_1) = \phi(x_2)\) implique \(x_1 = x_2\). Pour prouver qu'une
  fonction n'est \emph{pas} injective, il suffit de trouver un
  contre-exemple, c'est-à-dire deux éléments distincts \(x_1 \neq x_2\)
  dans l'ensemble de départ qui ont la même image par \(\phi\).
\item
  \textbf{Surjectivité :} Une fonction \(\phi: E \to F\) est surjective
  si et seulement si pour tout \(y \in F\) (l'ensemble d'arrivée), il
  existe au moins un \(x \in E\) (l'ensemble de départ) tel que
  \(\phi(x) = y\). Pour prouver qu'une fonction n'est \emph{pas}
  surjective, il suffit de trouver un élément \(y_0\) dans l'ensemble
  d'arrivée \(F\) pour lequel il n'existe aucun \(x\) dans l'ensemble de
  départ \(E\) tel que \(\phi(x) = y_0\).
\end{enumerate}

Les fonctions \(f\) et \(g\) sont des fonctions polynomiales simples. La
fonction \(f(x) = x-1\) est une translation, qui est bijective. La
fonction \(g(x) = x^2\) est une fonction quadratique, qui n'est ni
injective ni surjective sur \(\mathbb{R}\). La composition de ces deux
fonctions devrait hériter des propriétés de la fonction quadratique en
termes d'injectivité et de surjectivité.

\begin{center}\rule{0.5\linewidth}{0.5pt}\end{center}

\subsubsection{Correction exhaustive
pas-à-pas}\label{correction-exhaustive-pas-uxe0-pas}

\paragraph{\texorpdfstring{Question 1 : Déterminer l'expression
explicite de \(h(x)\) et préciser ses
ensembles.}{Question 1 : Déterminer l'expression explicite de h(x) et préciser ses ensembles.}}\label{question-1-duxe9terminer-lexpression-explicite-de-hx-et-pruxe9ciser-ses-ensembles.}

Soient les fonctions : * \(f: \mathbb{R} \to \mathbb{R}\), définie par
\(f(x) = x-1\). * \(g: \mathbb{R} \to \mathbb{R}\), définie par
\(g(x) = x^2\).

Nous cherchons la fonction composée \(h = g \circ f\).

\textbf{Étape 1 : Vérification de la compatibilité des ensembles.}
L'ensemble d'arrivée de \(f\) est \(\mathbb{R}\). L'ensemble de départ
de \(g\) est \(\mathbb{R}\). Puisque l'ensemble d'arrivée de \(f\) est
égal à l'ensemble de départ de \(g\) (\(\mathbb{R} = \mathbb{R}\)), la
composition \(g \circ f\) est bien définie.

\textbf{Étape 2 : Détermination des ensembles de départ et d'arrivée de
\(h\).} L'ensemble de départ de \(h = g \circ f\) est l'ensemble de
départ de \(f\), qui est \(\mathbb{R}\). L'ensemble d'arrivée de
\(h = g \circ f\) est l'ensemble d'arrivée de \(g\), qui est
\(\mathbb{R}\). Ainsi, \(h: \mathbb{R} \to \mathbb{R}\).

\textbf{Étape 3 : Calcul de l'expression explicite de \(h(x)\).} Par
définition de la composition de fonctions, pour tout
\(x \in \mathbb{R}\) : \[h(x) = (g \circ f)(x) = g(f(x))\] Nous
substituons l'expression de \(f(x)\) dans \(g\): \[h(x) = g(x-1)\]
Maintenant, nous appliquons la définition de \(g\), qui est
\(g(u) = u^2\) pour tout \(u \in \mathbb{R}\). Ici, \(u = x-1\).
\[h(x) = (x-1)^2\] Pour une expression développée, nous pouvons utiliser
l'identité remarquable \((a-b)^2 = a^2 - 2ab + b^2\):
\[h(x) = x^2 - 2 \cdot x \cdot 1 + 1^2\] \[h(x) = x^2 - 2x + 1\]

\textbf{Conclusion de la Question 1 :} La fonction \(h\) est définie par
\(h: \mathbb{R} \to \mathbb{R}\), et son expression explicite est
\(h(x) = (x-1)^2\) (ou \(h(x) = x^2 - 2x + 1\)).

\paragraph{\texorpdfstring{Question 2 : Étudier l'injectivité de la
fonction
\(h\).}{Question 2 : Étudier l'injectivité de la fonction h.}}\label{question-2-uxe9tudier-linjectivituxe9-de-la-fonction-h.}

Pour étudier l'injectivité de \(h: \mathbb{R} \to \mathbb{R}\) définie
par \(h(x) = (x-1)^2\), nous devons vérifier si pour tout
\(x_1, x_2 \in \mathbb{R}\), l'égalité \(h(x_1) = h(x_2)\) implique
\(x_1 = x_2\).

\textbf{Étape 1 : Poser l'hypothèse d'égalité des images.} Soient
\(x_1, x_2 \in \mathbb{R}\) tels que \(h(x_1) = h(x_2)\). Ceci signifie
: \[(x_1-1)^2 = (x_2-1)^2\]

\textbf{Étape 2 : Résoudre l'équation pour \(x_1\) et \(x_2\).}
L'équation \(A^2 = B^2\) est équivalente à \(A = B\) ou \(A = -B\).
Donc, nous avons deux cas possibles : * \textbf{Cas 1 :}
\(x_1-1 = x_2-1\) En ajoutant \(1\) aux deux membres de l'équation, nous
obtenons : \(x_1 = x_2\) Ce cas correspond à la condition d'injectivité.

\begin{itemize}
\tightlist
\item
  \textbf{Cas 2 :} \(x_1-1 = -(x_2-1)\) \(x_1-1 = -x_2+1\) En ajoutant
  \(1\) aux deux membres de l'équation : \(x_1 = -x_2+2\)
\end{itemize}

\textbf{Étape 3 : Chercher un contre-exemple (si la fonction n'est pas
injective).} Le Cas 2 montre qu'il est possible d'avoir
\(h(x_1) = h(x_2)\) sans que \(x_1 = x_2\). Pour prouver que \(h\) n'est
pas injective, il suffit de trouver un exemple concret de
\(x_1 \neq x_2\) tels que \(h(x_1) = h(x_2)\). Choisissons une valeur
pour \(x_1\), par exemple \(x_1 = 0\). Alors
\(h(0) = (0-1)^2 = (-1)^2 = 1\). Nous cherchons un \(x_2 \neq 0\) tel
que \(h(x_2) = 1\). \((x_2-1)^2 = 1\) Ceci implique \(x_2-1 = 1\) ou
\(x_2-1 = -1\). Si \(x_2-1 = 1\), alors \(x_2 = 1+1 = 2\). Si
\(x_2-1 = -1\), alors \(x_2 = -1+1 = 0\). Nous avons trouvé \(x_2 = 2\).
Vérifions : \(h(2) = (2-1)^2 = 1^2 = 1\). Nous avons donc \(h(0) = 1\)
et \(h(2) = 1\). Puisque \(0 \neq 2\) mais \(h(0) = h(2)\), la fonction
\(h\) n'est pas injective.

\textbf{Conclusion de la Question 2 :} La fonction \(h\) n'est pas
injective. Par exemple, \(h(0) = 1\) et \(h(2) = 1\), alors que
\(0 \neq 2\).

\paragraph{\texorpdfstring{Question 3 : Étudier la surjectivité de la
fonction
\(h\).}{Question 3 : Étudier la surjectivité de la fonction h.}}\label{question-3-uxe9tudier-la-surjectivituxe9-de-la-fonction-h.}

Pour étudier la surjectivité de \(h: \mathbb{R} \to \mathbb{R}\) définie
par \(h(x) = (x-1)^2\), nous devons vérifier si pour tout
\(y \in \mathbb{R}\) (l'ensemble d'arrivée), il existe au moins un
\(x \in \mathbb{R}\) (l'ensemble de départ) tel que \(h(x) = y\).

\textbf{Étape 1 : Poser l'équation \(h(x) = y\).} Nous cherchons à
savoir si, pour tout \(y \in \mathbb{R}\), l'équation \((x-1)^2 = y\)
admet au moins une solution \(x \in \mathbb{R}\).

\textbf{Étape 2 : Analyser l'équation \((x-1)^2 = y\).} Le terme
\((x-1)^2\) est le carré d'un nombre réel \((x-1)\). Par les propriétés
des nombres réels, le carré de tout nombre réel est toujours supérieur
ou égal à zéro. Donc, pour tout \(x \in \mathbb{R}\), nous avons
\((x-1)^2 \ge 0\). Ceci signifie que l'image de la fonction \(h\)
(l'ensemble des valeurs que \(h(x)\) peut prendre) est l'intervalle
\([0, +\infty[\).

\textbf{Étape 3 : Comparer l'image avec l'ensemble d'arrivée.}
L'ensemble d'arrivée de \(h\) est \(\mathbb{R}\). L'image de \(h\) est
\([0, +\infty[\). Puisque \([0, +\infty[ \neq \mathbb{R}\), la fonction
\(h\) n'est pas surjective.

\textbf{Étape 4 : Fournir un contre-exemple (si la fonction n'est pas
surjective).} Pour prouver que \(h\) n'est pas surjective, il suffit de
trouver un \(y_0 \in \mathbb{R}\) tel qu'il n'existe aucun
\(x \in \mathbb{R}\) vérifiant \(h(x) = y_0\). Choisissons un \(y_0\)
qui n'est pas dans l'image de \(h\), c'est-à-dire un nombre réel
strictement négatif. Par exemple, soit \(y_0 = -1\). Nous cherchons un
\(x \in \mathbb{R}\) tel que \(h(x) = -1\). \[(x-1)^2 = -1\] Comme nous
l'avons établi, le carré de tout nombre réel est non-négatif. Il est
donc impossible qu'un carré soit égal à \(-1\) dans l'ensemble des
nombres réels. L'équation \((x-1)^2 = -1\) n'admet aucune solution
réelle \(x\). Par conséquent, il n'existe aucun \(x \in \mathbb{R}\) tel
que \(h(x) = -1\).

\textbf{Conclusion de la Question 3 :} La fonction \(h\) n'est pas
surjective. Par exemple, la valeur \(-1\) appartient à l'ensemble
d'arrivée \(\mathbb{R}\), mais il n'existe aucun \(x \in \mathbb{R}\)
tel que \(h(x) = -1\).

\begin{center}\rule{0.5\linewidth}{0.5pt}\end{center}

\subsubsection{Liens avec l'Intelligence
Artificielle}\label{liens-avec-lintelligence-artificielle}

Les concepts d'applications, d'injectivité, de surjectivité et de
composition sont absolument fondamentaux en Intelligence Artificielle,
en particulier dans le domaine de l'apprentissage profond (Deep
Learning).

\begin{enumerate}
\def\labelenumi{\arabic{enumi}.}
\item
  \textbf{Composition de fonctions et Réseaux de Neurones :} Un réseau
  de neurones profond est, par essence, une composition de fonctions.
  Chaque ``couche'' du réseau applique une transformation (souvent une
  combinaison linéaire suivie d'une fonction d'activation non linéaire)
  aux données d'entrée. Si \(L_1, L_2, \dots, L_k\) sont les fonctions
  représentant les couches successives d'un réseau, alors la fonction
  globale apprise par le réseau est
  \(L_k \circ L_{k-1} \circ \dots \circ L_1\). Comprendre la composition
  est donc essentiel pour saisir comment les informations sont
  transformées et traitées à travers les différentes étapes d'un modèle
  d'IA.
\item
  \textbf{Injectivité et Surjectivité dans les Modèles Génératifs :}

  \begin{itemize}
  \tightlist
  \item
    \textbf{Modèles Génératifs (GANs, VAEs) :} Ces modèles apprennent à
    générer de nouvelles données (images, textes, etc.) à partir d'un
    ``espace latent'' (un espace de vecteurs de faible dimension). La
    fonction qui mappe l'espace latent à l'espace des données est une
    application complexe.
  \item
    \textbf{Injectivité :} Si cette application n'est pas injective,
    cela signifie que différents points dans l'espace latent peuvent
    produire la même sortie dans l'espace des données. C'est un problème
    connu sous le nom de ``mode collapse'' dans les GANs, où le
    générateur ne produit qu'une variété limitée d'échantillons,
    ignorant d'autres ``modes'' (types de données) de la distribution
    cible. Une perte d'injectivité implique une perte d'information lors
    de la transformation.
  \item
    \textbf{Surjectivité :} Si l'application de l'espace latent à
    l'espace des données n'est pas surjective (par rapport à la
    distribution de données réelles), cela signifie que le modèle ne
    peut pas générer certains types de données qui existent dans la
    distribution réelle. Le modèle ne couvre pas l'intégralité de
    l'espace des données, ce qui limite sa capacité à produire des
    échantillons diversifiés et réalistes.
  \end{itemize}
\item
  \textbf{Réseaux de Neurones Inversibles (Invertible Neural Networks -
  INN) :} Certains architectures de réseaux de neurones sont
  spécifiquement conçues pour être bijectives (à la fois injectives et
  surjectives). Ces INN sont utilisés dans des domaines comme
  l'estimation de densité, la compression de données sans perte, ou la
  génération d'images. La bijectivité garantit qu'il n'y a pas de perte
  d'information lors du passage à travers le réseau et que chaque sortie
  correspond à une entrée unique, et vice-versa. Cela permet, par
  exemple, de calculer la vraisemblance exacte des données, ce qui est
  crucial pour certaines tâches.
\item
  \textbf{Représentation et Réduction de Dimension :} Lorsque l'on
  utilise des fonctions pour transformer des données brutes en
  ``features'' (caractéristiques) plus significatives, les propriétés
  d'injectivité et de surjectivité sont importantes. Une transformation
  non injective peut entraîner une perte d'information, où des données
  d'entrée distinctes sont représentées de manière identique, rendant
  leur discrimination impossible par la suite. La surjectivité est liée
  à la capacité de la représentation à capturer toutes les variations
  pertinentes des données originales.
\end{enumerate}

En somme, une compréhension approfondie de ces concepts mathématiques
permet aux chercheurs et ingénieurs en IA de concevoir, d'analyser et de
diagnostiquer les comportements de modèles complexes, en particulier
lorsqu'il s'agit de transformations de données et de génération de
contenu.
\newpage
\section{Exercice 4}
\section{Exercice 4/10 : Jalon 5 - Applications, injections,
surjections, bijections et composition de
fonctions}\label{exercice-410-jalon-5---applications-injections-surjections-bijections-et-composition-de-fonctions}

\textbf{Niveau de difficulté :} \(\star\)\(\star\)$\circ$$\circ$$\circ$

\begin{center}\rule{0.5\linewidth}{0.5pt}\end{center}

\subsection{Énoncé}\label{uxe9noncuxe9}

Soient \(E\), \(F\) et \(G\) trois ensembles non vides. Soient
\(f: F \to G\) et \(g: E \to F\) deux applications.

Démontrer l'implication suivante :
\[ (f \circ g \text{ est injective}) \implies (g \text{ est injective}) \]

\begin{center}\rule{0.5\linewidth}{0.5pt}\end{center}

\subsection{Analyse de l'énoncé}\label{analyse-de-luxe9noncuxe9}

Cet exercice nous demande de prouver une propriété fondamentale de la
composition d'applications concernant l'injectivité. Nous sommes dans un
cadre purement abstrait, où les ensembles \(E, F, G\) et les
applications \(f, g\) ne sont pas spécifiés au-delà de leurs types.

\textbf{Rappel des définitions clés :}

\begin{enumerate}
\def\labelenumi{\arabic{enumi}.}
\tightlist
\item
  \textbf{Application \(h: A \to B\) :} Pour tout élément \(a \in A\),
  il existe un unique élément \(b \in B\) tel que \(h(a) = b\).
\item
  \textbf{Application \(h: A \to B\) injective :} Pour tous
  \(a_1, a_2 \in A\), si \(h(a_1) = h(a_2)\), alors \(a_1 = a_2\).
  Autrement dit, des éléments distincts de l'ensemble de départ sont
  toujours envoyés sur des éléments distincts de l'ensemble d'arrivée.
\item
  \textbf{Application composée \(f \circ g : E \to G\) :} Pour tout
  \(x \in E\), \((f \circ g)(x) = f(g(x))\). L'application \(g\) est
  appliquée en premier, puis \(f\) est appliquée au résultat de \(g\).
\end{enumerate}

\textbf{Ce que l'on nous donne (hypothèse) :} L'application
\(f \circ g: E \to G\) est injective. Cela signifie que pour tous
\(x_1, x_2 \in E\), si \((f \circ g)(x_1) = (f \circ g)(x_2)\), alors
\(x_1 = x_2\).

\textbf{Ce que l'on doit démontrer (conclusion) :} L'application
\(g: E \to F\) est injective. Cela signifie que pour tous
\(x_1, x_2 \in E\), si \(g(x_1) = g(x_2)\), alors \(x_1 = x_2\).

\textbf{Stratégie de démonstration :} Pour prouver que \(g\) est
injective, nous devons partir de l'hypothèse \(g(x_1) = g(x_2)\) pour
des éléments arbitraires \(x_1, x_2 \in E\), et en déduire que
\(x_1 = x_2\). L'information cruciale à utiliser sera l'injectivité de
\(f \circ g\).

\begin{center}\rule{0.5\linewidth}{0.5pt}\end{center}

\subsection{Correction exhaustive
``pas-à-pas''}\label{correction-exhaustive-pas-uxe0-pas}

Soient \(E\), \(F\) et \(G\) trois ensembles non vides. Soient
\(f: F \to G\) et \(g: E \to F\) deux applications.

Nous voulons démontrer que si \(f \circ g\) est injective, alors \(g\)
est injective.

\textbf{Étape 1 : Rappel des définitions formelles.}

\begin{itemize}
\tightlist
\item
  L'application composée \(f \circ g\) est définie pour tout \(x \in E\)
  par \((f \circ g)(x) = f(g(x))\). Son ensemble de départ est \(E\) et
  son ensemble d'arrivée est \(G\).
\item
  L'hypothèse que \(f \circ g\) est injective signifie que : Pour tous
  \(x_1, x_2 \in E\), si \((f \circ g)(x_1) = (f \circ g)(x_2)\), alors
  \(x_1 = x_2\).
\item
  La conclusion que nous voulons atteindre est que \(g\) est injective,
  ce qui signifie que : Pour tous \(x_1, x_2 \in E\), si
  \(g(x_1) = g(x_2)\), alors \(x_1 = x_2\).
\end{itemize}

\textbf{Étape 2 : Début de la démonstration de l'injectivité de \(g\).}

Pour démontrer que \(g\) est injective, nous devons prendre deux
éléments arbitraires \(x_1\) et \(x_2\) dans l'ensemble de départ de
\(g\), qui est \(E\), et supposer que leurs images par \(g\) sont
égales. Ensuite, nous devrons montrer que cela implique que \(x_1\) et
\(x_2\) sont en fait le même élément.

Soient \(x_1 \in E\) et \(x_2 \in E\) deux éléments quelconques.
Supposons que \(g(x_1) = g(x_2)\).

\textbf{Étape 3 : Application de \(f\) et utilisation de l'hypothèse
d'injectivité de \(f \circ g\).}

Puisque \(g(x_1)\) et \(g(x_2)\) sont des éléments de l'ensemble \(F\)
(l'ensemble d'arrivée de \(g\) et l'ensemble de départ de \(f\)), et que
nous avons supposé qu'ils sont égaux, nous pouvons appliquer
l'application \(f\) à ces deux images égales.

Appliquons \(f\) aux deux membres de l'égalité \(g(x_1) = g(x_2)\) :
\[ f(g(x_1)) = f(g(x_2)) \]

Par définition de la composition d'applications, nous savons que
\(f(g(x_1)) = (f \circ g)(x_1)\) et \(f(g(x_2)) = (f \circ g)(x_2)\).
Donc, l'égalité précédente peut être réécrite comme :
\[ (f \circ g)(x_1) = (f \circ g)(x_2) \]

À ce stade, nous avons établi que si \(g(x_1) = g(x_2)\), alors
\((f \circ g)(x_1) = (f \circ g)(x_2)\). Or, l'hypothèse de l'énoncé est
que l'application \(f \circ g\) est injective. Par définition de
l'injectivité de \(f \circ g\), si
\((f \circ g)(x_1) = (f \circ g)(x_2)\), alors cela implique que les
éléments de départ \(x_1\) et \(x_2\) doivent être égaux. Donc, en
utilisant l'injectivité de \(f \circ g\), nous déduisons :
\[ x_1 = x_2 \]

\textbf{Étape 4 : Conclusion.}

Nous avons commencé par supposer que \(g(x_1) = g(x_2)\) pour des
\(x_1, x_2 \in E\) arbitraires, et nous avons logiquement déduit que
\(x_1 = x_2\). Ceci correspond précisément à la définition de
l'injectivité de l'application \(g\).

Par conséquent, nous avons démontré que si \(f \circ g\) est injective,
alors \(g\) est injective.

\begin{center}\rule{0.5\linewidth}{0.5pt}\end{center}

\subsection{Liens avec l'Intelligence
Artificielle}\label{liens-avec-lintelligence-artificielle}

Le concept d'injectivité et de composition de fonctions est fondamental
en Intelligence Artificielle, en particulier dans le domaine de
l'apprentissage profond (Deep Learning).

\begin{enumerate}
\def\labelenumi{\arabic{enumi}.}
\item
  \textbf{Réseaux de Neurones comme Compositions de Fonctions :} Un
  réseau de neurones profond est intrinsèquement une composition de
  fonctions. Chaque couche du réseau peut être vue comme une application
  \(f_i: \mathbb{R}^{d_i} \to \mathbb{R}^{d_{i+1}}\), où \(d_i\) est la
  dimension de l'espace d'entrée de la couche \(i\). Un réseau de \(L\)
  couches est alors la composition
  \(F = f_L \circ f_{L-1} \circ \dots \circ f_1\). L'injectivité de la
  composition \(f \circ g\) implique l'injectivité de \(g\) est une
  propriété cruciale ici. Si la première couche (\(g\)) d'un réseau de
  neurones n'est pas injective, elle ``fusionne'' des informations
  distinctes dès le début. Aucune couche subséquente (\(f\)) ne pourra
  ``dé-fusionner'' ces informations, car l'information perdue est
  irrécupérable.
\item
  \textbf{Préservation de l'Information et Représentations Latentes :}
  Dans des architectures comme les auto-encodeurs, l'objectif est
  d'apprendre une représentation latente (un encodage) des données
  d'entrée. L'encodeur \(g: \text{Données} \to \text{Latent}\) cherche à
  compresser l'information. Le décodeur
  \(f: \text{Latent} \to \text{Données}\) tente de reconstruire l'entrée
  à partir de la représentation latente. Si l'encodeur \(g\) n'est pas
  injectif, cela signifie que différentes données d'entrée peuvent être
  mappées à la même représentation latente. Dans ce cas, même un
  décodeur parfait \(f\) ne pourrait pas reconstruire l'entrée originale
  de manière unique, car il y aurait une ambiguïté sur quelle entrée
  originale correspond à cette représentation latente. La propriété
  démontrée ici signifie que si l'auto-encodeur complet (\(f \circ g\))
  est capable de reconstruire l'entrée de manière unique (ce qui est une
  forme d'injectivité si l'on considère l'identité comme la fonction
  cible), alors l'encodeur \(g\) lui-même doit avoir été injectif pour
  ne pas perdre d'information essentielle.
\item
  \textbf{Fonctions de Hachage et Collisions :} Bien que les fonctions
  de hachage ne soient généralement pas injectives (elles sont conçues
  pour mapper un grand espace d'entrée à un espace de sortie plus petit,
  entraînant des ``collisions''), le concept est pertinent. Si une
  fonction de hachage \(g\) est utilisée comme première étape d'un
  processus, et qu'elle n'est pas injective, alors toute fonction \(f\)
  appliquée ensuite à la sortie de \(g\) ne pourra pas distinguer les
  entrées originales qui ont produit la même valeur de hachage.
  L'injectivité de \(g\) est une condition nécessaire pour la
  préservation de la distinction des entrées à travers une composition.
\end{enumerate}

En résumé, cette propriété mathématique souligne l'importance de
l'injectivité des premières étapes d'un pipeline de traitement de
l'information (comme un réseau de neurones) pour garantir que
l'information pertinente n'est pas perdue prématurément, ce qui
limiterait la capacité des étapes ultérieures à accomplir leur tâche.
\newpage
\section{Exercice 5}
\section{Exercice 5/10 : Jalon 5 - Applications, injections,
surjections, bijections et composition de
fonctions}\label{exercice-510-jalon-5---applications-injections-surjections-bijections-et-composition-de-fonctions}

\textbf{Professeur Émérite :} Chers étudiantes et étudiants, abordons
aujourd'hui un résultat fondamental concernant la surjectivité et la
composition de fonctions. Ce type de démonstration abstraite est
essentiel pour développer une intuition solide en algèbre et en analyse,
et pour comprendre la structure des transformations.

\begin{center}\rule{0.5\linewidth}{0.5pt}\end{center}

\subsubsection{\texorpdfstring{Énoncé de l'Exercice (Difficulté :
\(\star\)\(\star\)\(\star\)$\circ$$\circ$)}{Énoncé de l'Exercice (Difficulté : \textbackslash star\textbackslash star\textbackslash star$\circ$$\circ$)}}\label{uxe9noncuxe9-de-lexercice-difficultuxe9-starstarstar}

Soient \(E\), \(F\) et \(G\) trois ensembles non vides. Considérons deux
fonctions : * \(g: E \to F\) * \(f: F \to G\)

On suppose que la fonction composée \(f \circ g: E \to G\) est
surjective.

Démontrez rigoureusement que la fonction \(f: F \to G\) est surjective.

\begin{center}\rule{0.5\linewidth}{0.5pt}\end{center}

\subsubsection{Analyse de l'Énoncé}\label{analyse-de-luxe9noncuxe9}

Cet exercice nous demande de prouver une implication : si la composition
de deux fonctions est surjective, alors la seconde fonction (celle qui
est appliquée en dernier) doit nécessairement être surjective.

\begin{enumerate}
\def\labelenumi{\arabic{enumi}.}
\tightlist
\item
  \textbf{Comprendre la surjectivité :}

  \begin{itemize}
  \tightlist
  \item
    Une fonction \(h: A \to B\) est dite surjective si et seulement si
    pour tout élément \(b\) de l'ensemble d'arrivée \(B\), il existe au
    moins un élément \(a\) de l'ensemble de départ \(A\) tel que
    \(h(a) = b\). En d'autres termes, l'image de \(h\) est égale à son
    ensemble d'arrivée : \(\text{Im}(h) = B\).
  \end{itemize}
\item
  \textbf{Hypothèse :} \(f \circ g: E \to G\) est surjective.

  \begin{itemize}
  \tightlist
  \item
    Cela signifie que pour tout \(z \in G\), il existe au moins un
    \(x \in E\) tel que \((f \circ g)(x) = z\).
  \item
    Par définition de la composition, cela équivaut à \(f(g(x)) = z\).
  \end{itemize}
\item
  \textbf{Conclusion à démontrer :} \(f: F \to G\) est surjective.

  \begin{itemize}
  \tightlist
  \item
    Cela signifie que pour tout \(z \in G\), il existe au moins un
    \(y \in F\) tel que \(f(y) = z\).
  \end{itemize}
\item
  \textbf{Stratégie de démonstration :}

  \begin{itemize}
  \tightlist
  \item
    Pour prouver que \(f\) est surjective, nous devons partir d'un
    élément arbitraire \(z\) dans l'ensemble d'arrivée de \(f\) (qui est
    \(G\)).
  \item
    Notre objectif est de trouver un élément \(y\) dans l'ensemble de
    départ de \(f\) (qui est \(F\)) tel que \(f(y) = z\).
  \item
    L'hypothèse sur la surjectivité de \(f \circ g\) est notre seule
    source d'information. Elle nous donne un \(x \in E\) tel que
    \(f(g(x)) = z\).
  \item
    L'expression \(f(g(x)) = z\) contient déjà la forme
    \(f(\text{quelque chose}) = z\). Le ``quelque chose'' est \(g(x)\).
  \item
    Si nous posons \(y = g(x)\), alors \(y\) est un élément de \(F\)
    (car \(g: E \to F\)) et nous aurons \(f(y) = z\). C'est précisément
    ce que nous cherchons !
  \end{itemize}
\end{enumerate}

\begin{center}\rule{0.5\linewidth}{0.5pt}\end{center}

\subsubsection{Correction Exhaustive
Pas-à-Pas}\label{correction-exhaustive-pas-uxe0-pas}

Pour démontrer que \(f: F \to G\) est surjective, nous devons montrer
que pour tout élément \(z\) de l'ensemble d'arrivée \(G\), il existe au
moins un élément \(y\) de l'ensemble de départ \(F\) tel que
\(f(y) = z\).

\begin{enumerate}
\def\labelenumi{\arabic{enumi}.}
\item
  \textbf{Initialisation :} Soit \(z\) un élément quelconque et
  arbitrairement choisi de l'ensemble \(G\). Notre but est de trouver un
  \(y \in F\) tel que \(f(y) = z\).
\item
  \textbf{Utilisation de l'hypothèse :} Nous savons par hypothèse que la
  fonction composée \(f \circ g: E \to G\) est surjective.

  \begin{itemize}
  \tightlist
  \item
    Par définition de la surjectivité, cela signifie que pour tout
    élément de l'ensemble d'arrivée \(G\), il existe au moins un élément
    dans l'ensemble de départ \(E\) qui lui est appliqué.
  \item
    Puisque \(z \in G\) (choisi à l'étape 1), il existe donc au moins un
    élément \(x_0 \in E\) tel que \((f \circ g)(x_0) = z\).
  \end{itemize}
\item
  \textbf{Décomposition de la composition :} Par définition de la
  composition de fonctions, l'expression \((f \circ g)(x_0)\) est
  équivalente à \(f(g(x_0))\).

  \begin{itemize}
  \tightlist
  \item
    En substituant cette équivalence dans l'équation de l'étape 2, nous
    obtenons : \[f(g(x_0)) = z\]
  \end{itemize}
\item
  \textbf{Identification du candidat pour \(y\) :} Nous cherchons un
  élément \(y \in F\) tel que \(f(y) = z\). L'équation \(f(g(x_0)) = z\)
  nous fournit directement un tel candidat.

  \begin{itemize}
  \tightlist
  \item
    Posons \(y_0 = g(x_0)\).
  \end{itemize}
\item
  \textbf{Vérification de l'appartenance de \(y_0\) à \(F\) :} La
  fonction \(g\) est définie comme \(g: E \to F\).

  \begin{itemize}
  \tightlist
  \item
    Puisque \(x_0 \in E\), l'image de \(x_0\) par \(g\), qui est
    \(y_0 = g(x_0)\), est nécessairement un élément de l'ensemble \(F\).
  \item
    Donc, \(y_0 \in F\).
  \end{itemize}
\item
  \textbf{Vérification de la condition \(f(y_0) = z\) :} En substituant
  \(y_0\) dans l'équation de l'étape 3, nous obtenons :

  \begin{itemize}
  \tightlist
  \item
    \(f(y_0) = z\).
  \end{itemize}
\item
  \textbf{Conclusion :} Nous avons démontré que pour tout \(z \in G\)
  (choisi arbitrairement au début), il existe un élément \(y_0 \in F\)
  (spécifiquement \(y_0 = g(x_0)\) pour un certain \(x_0 \in E\)) tel
  que \(f(y_0) = z\).

  \begin{itemize}
  \tightlist
  \item
    Ceci est précisément la définition de la surjectivité de la fonction
    \(f: F \to G\).
  \end{itemize}
\end{enumerate}

Par conséquent, si la fonction composée \(f \circ g\) est surjective,
alors la fonction \(f\) est surjective.

\begin{center}\rule{0.5\linewidth}{0.5pt}\end{center}

\subsubsection{Liens avec l'Intelligence
Artificielle}\label{liens-avec-lintelligence-artificielle}

Le concept de surjectivité et de composition de fonctions est
omniprésent en Intelligence Artificielle, particulièrement dans le
domaine de l'apprentissage profond (Deep Learning).

\begin{enumerate}
\def\labelenumi{\arabic{enumi}.}
\tightlist
\item
  \textbf{Réseaux de Neurones comme Compositions de Fonctions :}

  \begin{itemize}
  \tightlist
  \item
    Un réseau de neurones est fondamentalement une composition de
    fonctions. Chaque couche du réseau peut être vue comme une fonction
    \(f_i: H_{i-1} \to H_i\), où \(H_{i-1}\) et \(H_i\) sont les espaces
    des activations de la couche précédente et de la couche actuelle,
    respectivement.
  \item
    Le réseau entier est alors une fonction composée
    \(F = f_L \circ f_{L-1} \circ \dots \circ f_1: X \to Y\), où \(X\)
    est l'espace d'entrée et \(Y\) est l'espace de sortie.
  \item
    Dans notre exercice, \(g\) pourrait représenter les premières
    couches du réseau (\(f_{L-1} \circ \dots \circ f_1\)), et \(f\) la
    dernière couche (ou les dernières couches) du réseau (\(f_L\)).
  \end{itemize}
\item
  \textbf{Surjectivité et Capacité Représentative :}

  \begin{itemize}
  \tightlist
  \item
    Si le réseau global \(F = f \circ g\) est capable de produire
    n'importe quelle sortie désirée dans l'espace \(G\) (c'est-à-dire
    qu'il est ``surjectif'' sur l'espace des cibles), alors notre
    démonstration implique que la dernière partie du réseau, \(f\), doit
    elle-même être surjective.
  \item
    Cela signifie que la couche de sortie (ou le ``head'' du modèle)
    doit avoir une capacité suffisante pour mapper les représentations
    internes (les ``features'' produites par \(g\)) à l'ensemble complet
    des sorties possibles. Si \(f\) n'était pas surjective, il
    existerait des sorties cibles que le réseau ne pourrait jamais
    atteindre, peu importe la qualité des représentations générées par
    \(g\).
  \end{itemize}
\item
  \textbf{Modèles Génératifs (GANs, VAEs) :}

  \begin{itemize}
  \tightlist
  \item
    Dans les modèles génératifs, comme les Generative Adversarial
    Networks (GANs) ou les Variational Autoencoders (VAEs), le
    générateur \(G\) est une fonction qui mappe un espace latent \(Z\)
    (souvent un bruit aléatoire) vers l'espace des données réelles
    \(X\). L'objectif est que \(G\) soit ``surjectif'' sur la variété
    des données réelles, c'est-à-dire qu'il puisse générer n'importe
    quelle donnée réaliste.
  \item
    Si le générateur \(G\) est lui-même un réseau profond,
    \(G = G_k \circ \dots \circ G_1\), alors la surjectivité de
    l'ensemble \(G\) implique que la dernière couche \(G_k\) doit être
    capable de couvrir l'espace des données. Si \(G_k\) est trop
    restrictive, le modèle ne pourra pas générer toute la diversité des
    données, même si les couches précédentes ont appris des
    représentations latentes riches.
  \end{itemize}
\item
  \textbf{Implications pour la Conception de Modèles :}

  \begin{itemize}
  \tightlist
  \item
    Ce théorème souligne l'importance de la conception de la couche de
    sortie. Si l'on sait que le problème requiert une couverture
    complète de l'espace de sortie (par exemple, classification
    multi-classes où toutes les classes sont possibles, ou génération
    d'images avec une grande diversité), alors la fonction de la
    dernière couche doit être choisie de manière à être potentiellement
    surjective.
  \item
    Par exemple, pour une classification, une couche
    \passthrough{\lstinline!softmax!} est souvent utilisée, qui peut en
    principe produire n'importe quelle distribution de probabilité sur
    les classes, permettant ainsi à la fonction \(f\) d'être
    ``surjective'' sur l'espace des distributions de probabilité.
  \end{itemize}
\end{enumerate}

En somme, ce résultat mathématique simple a des répercussions profondes
sur la compréhension de la capacité et des limitations des architectures
d'apprentissage profond, en soulignant que la ``puissance'' de la
transformation finale est cruciale pour la capacité globale du système à
atteindre toutes les cibles possibles.
\newpage
\section{Exercice 6}
\section{\texorpdfstring{Exercice 6/10 : Construction d'une bijection
entre \(\mathbb{N}\) et
\(\mathbb{Z}\)}{Exercice 6/10 : Construction d'une bijection entre \textbackslash mathbb\{N\} et \textbackslash mathbb\{Z\}}}\label{exercice-610-construction-dune-bijection-entre-mathbbn-et-mathbbz}

\textbf{Jalon 5 : Applications, injections, surjections, bijections et
composition de fonctions}

\textbf{Niveau de difficulté :} \(\star\)\(\star\)\(\star\)$\circ$$\circ$

\begin{center}\rule{0.5\linewidth}{0.5pt}\end{center}

\subsection{Énoncé}\label{uxe9noncuxe9}

Soient les ensembles \(\mathbb{N} = \{0, 1, 2, 3, \dots\}\) l'ensemble
des nombres entiers naturels et
\(\mathbb{Z} = \{\dots, -2, -1, 0, 1, 2, \dots\}\) l'ensemble des
nombres entiers relatifs.

\begin{enumerate}
\def\labelenumi{\arabic{enumi}.}
\tightlist
\item
  Construire explicitement une fonction
  \(f: \mathbb{N} \to \mathbb{Z}\).
\item
  Démontrer rigoureusement que la fonction \(f\) ainsi construite est
  une bijection.
\end{enumerate}

\begin{center}\rule{0.5\linewidth}{0.5pt}\end{center}

\subsection{Analyse de l'énoncé}\label{analyse-de-luxe9noncuxe9}

L'objectif de cet exercice est de démontrer l'équipotence des ensembles
\(\mathbb{N}\) et \(\mathbb{Z}\) en construisant une bijection explicite
entre eux. Bien que l'on sache intuitivement que ces deux ensembles sont
``infinis'', la notion d'équipotence pour les ensembles infinis est plus
subtile et nécessite la construction d'une fonction qui soit à la fois
injective et surjective.

\textbf{1. Construction explicite de \(f: \mathbb{N} \to \mathbb{Z}\) :}
Nous devons trouver une règle qui associe à chaque entier naturel \(n\)
un unique entier relatif \(f(n)\), de telle sorte que tous les entiers
relatifs soient atteints et qu'aucun ne soit atteint plus d'une fois.
Une stratégie courante pour construire une bijection entre
\(\mathbb{N}\) et \(\mathbb{Z}\) est de ``zigzaguer'' entre les entiers
positifs (y compris zéro) et les entiers négatifs. Par exemple : *
\(0 \mapsto 0\) * \(1 \mapsto -1\) * \(2 \mapsto 1\) * \(3 \mapsto -2\)
* \(4 \mapsto 2\) * \(5 \mapsto -3\) etc.

Cette observation suggère une distinction basée sur la parité de \(n\):
* Si \(n\) est pair, \(n = 2k\) pour un certain \(k \in \mathbb{N}\),
alors \(f(n)\) devrait être un entier non-négatif. En observant la
séquence, \(f(0)=0\), \(f(2)=1\), \(f(4)=2\), on voit que \(f(2k) = k\).
* Si \(n\) est impair, \(n = 2k+1\) pour un certain
\(k \in \mathbb{N}\), alors \(f(n)\) devrait être un entier négatif. En
observant la séquence, \(f(1)=-1\), \(f(3)=-2\), \(f(5)=-3\), on voit
que \(f(2k+1) = -(k+1)\).

Nous allons formaliser cette fonction par morceaux.

\textbf{2. Démonstration de la bijectivité :} Pour prouver que \(f\) est
une bijection, nous devons démontrer deux propriétés : *
\textbf{Injectivité (ou injection) :} Pour tous
\(n_1, n_2 \in \mathbb{N}\), si \(f(n_1) = f(n_2)\), alors
\(n_1 = n_2\). Cela signifie que deux éléments distincts de
\(\mathbb{N}\) ne peuvent pas avoir la même image dans \(\mathbb{Z}\). *
\textbf{Surjectivité (ou surjection) :} Pour tout \(z \in \mathbb{Z}\),
il existe au moins un \(n \in \mathbb{N}\) tel que \(f(n) = z\). Cela
signifie que chaque élément de \(\mathbb{Z}\) est l'image d'au moins un
élément de \(\mathbb{N}\).

La preuve de l'injectivité et de la surjectivité nécessitera une analyse
par cas, en fonction de la parité de \(n\) pour l'injectivité, et du
signe de \(z\) pour la surjectivité.

\begin{center}\rule{0.5\linewidth}{0.5pt}\end{center}

\subsection{Correction exhaustive
pas-à-pas}\label{correction-exhaustive-pas-uxe0-pas}

\subsubsection{\texorpdfstring{1. Construction explicite de la fonction
\(f: \mathbb{N} \to \mathbb{Z}\)}{1. Construction explicite de la fonction f: \textbackslash mathbb\{N\} \textbackslash to \textbackslash mathbb\{Z\}}}\label{construction-explicite-de-la-fonction-f-mathbbn-to-mathbbz}

Nous proposons la fonction \(f\) définie par :
\[ f(n) = \begin{cases} \frac{n}{2} & \text{si } n \text{ est pair} \\ -\frac{n+1}{2} & \text{si } n \text{ est impair} \end{cases} \]

Vérifions quelques valeurs pour s'assurer de la cohérence avec la
stratégie de ``zigzag'' : * Pour \(n=0\) (pair) :
\(f(0) = \frac{0}{2} = 0\). * Pour \(n=1\) (impair) :
\(f(1) = -\frac{1+1}{2} = -\frac{2}{2} = -1\). * Pour \(n=2\) (pair) :
\(f(2) = \frac{2}{2} = 1\). * Pour \(n=3\) (impair) :
\(f(3) = -\frac{3+1}{2} = -\frac{4}{2} = -2\). * Pour \(n=4\) (pair) :
\(f(4) = \frac{4}{2} = 2\).

La fonction semble correcte et correspond à notre intuition.

\subsubsection{\texorpdfstring{2. Démonstration que \(f\) est une
bijection}{2. Démonstration que f est une bijection}}\label{duxe9monstration-que-f-est-une-bijection}

Pour prouver que \(f\) est une bijection, nous devons montrer qu'elle
est injective et surjective.

\paragraph{a) Preuve de
l'injectivité}\label{a-preuve-de-linjectivituxe9}

Soient \(n_1, n_2 \in \mathbb{N}\) tels que \(f(n_1) = f(n_2)\). Nous
devons montrer que \(n_1 = n_2\). Soit \(z = f(n_1) = f(n_2)\).

Observons les images de \(f\) en fonction de la parité de \(n\): * Si
\(n\) est pair, \(n=2k\) pour un \(k \in \mathbb{N}\), alors
\(f(n) = k\). Les images sont \(\{0, 1, 2, \dots\} = \mathbb{N}\). * Si
\(n\) est impair, \(n=2k+1\) pour un \(k \in \mathbb{N}\), alors
\(f(n) = -(k+1)\). Les images sont
\(\{-1, -2, -3, \dots\} = \mathbb{Z}^-\).

Les ensembles d'images pour les \(n\) pairs (\(\mathbb{N}\)) et les
\(n\) impairs (\(\mathbb{Z}^-\)) sont disjoints. En effet,
\(\mathbb{N} \cap \mathbb{Z}^- = \emptyset\). Par conséquent, si
\(f(n_1) = f(n_2)\), alors \(f(n_1)\) et \(f(n_2)\) doivent
nécessairement appartenir au même sous-ensemble d'images. Cela implique
que \(n_1\) et \(n_2\) doivent avoir la même parité.

Nous considérons deux cas :

\textbf{Cas 1 : \(n_1\) et \(n_2\) sont tous les deux pairs.} Soient
\(n_1 = 2k_1\) et \(n_2 = 2k_2\) pour certains
\(k_1, k_2 \in \mathbb{N}\). Alors
\(f(n_1) = \frac{n_1}{2} = \frac{2k_1}{2} = k_1\). Et
\(f(n_2) = \frac{n_2}{2} = \frac{2k_2}{2} = k_2\). Puisque
\(f(n_1) = f(n_2)\), nous avons \(k_1 = k_2\). En multipliant par 2,
nous obtenons \(2k_1 = 2k_2\), ce qui signifie \(n_1 = n_2\).

\textbf{Cas 2 : \(n_1\) et \(n_2\) sont tous les deux impairs.} Soient
\(n_1 = 2k_1+1\) et \(n_2 = 2k_2+1\) pour certains
\(k_1, k_2 \in \mathbb{N}\). Alors
\(f(n_1) = -\frac{n_1+1}{2} = -\frac{(2k_1+1)+1}{2} = -\frac{2k_1+2}{2} = -(k_1+1)\).
Et
\(f(n_2) = -\frac{n_2+1}{2} = -\frac{(2k_2+1)+1}{2} = -\frac{2k_2+2}{2} = -(k_2+1)\).
Puisque \(f(n_1) = f(n_2)\), nous avons \(-(k_1+1) = -(k_2+1)\). En
multipliant par \(-1\), nous obtenons \(k_1+1 = k_2+1\). En soustrayant
1, nous obtenons \(k_1 = k_2\). En multipliant par 2 et en ajoutant 1,
nous obtenons \(2k_1+1 = 2k_2+1\), ce qui signifie \(n_1 = n_2\).

Dans tous les cas possibles où \(f(n_1) = f(n_2)\), nous avons montré
que \(n_1 = n_2\). Par conséquent, la fonction \(f\) est injective.

\paragraph{b) Preuve de la
surjectivité}\label{b-preuve-de-la-surjectivituxe9}

Soit \(z \in \mathbb{Z}\). Nous devons trouver un \(n \in \mathbb{N}\)
tel que \(f(n) = z\). Nous considérons deux cas pour \(z\):

\textbf{Cas 1 : \(z \ge 0\).} (C'est-à-dire \(z \in \mathbb{N}\)) Nous
cherchons un \(n \in \mathbb{N}\) tel que \(f(n) = z\). Puisque
\(z \ge 0\), l'image \(z\) doit provenir d'un \(n\) pair, selon la
définition de \(f\). Nous posons donc \(f(n) = \frac{n}{2}\). Pour que
\(\frac{n}{2} = z\), il faut que \(n = 2z\). Vérifions que ce \(n\) est
valide : * Puisque \(z \in \mathbb{N}\), \(z \ge 0\). Donc \(2z \ge 0\).
Ainsi \(n = 2z \in \mathbb{N}\). * Puisque \(n = 2z\), \(n\) est un
nombre pair. * En appliquant \(f\) à ce \(n\):
\(f(n) = f(2z) = \frac{2z}{2} = z\). Nous avons bien trouvé un
\(n \in \mathbb{N}\) (en l'occurrence \(n=2z\)) tel que \(f(n)=z\) pour
tout \(z \in \mathbb{N}\).

\textbf{Cas 2 : \(z < 0\).} (C'est-à-dire \(z \in \mathbb{Z}^-\)) Nous
cherchons un \(n \in \mathbb{N}\) tel que \(f(n) = z\). Puisque
\(z < 0\), l'image \(z\) doit provenir d'un \(n\) impair, selon la
définition de \(f\). Nous posons donc \(f(n) = -\frac{n+1}{2}\). Pour
que \(-\frac{n+1}{2} = z\), nous résolvons pour \(n\):
\(-\frac{n+1}{2} = z\) \(\frac{n+1}{2} = -z\) \(n+1 = -2z\)
\(n = -2z-1\). Vérifions que ce \(n\) est valide : * Puisque
\(z \in \mathbb{Z}^-\), \(z \le -1\). * Alors \(-z \ge 1\). * Donc
\(-2z \ge 2\). * Par conséquent, \(n = -2z-1 \ge 2-1 = 1\). * Ainsi
\(n = -2z-1 \in \mathbb{N}\) (c'est un entier naturel strictement
positif). * Puisque \(n = -2z-1\), \(n\) est de la forme
\(2(\text{entier})+1\) (car \(-2z\) est pair, donc \(-2z-1\) est
impair). Donc \(n\) est un nombre impair. * En appliquant \(f\) à ce
\(n\):
\(f(n) = f(-2z-1) = -\frac{(-2z-1)+1}{2} = -\frac{-2z}{2} = -(-z) = z\).
Nous avons bien trouvé un \(n \in \mathbb{N}\) (en l'occurrence
\(n=-2z-1\)) tel que \(f(n)=z\) pour tout \(z \in \mathbb{Z}^-\).

Dans tous les cas, pour tout \(z \in \mathbb{Z}\), nous avons trouvé un
\(n \in \mathbb{N}\) tel que \(f(n)=z\). Par conséquent, la fonction
\(f\) est surjective.

\paragraph{c) Conclusion}\label{c-conclusion}

Puisque la fonction \(f: \mathbb{N} \to \mathbb{Z}\) est à la fois
injective et surjective, elle est une bijection. Ceci démontre que les
ensembles \(\mathbb{N}\) et \(\mathbb{Z}\) ont la même cardinalité,
c'est-à-dire qu'ils sont équipotents.

\begin{center}\rule{0.5\linewidth}{0.5pt}\end{center}

\subsection{Liens avec l'Intelligence
Artificielle}\label{liens-avec-lintelligence-artificielle}

La construction d'une bijection entre \(\mathbb{N}\) et \(\mathbb{Z}\)
est un concept fondamental en mathématiques discrètes et en théorie de
la calculabilité, qui sont des piliers théoriques de l'Intelligence
Artificielle.

\begin{enumerate}
\def\labelenumi{\arabic{enumi}.}
\item
  \textbf{Numérisation et Encodage des Données :} En IA, de nombreuses
  données sont de nature discrète (catégories, mots, symboles). Pour
  qu'un algorithme puisse les traiter, elles doivent souvent être
  converties en représentations numériques. Une bijection comme celle-ci
  illustre le principe d'encodage : chaque élément d'un ensemble (ici
  \(\mathbb{Z}\)) peut être mappé de manière unique et réversible à un
  élément d'un autre ensemble (ici \(\mathbb{N}\)). En PNL (Traitement
  du Langage Naturel), par exemple, les mots d'un vocabulaire sont
  souvent encodés en entiers (indices), et des techniques plus avancées
  comme le ``word embedding'' construisent des vecteurs numériques pour
  chaque mot. La capacité de passer d'une représentation à l'autre sans
  perte d'information est cruciale.
\item
  \textbf{Théorie de la Calculabilité et Limites de l'IA :} La
  démonstration que \(\mathbb{N}\) et \(\mathbb{Z}\) sont équipotents
  est un exemple simple de la notion de ``dénombrabilité''. La théorie
  de la calculabilité, développée par des figures comme Alan Turing,
  repose sur l'idée que tout ce qui est ``calculable'' par un algorithme
  peut être représenté par des opérations sur des nombres naturels.
  L'ensemble de tous les programmes informatiques possibles est
  dénombrable (on peut les énumérer avec des entiers naturels).
  Comprendre les propriétés des ensembles dénombrables et non
  dénombrables (comme \(\mathbb{R}\)) est essentiel pour définir les
  limites théoriques de ce que les ordinateurs et, par extension, les
  systèmes d'IA peuvent accomplir. Par exemple, le problème de l'arrêt
  de Turing démontre qu'il n'existe pas d'algorithme général pour
  déterminer si un programme donné se terminera ou non, une limitation
  fondamentale qui s'applique à toute IA.
\item
  \textbf{Structures de Données et Hachage :} Bien que la bijection
  entre \(\mathbb{N}\) et \(\mathbb{Z}\) soit un cas spécifique, le
  principe de mappage unique est lié aux fonctions de hachage utilisées
  dans les structures de données (tables de hachage) qui sont
  omniprésentes en IA pour l'efficacité. Une fonction de hachage tente
  de mapper des données d'entrée de taille arbitraire à une valeur de
  hachage de taille fixe (souvent un entier). Bien que les fonctions de
  hachage ne soient pas des bijections (des collisions peuvent se
  produire), l'objectif est de s'en rapprocher le plus possible pour
  garantir une distribution uniforme et des recherches rapides, ce qui
  est vital pour la performance des algorithmes d'apprentissage
  automatique et des bases de données utilisées par l'IA.
\item
  \textbf{Numérotation de Gödel :} Un exemple plus avancé de bijection
  entre des objets complexes et les nombres naturels est la numérotation
  de Gödel. Cette technique attribue un nombre naturel unique à chaque
  formule et preuve dans un système formel. Cela permet de traduire des
  énoncés sur les propriétés des systèmes formels en énoncés sur les
  propriétés des nombres naturels, ce qui a conduit aux célèbres
  théorèmes d'incomplétude de Gödel. Ces théorèmes ont des implications
  profondes pour la logique, les fondements des mathématiques et,
  indirectement, pour la compréhension des limites de la raisonnement
  formel et de l'intelligence artificielle symbolique.
\end{enumerate}
\newpage
\section{Exercice 7}
\section{Exercice 7/10 : Images directes et images réciproques
d'ensembles}\label{exercice-710-images-directes-et-images-ruxe9ciproques-densembles}

\textbf{Jalon 5 : Applications, injections, surjections, bijections et
composition de fonctions}

\textbf{Niveau de difficulté :} \(\star\)\(\star\)\(\star\)\(\star\)$\circ$

\begin{center}\rule{0.5\linewidth}{0.5pt}\end{center}

\subsubsection{Énoncé}\label{uxe9noncuxe9}

Soient \(E\) et \(F\) deux ensembles non vides. Soit \(f: E \to F\) une
application.

\begin{enumerate}
\def\labelenumi{\arabic{enumi}.}
\item
  Soient \(A\) et \(B\) deux sous-ensembles de \(E\) (i.e.,
  \(A \subseteq E\) et \(B \subseteq E\)). Démontrer que l'image directe
  de l'union de \(A\) et \(B\) par \(f\) est égale à l'union des images
  directes de \(A\) et \(B\) par \(f\). Autrement dit, démontrer que
  \(f(A \cup B) = f(A) \cup f(B)\).
\item
  Soient \(A'\) et \(B'\) deux sous-ensembles de \(F\) (i.e.,
  \(A' \subseteq F\) et \(B' \subseteq F\)). Démontrer que l'image
  réciproque de l'intersection de \(A'\) et \(B'\) par \(f\) est égale à
  l'intersection des images réciproques de \(A'\) et \(B'\) par \(f\).
  Autrement dit, démontrer que
  \(f^{-1}(A' \cap B') = f^{-1}(A') \cap f^{-1}(B')\).
\end{enumerate}

\begin{center}\rule{0.5\linewidth}{0.5pt}\end{center}

\subsubsection{Analyse de l'énoncé}\label{analyse-de-luxe9noncuxe9}

Cet exercice vise à consolider la compréhension des définitions d'image
directe et d'image réciproque d'ensembles par une application, ainsi que
leur interaction avec les opérations ensemblistes fondamentales que sont
l'union et l'intersection.

Pour démontrer l'égalité de deux ensembles \(X\) et \(Y\) (par exemple,
\(X = f(A \cup B)\) et \(Y = f(A) \cup f(B)\)), la méthode standard
consiste à prouver la double inclusion : 1. \(X \subseteq Y\) (tout
élément de \(X\) est un élément de \(Y\)). 2. \(Y \subseteq X\) (tout
élément de \(Y\) est un élément de \(X\)).

Nous devrons nous appuyer rigoureusement sur les définitions : *
\textbf{Image directe d'un ensemble \(S \subseteq E\) :}
\(f(S) = \{y \in F \mid \exists x \in S, y = f(x)\}\). * \textbf{Image
réciproque d'un ensemble \(S' \subseteq F\) :}
\(f^{-1}(S') = \{x \in E \mid f(x) \in S'\}\). * \textbf{Union
d'ensembles :} \(x \in X \cup Y \iff x \in X \text{ ou } x \in Y\). *
\textbf{Intersection d'ensembles :}
\(x \in X \cap Y \iff x \in X \text{ et } x \in Y\).

La difficulté réside dans la manipulation précise des quantificateurs
(\(\exists\), \(\forall\)) et des connecteurs logiques (et, ou) à chaque
étape de la démonstration, sans aucune ellipse mathématique.

\begin{center}\rule{0.5\linewidth}{0.5pt}\end{center}

\subsubsection{Correction exhaustive
pas-à-pas}\label{correction-exhaustive-pas-uxe0-pas}

\paragraph{\texorpdfstring{Question 1 : Démontrer que
\(f(A \cup B) = f(A) \cup f(B)\)}{Question 1 : Démontrer que f(A \textbackslash cup B) = f(A) \textbackslash cup f(B)}}\label{question-1-duxe9montrer-que-fa-cup-b-fa-cup-fb}

Pour démontrer cette égalité, nous allons prouver la double inclusion :
\(f(A \cup B) \subseteq f(A) \cup f(B)\) et
\(f(A) \cup f(B) \subseteq f(A \cup B)\).

\subparagraph{\texorpdfstring{Étape 1.1 : Démonstration de
\(f(A \cup B) \subseteq f(A) \cup f(B)\)}{Étape 1.1 : Démonstration de f(A \textbackslash cup B) \textbackslash subseteq f(A) \textbackslash cup f(B)}}\label{uxe9tape-1.1-duxe9monstration-de-fa-cup-b-subseteq-fa-cup-fb}

Soit \(y\) un élément quelconque de l'ensemble \(f(A \cup B)\). Par
définition de l'image directe, cela signifie qu'il existe au moins un
élément \(x\) dans l'ensemble \(A \cup B\) tel que \(y = f(x)\).
\[y \in f(A \cup B) \implies \exists x \in A \cup B \text{ tel que } y = f(x)\]
Puisque \(x \in A \cup B\), par définition de l'union, cela signifie que
\(x\) appartient à \(A\) ou \(x\) appartient à \(B\).
\[x \in A \cup B \iff (x \in A \text{ ou } x \in B)\] Nous avons donc
deux cas possibles pour cet élément \(x\) :

\textbf{Cas 1 :} \(x \in A\). Si \(x \in A\) et \(y = f(x)\), alors par
définition de l'image directe, \(y\) est un élément de \(f(A)\).
\[x \in A \text{ et } y = f(x) \implies y \in f(A)\]

\textbf{Cas 2 :} \(x \in B\). Si \(x \in B\) et \(y = f(x)\), alors par
définition de l'image directe, \(y\) est un élément de \(f(B)\).
\[x \in B \text{ et } y = f(x) \implies y \in f(B)\]

Dans les deux cas (soit \(y \in f(A)\), soit \(y \in f(B)\)), nous
pouvons conclure que \(y\) appartient à l'union de \(f(A)\) et \(f(B)\).
\[(y \in f(A) \text{ ou } y \in f(B)) \implies y \in f(A) \cup f(B)\]
Puisque nous avons montré que si \(y \in f(A \cup B)\), alors
\(y \in f(A) \cup f(B)\), nous avons prouvé l'inclusion :
\[f(A \cup B) \subseteq f(A) \cup f(B)\]

\subparagraph{\texorpdfstring{Étape 1.2 : Démonstration de
\(f(A) \cup f(B) \subseteq f(A \cup B)\)}{Étape 1.2 : Démonstration de f(A) \textbackslash cup f(B) \textbackslash subseteq f(A \textbackslash cup B)}}\label{uxe9tape-1.2-duxe9monstration-de-fa-cup-fb-subseteq-fa-cup-b}

Soit \(y\) un élément quelconque de l'ensemble \(f(A) \cup f(B)\). Par
définition de l'union, cela signifie que \(y\) appartient à \(f(A)\) ou
\(y\) appartient à \(f(B)\).
\[y \in f(A) \cup f(B) \iff (y \in f(A) \text{ ou } y \in f(B))\] Nous
avons donc deux cas possibles pour cet élément \(y\) :

\textbf{Cas 1 :} \(y \in f(A)\). Si \(y \in f(A)\), alors par définition
de l'image directe, il existe au moins un élément \(x_A\) dans \(A\) tel
que \(y = f(x_A)\).
\[y \in f(A) \implies \exists x_A \in A \text{ tel que } y = f(x_A)\]
Puisque \(x_A \in A\), par définition de l'union, il est également vrai
que \(x_A \in A \cup B\). \[x_A \in A \implies x_A \in A \cup B\] Ainsi,
nous avons un élément \(x_A \in A \cup B\) tel que \(y = f(x_A)\). Par
définition de l'image directe, cela signifie que \(y \in f(A \cup B)\).
\[(\exists x_A \in A \cup B \text{ tel que } y = f(x_A)) \implies y \in f(A \cup B)\]

\textbf{Cas 2 :} \(y \in f(B)\). Si \(y \in f(B)\), alors par définition
de l'image directe, il existe au moins un élément \(x_B\) dans \(B\) tel
que \(y = f(x_B)\).
\[y \in f(B) \implies \exists x_B \in B \text{ tel que } y = f(x_B)\]
Puisque \(x_B \in B\), par définition de l'union, il est également vrai
que \(x_B \in A \cup B\). \[x_B \in B \implies x_B \in A \cup B\] Ainsi,
nous avons un élément \(x_B \in A \cup B\) tel que \(y = f(x_B)\). Par
définition de l'image directe, cela signifie que \(y \in f(A \cup B)\).
\[(\exists x_B \in A \cup B \text{ tel que } y = f(x_B)) \implies y \in f(A \cup B)\]

Dans les deux cas (soit \(y \in f(A)\), soit \(y \in f(B)\)), nous avons
montré que \(y \in f(A \cup B)\). Puisque nous avons montré que si
\(y \in f(A) \cup f(B)\), alors \(y \in f(A \cup B)\), nous avons prouvé
l'inclusion : \[f(A) \cup f(B) \subseteq f(A \cup B)\]

\subparagraph{Étape 1.3 : Conclusion pour la Question
1}\label{uxe9tape-1.3-conclusion-pour-la-question-1}

Puisque nous avons démontré les deux inclusions
\(f(A \cup B) \subseteq f(A) \cup f(B)\) et
\(f(A) \cup f(B) \subseteq f(A \cup B)\), nous pouvons conclure que les
deux ensembles sont égaux : \[f(A \cup B) = f(A) \cup f(B)\]

\paragraph{\texorpdfstring{Question 2 : Démontrer que
\(f^{-1}(A' \cap B') = f^{-1}(A') \cap f^{-1}(B')\)}{Question 2 : Démontrer que f\^{}\{-1\}(A\textquotesingle{} \textbackslash cap B\textquotesingle) = f\^{}\{-1\}(A\textquotesingle) \textbackslash cap f\^{}\{-1\}(B\textquotesingle)}}\label{question-2-duxe9montrer-que-f-1a-cap-b-f-1a-cap-f-1b}

Pour démontrer cette égalité, nous allons prouver la double inclusion :
\(f^{-1}(A' \cap B') \subseteq f^{-1}(A') \cap f^{-1}(B')\) et
\(f^{-1}(A') \cap f^{-1}(B') \subseteq f^{-1}(A' \cap B')\).

\subparagraph{\texorpdfstring{Étape 2.1 : Démonstration de
\(f^{-1}(A' \cap B') \subseteq f^{-1}(A') \cap f^{-1}(B')\)}{Étape 2.1 : Démonstration de f\^{}\{-1\}(A\textquotesingle{} \textbackslash cap B\textquotesingle) \textbackslash subseteq f\^{}\{-1\}(A\textquotesingle) \textbackslash cap f\^{}\{-1\}(B\textquotesingle)}}\label{uxe9tape-2.1-duxe9monstration-de-f-1a-cap-b-subseteq-f-1a-cap-f-1b}

Soit \(x\) un élément quelconque de l'ensemble \(f^{-1}(A' \cap B')\).
Par définition de l'image réciproque, cela signifie que l'image de \(x\)
par \(f\), c'est-à-dire \(f(x)\), appartient à l'ensemble
\(A' \cap B'\).
\[x \in f^{-1}(A' \cap B') \implies f(x) \in A' \cap B'\] Puisque
\(f(x) \in A' \cap B'\), par définition de l'intersection, cela signifie
que \(f(x)\) appartient à \(A'\) et \(f(x)\) appartient à \(B'\).
\[f(x) \in A' \cap B' \iff (f(x) \in A' \text{ et } f(x) \in B')\] Nous
pouvons séparer cette conjonction en deux affirmations :

\begin{enumerate}
\def\labelenumi{\arabic{enumi}.}
\item
  \(f(x) \in A'\). Par définition de l'image réciproque, si
  \(f(x) \in A'\), alors \(x\) est un élément de \(f^{-1}(A')\).
  \[f(x) \in A' \implies x \in f^{-1}(A')\]
\item
  \(f(x) \in B'\). Par définition de l'image réciproque, si
  \(f(x) \in B'\), alors \(x\) est un élément de \(f^{-1}(B')\).
  \[f(x) \in B' \implies x \in f^{-1}(B')\]
\end{enumerate}

Puisque \(x \in f^{-1}(A')\) et \(x \in f^{-1}(B')\), par définition de
l'intersection, nous pouvons conclure que \(x\) appartient à
l'intersection de \(f^{-1}(A')\) et \(f^{-1}(B')\).
\[(x \in f^{-1}(A') \text{ et } x \in f^{-1}(B')) \implies x \in f^{-1}(A') \cap f^{-1}(B')\]
Puisque nous avons montré que si \(x \in f^{-1}(A' \cap B')\), alors
\(x \in f^{-1}(A') \cap f^{-1}(B')\), nous avons prouvé l'inclusion :
\[f^{-1}(A' \cap B') \subseteq f^{-1}(A') \cap f^{-1}(B')\]

\subparagraph{\texorpdfstring{Étape 2.2 : Démonstration de
\(f^{-1}(A') \cap f^{-1}(B') \subseteq f^{-1}(A' \cap B')\)}{Étape 2.2 : Démonstration de f\^{}\{-1\}(A\textquotesingle) \textbackslash cap f\^{}\{-1\}(B\textquotesingle) \textbackslash subseteq f\^{}\{-1\}(A\textquotesingle{} \textbackslash cap B\textquotesingle)}}\label{uxe9tape-2.2-duxe9monstration-de-f-1a-cap-f-1b-subseteq-f-1a-cap-b}

Soit \(x\) un élément quelconque de l'ensemble
\(f^{-1}(A') \cap f^{-1}(B')\). Par définition de l'intersection, cela
signifie que \(x\) appartient à \(f^{-1}(A')\) et \(x\) appartient à
\(f^{-1}(B')\).
\[x \in f^{-1}(A') \cap f^{-1}(B') \iff (x \in f^{-1}(A') \text{ et } x \in f^{-1}(B'))\]
Nous pouvons séparer cette conjonction en deux affirmations :

\begin{enumerate}
\def\labelenumi{\arabic{enumi}.}
\item
  \(x \in f^{-1}(A')\). Par définition de l'image réciproque, si
  \(x \in f^{-1}(A')\), alors l'image de \(x\) par \(f\), c'est-à-dire
  \(f(x)\), appartient à \(A'\).
  \[x \in f^{-1}(A') \implies f(x) \in A'\]
\item
  \(x \in f^{-1}(B')\). Par définition de l'image réciproque, si
  \(x \in f^{-1}(B')\), alors l'image de \(x\) par \(f\), c'est-à-dire
  \(f(x)\), appartient à \(B'\).
  \[x \in f^{-1}(B') \implies f(x) \in B'\]
\end{enumerate}

Puisque \(f(x) \in A'\) et \(f(x) \in B'\), par définition de
l'intersection, nous pouvons conclure que \(f(x)\) appartient à
l'intersection de \(A'\) et \(B'\).
\[(f(x) \in A' \text{ et } f(x) \in B') \implies f(x) \in A' \cap B'\]
Enfin, si \(f(x) \in A' \cap B'\), par définition de l'image réciproque,
cela signifie que \(x\) est un élément de \(f^{-1}(A' \cap B')\).
\[f(x) \in A' \cap B' \implies x \in f^{-1}(A' \cap B')\] Puisque nous
avons montré que si \(x \in f^{-1}(A') \cap f^{-1}(B')\), alors
\(x \in f^{-1}(A' \cap B')\), nous avons prouvé l'inclusion :
\[f^{-1}(A') \cap f^{-1}(B') \subseteq f^{-1}(A' \cap B')\]

\subparagraph{Étape 2.3 : Conclusion pour la Question
2}\label{uxe9tape-2.3-conclusion-pour-la-question-2}

Puisque nous avons démontré les deux inclusions
\(f^{-1}(A' \cap B') \subseteq f^{-1}(A') \cap f^{-1}(B')\) et
\(f^{-1}(A') \cap f^{-1}(B') \subseteq f^{-1}(A' \cap B')\), nous
pouvons conclure que les deux ensembles sont égaux :
\[f^{-1}(A' \cap B') = f^{-1}(A') \cap f^{-1}(B')\]

\begin{center}\rule{0.5\linewidth}{0.5pt}\end{center}

\subsubsection{Liens avec l'Intelligence
Artificielle}\label{liens-avec-lintelligence-artificielle}

Les concepts d'images directes et réciproques d'ensembles, ainsi que
leurs propriétés en relation avec les opérations ensemblistes, sont des
fondations mathématiques omniprésentes en Intelligence Artificielle,
bien que souvent implicites dans les applications de haut niveau.

\begin{enumerate}
\def\labelenumi{\arabic{enumi}.}
\item
  \textbf{Traitement et Transformation des Données (Feature Engineering)
  :} En IA, les données sont fréquemment transformées d'un espace à un
  autre (par exemple, réduction de dimensionnalité, encodage,
  normalisation). Une telle transformation peut être modélisée comme une
  application \(f: E \to F\), où \(E\) est l'espace des données brutes
  et \(F\) l'espace des caractéristiques transformées.

  \begin{itemize}
  \tightlist
  \item
    La propriété \(f(A \cup B) = f(A) \cup f(B)\) signifie que si nous
    avons deux groupes de données \(A\) et \(B\) (par exemple, des
    images de chats et des images de chiens), l'ensemble des
    caractéristiques transformées pour l'ensemble combiné (chats ou
    chiens) est simplement l'union des caractéristiques transformées
    pour les chats et des caractéristiques transformées pour les chiens.
    Cette propriété intuitive est cruciale pour la cohérence des
    pipelines de traitement de données.
  \end{itemize}
\item
  \textbf{Classification et Reconnaissance de Formes :} Dans les tâches
  de classification, une fonction \(f\) (le modèle de classification)
  mappe un ensemble d'entrées \(E\) (par exemple, des images) à un
  ensemble de sorties \(F\) (par exemple, des étiquettes de classe).

  \begin{itemize}
  \tightlist
  \item
    La propriété \(f^{-1}(A' \cap B') = f^{-1}(A') \cap f^{-1}(B')\) est
    particulièrement pertinente. Supposons que \(A'\) représente
    l'ensemble des sorties ``chat'' et \(B'\) l'ensemble des sorties
    ``animal à quatre pattes''. L'intersection \(A' \cap B'\) serait
    l'ensemble des sorties ``chat ET animal à quatre pattes''. La
    propriété démontrée signifie que l'ensemble des images qui sont
    classées comme ``chat ET animal à quatre pattes'' est exactement
    l'intersection des images classées comme ``chat'' et des images
    classées comme ``animal à quatre pattes''. Cela garantit une
    interprétation logique et cohérente des prédictions du modèle,
    essentielle pour la robustesse et l'explicabilité des systèmes d'IA.
  \end{itemize}
\item
  \textbf{Vérification Formelle et Robustesse des Modèles :} Pour les
  systèmes d'IA critiques (véhicules autonomes, médecine), il est vital
  de prouver formellement certaines propriétés de leur comportement. Les
  identités ensemblistes comme celles-ci sont des briques de base pour
  des preuves plus complexes. Par exemple, pour analyser la robustesse
  d'un réseau de neurones face à des perturbations adverses, on pourrait
  étudier comment des ensembles d'entrées ``sûres'' (ou ``dangereuses'')
  se transforment à travers les couches du réseau, ou comment des
  ensembles de sorties désirées (ou indésirables) se
  ``rétro-projettent'' dans l'espace d'entrée.
\item
  \textbf{Représentation des Connaissances et Raisonnement Symbolique :}
  Dans l'IA symbolique, les connaissances sont souvent représentées par
  des ensembles et des relations. Les opérations ensemblistes et leurs
  propriétés sous des transformations fonctionnelles sont fondamentales
  pour l'inférence logique et la manipulation de concepts. Par exemple,
  si ``mammifère'' est un ensemble \(A'\) et ``carnivore'' est un
  ensemble \(B'\), alors ``mammifère carnivore'' est \(A' \cap B'\).
  Comprendre comment les fonctions (relations) interagissent avec ces
  ensembles est au cœur du raisonnement automatique.
\end{enumerate}

En somme, ces propriétés fondamentales de la théorie des ensembles et
des fonctions, bien que simples en apparence, sont les piliers logiques
qui sous-tendent la capacité des systèmes d'IA à traiter, transformer et
interpréter des informations de manière cohérente et prévisible.
\newpage
\section{Exercice 8}
\subsection{Exercice 8/10 : Injectivité et existence d'une fonction
inverse à
gauche}\label{exercice-810-injectivituxe9-et-existence-dune-fonction-inverse-uxe0-gauche}

\textbf{Niveau de difficulté :} \(\star\)\(\star\)\(\star\)\(\star\)$\circ$

\begin{center}\rule{0.5\linewidth}{0.5pt}\end{center}

\subsubsection{Énoncé}\label{uxe9noncuxe9}

Soient \(E\) et \(F\) deux ensembles non vides. Soit \(f: E \to F\) une
application.

Démontrer que \(f\) est injective si et seulement si il existe une
application \(g: F \to E\) telle que \(g \circ f = \text{Id}_E\).

Où \(\text{Id}_E: E \to E\) est l'application identité sur \(E\),
définie par \(\text{Id}_E(x) = x\) pour tout \(x \in E\).

\begin{center}\rule{0.5\linewidth}{0.5pt}\end{center}

\subsubsection{Analyse de l'énoncé}\label{analyse-de-luxe9noncuxe9}

Cet exercice nous demande de prouver une équivalence fondamentale en
théorie des ensembles et des fonctions, reliant la propriété
d'injectivité d'une fonction à l'existence d'une ``inverse à gauche''.

\begin{enumerate}
\def\labelenumi{\arabic{enumi}.}
\item
  \textbf{Comprendre l'injectivité :} Une application \(f: E \to F\) est
  dite injective si et seulement si pour tous \(x_1, x_2 \in E\),
  l'égalité \(f(x_1) = f(x_2)\) implique \(x_1 = x_2\). En d'autres
  termes, des éléments distincts de l'ensemble de départ \(E\) sont
  toujours envoyés sur des éléments distincts de l'ensemble d'arrivée
  \(F\). Il n'y a pas de ``collisions'' d'images.
\item
  \textbf{Comprendre l'inverse à gauche :} L'existence d'une application
  \(g: F \to E\) telle que \(g \circ f = \text{Id}_E\) signifie que la
  composition de \(g\) avec \(f\) redonne l'identité sur \(E\). Cela
  implique que pour tout \(x \in E\), \((g \circ f)(x) = g(f(x)) = x\).
  L'application \(g\) ``annule'' l'effet de \(f\) pour tout élément de
  \(E\).
\item
  \textbf{Structure de la preuve :} L'énoncé est un ``si et seulement
  si'' (équivalence logique), ce qui signifie que nous devons prouver
  deux implications distinctes :

  \begin{itemize}
  \tightlist
  \item
    \textbf{Implication 1 (\(\implies\)) :} Si \(f\) est injective,
    alors il existe une application \(g: F \to E\) telle que
    \(g \circ f = \text{Id}_E\). Pour cette partie, nous devrons
    \emph{construire} explicitement une telle fonction \(g\).
  \item
    \textbf{Implication 2 (\(\implies\)) :} S'il existe une
    application \(g: F \to E\) telle que \(g \circ f = \text{Id}_E\),
    alors \(f\) est injective. Pour cette partie, nous devrons utiliser
    la propriété de \(g\) pour démontrer l'injectivité de \(f\).
  \end{itemize}
\item
  \textbf{Considérations sur les ensembles :} L'énoncé précise que \(E\)
  et \(F\) sont des ensembles non vides. Cette condition est importante,
  notamment pour la construction de \(g\) dans la première implication,
  où nous pourrions avoir besoin de choisir un élément arbitraire de
  \(E\).
\end{enumerate}

\begin{center}\rule{0.5\linewidth}{0.5pt}\end{center}

\subsubsection{Correction exhaustive
pas-à-pas}\label{correction-exhaustive-pas-uxe0-pas}

Nous allons prouver les deux implications séparément.

\paragraph{\texorpdfstring{Partie 1 : \(f\) est injective \(\implies\)
il existe \(g: F \to E\) telle que
\(g \circ f = \text{Id}_E\).}{Partie 1 : f est injective \textbackslash implies il existe g: F \textbackslash to E telle que g \textbackslash circ f = \textbackslash text\{Id\}\_E.}}\label{partie-1-f-est-injective-implies-il-existe-g-f-to-e-telle-que-g-circ-f-textid_e.}

\textbf{Hypothèse :} \(f: E \to F\) est une application injective.
\textbf{Objectif :} Construire une application \(g: F \to E\) telle que
\(g \circ f = \text{Id}_E\).

Puisque \(E\) est un ensemble non vide, nous pouvons choisir et fixer un
élément arbitraire \(x_0 \in E\). Cet élément nous servira pour définir
\(g(y)\) lorsque \(y\) n'est pas une image par \(f\).

Nous allons définir l'application \(g: F \to E\) de la manière suivante
:

Pour tout \(y \in F\): * \textbf{Cas 1 :} Si \(y \in \text{Im}(f)\)
(l'image de \(f\)), cela signifie qu'il existe au moins un \(x \in E\)
tel que \(f(x) = y\). Puisque \(f\) est injective, si \(f(x_1) = y\) et
\(f(x_2) = y\), alors \(f(x_1) = f(x_2)\), ce qui implique \(x_1 = x_2\)
par définition de l'injectivité. Par conséquent, pour chaque
\(y \in \text{Im}(f)\), il existe un \emph{unique} \(x \in E\) tel que
\(f(x) = y\). Dans ce cas, nous définissons \(g(y) = x\), où \(x\) est
l'unique antécédent de \(y\) par \(f\).

\begin{itemize}
\tightlist
\item
  \textbf{Cas 2 :} Si \(y \notin \text{Im}(f)\), cela signifie qu'il
  n'existe aucun \(x \in E\) tel que \(f(x) = y\). Dans ce cas, nous
  définissons \(g(y) = x_0\), où \(x_0\) est l'élément arbitraire de
  \(E\) que nous avons fixé au début.
\end{itemize}

Cette définition de \(g\) est bien une application de \(F\) vers \(E\),
car pour chaque élément \(y \in F\), \(g(y)\) est défini de manière
unique et appartient à \(E\).

Maintenant, nous devons vérifier que \(g \circ f = \text{Id}_E\). Pour
cela, nous devons montrer que pour tout \(x \in E\),
\((g \circ f)(x) = x\).

Soit \(x \in E\) un élément quelconque. Calculons \((g \circ f)(x)\): 1.
D'abord, nous évaluons \(f(x)\). Soit \(y_x = f(x)\). 2. Par définition
de l'image, \(y_x = f(x)\) est un élément de \(\text{Im}(f)\). 3.
Puisque \(y_x \in \text{Im}(f)\), la définition de \(g\) pour ce \(y_x\)
relève du Cas 1. 4. Selon le Cas 1, \(g(y_x)\) est défini comme l'unique
antécédent de \(y_x\) par \(f\). 5. Nous savons que \(x\) est un
antécédent de \(y_x\) par \(f\), car \(f(x) = y_x\). 6. Puisque \(x\)
est l'unique antécédent de \(y_x\) par \(f\), il s'ensuit que
\(g(y_x) = x\). 7. En substituant \(y_x = f(x)\), nous obtenons
\(g(f(x)) = x\).

Cette égalité \(g(f(x)) = x\) est vraie pour tout \(x \in E\). Par
conséquent, \(g \circ f = \text{Id}_E\).

Nous avons donc construit une application \(g: F \to E\) telle que
\(g \circ f = \text{Id}_E\). L'implication
\(f \text{ est injective} \implies \text{il existe } g: F \to E \text{ telle que } g \circ f = \text{Id}_E\)
est démontrée.

\paragraph{\texorpdfstring{Partie 2 : Il existe \(g: F \to E\) telle que
\(g \circ f = \text{Id}_E \implies f\) est
injective.}{Partie 2 : Il existe g: F \textbackslash to E telle que g \textbackslash circ f = \textbackslash text\{Id\}\_E \textbackslash implies f est injective.}}\label{partie-2-il-existe-g-f-to-e-telle-que-g-circ-f-textid_e-implies-f-est-injective.}

\textbf{Hypothèse :} Il existe une application \(g: F \to E\) telle que
\(g \circ f = \text{Id}_E\). \textbf{Objectif :} Démontrer que \(f\) est
injective.

Pour prouver que \(f\) est injective, nous devons montrer que pour tous
\(x_1, x_2 \in E\), si \(f(x_1) = f(x_2)\), alors \(x_1 = x_2\).

Soient \(x_1, x_2 \in E\) deux éléments quelconques de \(E\). Supposons
que \(f(x_1) = f(x_2)\).

Puisque \(f(x_1)\) et \(f(x_2)\) sont des éléments de \(F\), nous
pouvons appliquer l'application \(g\) à ces deux éléments. En appliquant
\(g\) aux deux membres de l'égalité \(f(x_1) = f(x_2)\), nous obtenons :
\(g(f(x_1)) = g(f(x_2))\).

Par définition de la composition d'applications, nous avons :
\((g \circ f)(x_1) = (g \circ f)(x_2)\).

Par notre hypothèse, nous savons que \(g \circ f = \text{Id}_E\). Donc,
nous pouvons remplacer \((g \circ f)(x_1)\) par \(\text{Id}_E(x_1)\) et
\((g \circ f)(x_2)\) par \(\text{Id}_E(x_2)\). L'égalité devient :
\(\text{Id}_E(x_1) = \text{Id}_E(x_2)\).

Par définition de l'application identité \(\text{Id}_E\), nous avons
\(\text{Id}_E(x_1) = x_1\) et \(\text{Id}_E(x_2) = x_2\). Par
conséquent, l'égalité se simplifie en : \(x_1 = x_2\).

Nous avons donc montré que si \(f(x_1) = f(x_2)\), alors \(x_1 = x_2\).
Ceci est précisément la définition de l'injectivité de \(f\).

L'implication
\(\text{il existe } g: F \to E \text{ telle que } g \circ f = \text{Id}_E \implies f \text{ est injective}\)
est démontrée.

\paragraph{Conclusion}\label{conclusion}

Puisque nous avons prouvé les deux implications, nous pouvons conclure
que \(f\) est injective si et seulement si il existe une application
\(g: F \to E\) telle que \(g \circ f = \text{Id}_E\).

\begin{center}\rule{0.5\linewidth}{0.5pt}\end{center}

\subsubsection{Liens avec l'Intelligence
Artificielle}\label{liens-avec-lintelligence-artificielle}

Le concept d'injectivité et d'inverse à gauche trouve des résonances
profondes dans plusieurs domaines de l'Intelligence Artificielle,
notamment en ce qui concerne la représentation des données, la
compression et la robustesse des modèles.

\begin{enumerate}
\def\labelenumi{\arabic{enumi}.}
\tightlist
\item
  \textbf{Représentation et Encodage des Données (Autoencodeurs) :}

  \begin{itemize}
  \tightlist
  \item
    Dans les réseaux de neurones, en particulier les autoencodeurs,
    l'objectif est d'apprendre une représentation compacte (un encodage)
    des données d'entrée. Un autoencodeur est composé d'un encodeur
    \(f: X \to Z\) (où \(X\) est l'espace d'entrée et \(Z\) est l'espace
    latent) et d'un décodeur \(g: Z \to X\).
  \item
    Idéalement, nous souhaiterions que l'encodeur \(f\) soit une
    fonction injective. Si \(f\) est injective, cela signifie que chaque
    donnée d'entrée unique \(x \in X\) est mappée à une représentation
    latente unique \(z \in Z\). Aucune information n'est perdue lors de
    l'encodage.
  \item
    Dans ce scénario idéal, le décodeur \(g\) agirait comme une inverse
    à gauche (et même une inverse tout court si \(f\) est surjective sur
    \(Z\)), permettant de reconstruire parfaitement l'entrée originale
    \(x\) à partir de sa représentation latente \(z = f(x)\),
    c'est-à-dire \(g(f(x)) = x\).
  \item
    En pratique, les autoencodeurs ne garantissent pas l'injectivité
    (surtout si la dimension de \(Z\) est inférieure à celle de \(X\)),
    mais la recherche d'une reconstruction fidèle
    (\(g \circ f \approx \text{Id}_X\)) est une tentative d'approcher
    cette propriété.
  \end{itemize}
\item
  \textbf{Compression de Données sans Perte (Lossless Compression) :}

  \begin{itemize}
  \tightlist
  \item
    Les algorithmes de compression sans perte visent à transformer des
    données (par exemple, un fichier texte ou une image) en une
    représentation plus compacte sans aucune perte d'information.
  \item
    Le processus de compression peut être vu comme une fonction
    \(f: \text{Données Originales} \to \text{Données Compressées}\).
    Pour que la compression soit sans perte, \(f\) doit être injective.
    Si \(f\) n'était pas injective, deux ensembles de données originaux
    distincts pourraient être compressés en la même représentation,
    rendant impossible de déterminer l'original lors de la
    décompression.
  \item
    La décompression est alors l'application
    \(g: \text{Données Compressées} \to \text{Données Originales}\), qui
    doit agir comme une inverse à gauche (et même une inverse tout
    court) pour restaurer les données originales :
    \(g \circ f = \text{Id}_{\text{Données Originales}}\).
  \end{itemize}
\item
  \textbf{Fonctions de Hachage et Indexation :}

  \begin{itemize}
  \tightlist
  \item
    Les fonctions de hachage sont utilisées pour mapper des données de
    taille arbitraire à des valeurs de taille fixe (les ``hachages'').
    Elles sont fondamentales pour l'indexation de données (tables de
    hachage) et la vérification d'intégrité.
  \item
    Idéalement, une fonction de hachage ``parfaite'' serait injective
    sur l'ensemble des clés utilisées, c'est-à-dire qu'elle ne
    produirait jamais de collisions (deux clés différentes donnant le
    même hachage).
  \item
    Si une fonction de hachage \(h: K \to V\) (où \(K\) est l'ensemble
    des clés et \(V\) l'ensemble des valeurs de hachage) était
    injective, alors il existerait une fonction \(g: V \to K\) telle que
    \(g \circ h = \text{Id}_K\), permettant de retrouver la clé
    originale à partir de son hachage. En pratique, les fonctions de
    hachage ne sont pas injectives sur l'ensemble de toutes les entrées
    possibles, mais elles sont conçues pour minimiser les collisions sur
    l'ensemble des entrées attendues.
  \end{itemize}
\item
  \textbf{Cryptographie :}

  \begin{itemize}
  \tightlist
  \item
    Les fonctions injectives sont cruciales en cryptographie. Par
    exemple, une fonction de chiffrement
    \(E: \text{Message Clair} \to \text{Message Chiffré}\) doit être
    injective. Si deux messages clairs différents pouvaient être
    chiffrés en le même message chiffré, il serait impossible de
    déchiffrer de manière unique le message original.
  \item
    La fonction de déchiffrement
    \(D: \text{Message Chiffré} \to \text{Message Clair}\) agit alors
    comme une inverse à gauche (et une inverse à droite) de la fonction
    de chiffrement, assurant que
    \(D \circ E = \text{Id}_{\text{Message Clair}}\).
  \end{itemize}
\end{enumerate}

En somme, la propriété d'injectivité garantit qu'une transformation ne
perd pas d'information, et l'existence d'une inverse à gauche est la
manifestation concrète de cette capacité à ``revenir en arrière'' ou à
``annuler'' la transformation pour les éléments qui ont été
effectivement transformés. C'est un principe fondamental pour la
réversibilité et la fidélité des processus en IA.
\newpage
\section{Exercice 9}
\section{Exercice 9/10 : Surjectivité et existence d'une fonction
inverse à
droite}\label{exercice-910-surjectivituxe9-et-existence-dune-fonction-inverse-uxe0-droite}

\textbf{Jalon 5 : Applications, injections, surjections, bijections et
composition de fonctions}

\textbf{Niveau de difficulté :} \(\star\star\star\star\star\)

\begin{center}\rule{0.5\linewidth}{0.5pt}\end{center}

\subsection{Énoncé Rigoureux}\label{uxe9noncuxe9-rigoureux}

Soient \(E\) et \(F\) deux ensembles non vides. Soit \(f: E \to F\) une
application (ou fonction).

Démontrer l'équivalence suivante :

\[
f \text{ est surjective } \iff \exists g: F \to E \text{ telle que } f \circ g = \text{Id}_F
\]

où \(\text{Id}_F: F \to F\) est l'application identité sur l'ensemble
\(F\), définie par \(\text{Id}_F(y) = y\) pour tout \(y \in F\).

\begin{center}\rule{0.5\linewidth}{0.5pt}\end{center}

\subsection{Analyse de l'Énoncé}\label{analyse-de-luxe9noncuxe9}

Cet exercice fondamental établit un lien direct et profond entre la
propriété de surjectivité d'une application et l'existence d'une
application ``inverse à droite''. C'est une propriété cruciale en
théorie des ensembles et en algèbre, dont la compréhension est
essentielle pour aborder des concepts plus avancés en mathématiques.

\textbf{Définitions clés à rappeler :}

\begin{enumerate}
\def\labelenumi{\arabic{enumi}.}
\item
  \textbf{Surjectivité de \(f\) :} Une application \(f: E \to F\) est
  dite surjective si et seulement si pour tout élément \(y\) de
  l'ensemble d'arrivée \(F\), il existe au moins un élément \(x\) de
  l'ensemble de départ \(E\) tel que \(f(x) = y\). Formellement : \[
  \forall y \in F, \exists x \in E, f(x) = y
  \] Cela signifie que l'image de \(f\), notée
  \(\text{Im}(f) = \{f(x) \mid x \in E\}\), est égale à l'ensemble
  d'arrivée \(F\).
\item
  \textbf{Application identité \(\text{Id}_F\) :} L'application identité
  sur \(F\) est l'application qui renvoie chaque élément de \(F\) sur
  lui-même. \[
  \text{Id}_F: F \to F, \quad y \mapsto y
  \]
\item
  \textbf{Composition d'applications \(f \circ g\) :} Si \(g: F \to E\)
  et \(f: E \to F\), alors \(f \circ g: F \to F\) est l'application
  définie par \((f \circ g)(y) = f(g(y))\) pour tout \(y \in F\).
  L'application \(g\) est appelée un \emph{inverse à droite} de \(f\).
\end{enumerate}

\textbf{L'énoncé nous demande de prouver une équivalence logique, ce qui
implique la démonstration de deux implications distinctes :}

\begin{itemize}
\tightlist
\item
  \textbf{Implication directe (\(\Rightarrow\)) :} Si \(f\) est
  surjective, alors il existe une application \(g: F \to E\) telle que
  \(f \circ g = \text{Id}_F\).

  \begin{itemize}
  \tightlist
  \item
    Cette partie de la preuve nécessite de \emph{construire}
    explicitement une telle application \(g\). La surjectivité de \(f\)
    garantit que pour chaque \(y \in F\), l'ensemble des antécédents
    \(f^{-1}(\{y\}) = \{x \in E \mid f(x) = y\}\) est non vide. Pour
    définir \(g(y)\), nous devrons \emph{choisir} un élément dans cet
    ensemble \(f^{-1}(\{y\})\). Si l'ensemble \(F\) est infini, cette
    ``collection de choix'' simultanés pour chaque \(y \in F\) requiert
    l'invocation de l'\textbf{Axiome du Choix}, un axiome fondamental de
    la théorie des ensembles. C'est ce point qui justifie en grande
    partie la difficulté 5 étoiles de l'exercice.
  \end{itemize}
\item
  \textbf{Implication réciproque (\(\Leftarrow\)) :} S'il existe une
  application \(g: F \to E\) telle que \(f \circ g = \text{Id}_F\),
  alors \(f\) est surjective.

  \begin{itemize}
  \tightlist
  \item
    Cette partie est généralement plus directe et ne nécessite pas
    l'Axiome du Choix. Nous devrons utiliser la propriété
    \(f \circ g = \text{Id}_F\) pour montrer que tout élément de \(F\)
    possède un antécédent par \(f\).
  \end{itemize}
\end{itemize}

\begin{center}\rule{0.5\linewidth}{0.5pt}\end{center}

\subsection{Correction Exhaustive
Pas-à-Pas}\label{correction-exhaustive-pas-uxe0-pas}

Nous allons démontrer l'équivalence en prouvant les deux implications
séparément, en respectant la règle de la ``Zéro Ellipse Mathématique''.

\subsubsection{\texorpdfstring{Partie 1 : Implication directe
(\(\Rightarrow\))}{Partie 1 : Implication directe (\textbackslash Rightarrow)}}\label{partie-1-implication-directe-rightarrow}

\textbf{Hypothèse :} L'application \(f: E \to F\) est surjective.
\textbf{Conclusion à démontrer :} Il existe une application
\(g: F \to E\) telle que \(f \circ g = \text{Id}_F\).

\begin{enumerate}
\def\labelenumi{\arabic{enumi}.}
\item
  \textbf{Analyse de l'hypothèse de surjectivité :} Par définition de la
  surjectivité, pour tout élément \(y\) de l'ensemble d'arrivée \(F\),
  il existe au moins un élément \(x\) de l'ensemble de départ \(E\) tel
  que \(f(x) = y\). Cela signifie que pour chaque \(y \in F\),
  l'ensemble de ses antécédents par \(f\), noté
  \(f^{-1}(\{y\}) = \{x \in E \mid f(x) = y\}\), est non vide. \[
  \forall y \in F, f^{-1}(\{y\}) \neq \emptyset
  \]
\item
  \textbf{Construction de l'application \(g: F \to E\) :} Nous devons
  définir l'image \(g(y)\) pour chaque élément \(y \in F\). Pour chaque
  \(y \in F\), nous savons que l'ensemble \(f^{-1}(\{y\})\) est non
  vide. Nous devons choisir un élément spécifique dans cet ensemble pour
  qu'il soit la valeur de \(g(y)\).

  \begin{itemize}
  \item
    \textbf{Si \(F\) est un ensemble fini :} Supposons que
    \(F = \{y_1, y_2, \dots, y_n\}\). Pour chaque \(y_i \in F\),
    l'ensemble \(f^{-1}(\{y_i\})\) est non vide. Nous pouvons alors
    choisir un élément \(x_i \in E\) tel que \(f(x_i) = y_i\). Nous
    définissons ensuite l'application \(g: F \to E\) en posant
    \(g(y_i) = x_i\) pour chaque \(i \in \{1, \dots, n\}\).
  \item
    \textbf{Si \(F\) est un ensemble infini :} Dans ce cas, nous avons
    une famille potentiellement infinie d'ensembles non vides
    \((f^{-1}(\{y\}))_{y \in F}\). Pour pouvoir effectuer un choix
    simultané d'un élément dans chacun de ces ensembles, nous faisons
    appel à l'\textbf{Axiome du Choix}. L'Axiome du Choix stipule que
    pour toute famille non vide d'ensembles non vides
    \((A_i)_{i \in I}\), il existe une fonction de choix
    \(c: I \to \bigcup_{i \in I} A_i\) telle que \(c(i) \in A_i\) pour
    tout \(i \in I\). Dans notre contexte, l'ensemble d'indices est
    \(I = F\), et pour chaque \(y \in F\), l'ensemble correspondant est
    \(A_y = f^{-1}(\{y\})\). L'Axiome du Choix garantit donc l'existence
    d'une fonction \(g: F \to E\) telle que pour tout \(y \in F\),
    \(g(y) \in f^{-1}(\{y\})\). Par définition de l'ensemble
    \(f^{-1}(\{y\})\), la condition \(g(y) \in f^{-1}(\{y\})\) signifie
    précisément que \(f(g(y)) = y\) pour tout \(y \in F\).
  \end{itemize}
\item
  \textbf{Vérification de la propriété \(f \circ g = \text{Id}_F\) :}
  Nous avons construit une application \(g: F \to E\) telle que pour
  tout \(y \in F\), la relation \(f(g(y)) = y\) est satisfaite. Par
  définition de la composition d'applications, pour tout \(y \in F\),
  l'image de \(y\) par \(f \circ g\) est donnée par
  \((f \circ g)(y) = f(g(y))\). Par définition de l'application identité
  sur \(F\), pour tout \(y \in F\), l'image de \(y\) par \(\text{Id}_F\)
  est donnée par \(\text{Id}_F(y) = y\). En combinant ces deux points,
  nous obtenons que pour tout \(y \in F\): \[
  (f \circ g)(y) = f(g(y)) = y = \text{Id}_F(y)
  \] Puisque cette égalité est vraie pour tout \(y \in F\), nous pouvons
  conclure que les applications \(f \circ g\) et \(\text{Id}_F\) sont
  égales. \[
  f \circ g = \text{Id}_F
  \]
\end{enumerate}

Ainsi, si \(f\) est surjective, il existe bien une application
\(g: F \to E\) telle que \(f \circ g = \text{Id}_F\).

\subsubsection{\texorpdfstring{Partie 2 : Implication réciproque
(\(\Leftarrow\))}{Partie 2 : Implication réciproque (\textbackslash Leftarrow)}}\label{partie-2-implication-ruxe9ciproque-leftarrow}

\textbf{Hypothèse :} Il existe une application \(g: F \to E\) telle que
\(f \circ g = \text{Id}_F\). \textbf{Conclusion à démontrer :}
L'application \(f: E \to F\) est surjective.

\begin{enumerate}
\def\labelenumi{\arabic{enumi}.}
\item
  \textbf{Analyse de l'hypothèse :} Nous avons une application
  \(g: F \to E\) telle que la composition \(f \circ g\) est égale à
  l'application identité sur \(F\). Par définition de l'égalité
  d'applications, cela signifie que pour tout élément \(y \in F\),
  l'image de \(y\) par \(f \circ g\) est égale à l'image de \(y\) par
  \(\text{Id}_F\). \[
  \forall y \in F, (f \circ g)(y) = \text{Id}_F(y)
  \] En utilisant les définitions de la composition et de l'identité,
  ceci se traduit par : \[
  \forall y \in F, f(g(y)) = y
  \]
\item
  \textbf{Démonstration de la surjectivité de \(f\) :} Pour montrer que
  \(f\) est surjective, nous devons prouver que pour tout élément \(y\)
  de l'ensemble d'arrivée \(F\), il existe au moins un élément \(x\) de
  l'ensemble de départ \(E\) tel que \(f(x) = y\).

  \begin{itemize}
  \tightlist
  \item
    Soit \(y \in F\) un élément arbitraire, mais fixé, de l'ensemble
    d'arrivée \(F\).
  \item
    Notre objectif est de trouver un antécédent \(x \in E\) pour ce
    \(y\) sous l'application \(f\).
  \item
    Considérons l'élément \(g(y)\). Puisque \(g\) est une application de
    \(F\) vers \(E\), l'image \(g(y)\) est nécessairement un élément de
    l'ensemble \(E\).
  \item
    Posons \(x_0 = g(y)\). Nous avons donc \(x_0 \in E\).
  \item
    Calculons l'image de cet \(x_0\) par l'application \(f\): \[
    f(x_0) = f(g(y))
    \]
  \item
    D'après notre hypothèse (analysée au point 1), nous savons que
    \(f(g(y)) = y\).
  \item
    Par conséquent, nous avons \(f(x_0) = y\).
  \item
    Nous avons ainsi démontré que pour tout \(y \in F\) (puisque \(y\)
    était arbitraire), il existe un élément \(x_0 \in E\) (à savoir
    \(g(y)\)) tel que \(f(x_0) = y\).
  \end{itemize}
\item
  \textbf{Conclusion :} Par définition, l'application \(f\) est
  surjective.
\end{enumerate}

Ainsi, s'il existe une application \(g: F \to E\) telle que
\(f \circ g = \text{Id}_F\), alors \(f\) est surjective.

\subsubsection{Synthèse}\label{synthuxe8se}

Les deux implications ayant été démontrées avec rigueur, nous pouvons
conclure que l'équivalence est vraie :

\[
f \text{ est surjective } \iff \exists g: F \to E \text{ telle que } f \circ g = \text{Id}_F
\]

\begin{center}\rule{0.5\linewidth}{0.5pt}\end{center}

\subsection{Liens avec l'Intelligence
Artificielle}\label{liens-avec-lintelligence-artificielle}

Bien que cet exercice soit de nature purement mathématique et
fondamentale, les concepts de surjectivité et d'existence d'un inverse à
droite trouvent des échos et des applications indirectes, mais
significatives, dans divers domaines de l'Intelligence Artificielle,
notamment dans la conception et l'analyse de modèles.

\begin{enumerate}
\def\labelenumi{\arabic{enumi}.}
\tightlist
\item
  \textbf{Modèles Génératifs (GANs, VAEs) et la ``Manifold Hypothesis''
  :}

  \begin{itemize}
  \tightlist
  \item
    Dans les modèles génératifs (comme les Generative Adversarial
    Networks ou les Variational Autoencoders), l'objectif est
    d'apprendre une application \(G: Z \to X\), où \(Z\) est un espace
    latent (souvent un espace de faible dimension, comme
    \(\mathbb{R}^d\)) et \(X\) est l'espace des données réelles (par
    exemple, des images de haute dimension).
  \item
    Idéalement, on voudrait que le générateur \(G\) soit ``surjectif''
    sur la \emph{variété des données réelles} (data manifold). Si \(G\)
    n'est pas surjective sur cette variété, cela signifie qu'il existe
    des données réelles valides et réalistes que le modèle ne pourra
    jamais générer. Un générateur parfaitement surjectif sur la variété
    des données serait capable de produire n'importe quelle donnée
    réaliste.
  \item
    L'existence d'un ``inverse à droite'' \(g: X \to Z\) (souvent appelé
    \emph{encodeur} ou \emph{inférence inverse}) tel que
    \(G \circ g = \text{Id}_X\) (sur la variété des données)
    signifierait qu'on peut prendre une donnée réelle \(x\), la mapper à
    un point latent \(z=g(x)\), et que le générateur \(G\) peut
    parfaitement reconstruire \(x\) à partir de \(z\). C'est précisément
    l'objectif des \textbf{auto-encodeurs}, où l'encodeur \(g\) et le
    décodeur \(f\) (ici \(G\)) sont entraînés pour minimiser l'erreur de
    reconstruction \(||x - f(g(x))||\). Un auto-encodeur parfait serait
    un exemple concret de cette relation \(f \circ g = \text{Id}_X\).
  \end{itemize}
\item
  \textbf{Systèmes de Contrôle et Robotique (Cinématique Inverse) :}

  \begin{itemize}
  \tightlist
  \item
    En robotique, la cinématique directe est une fonction
    \(f: \Theta \to \mathcal{P}\) qui mappe un ensemble d'angles
    articulaires \(\Theta\) (l'espace de départ \(E\)) à une pose
    d'effecteur final \(\mathcal{P}\) (l'espace d'arrivée \(F\)).
  \item
    Si un robot peut atteindre n'importe quelle pose dans son espace de
    travail (c'est-à-dire que \(f\) est surjective sur l'espace de
    travail), alors il existe une fonction de cinématique inverse
    \(g: \mathcal{P} \to \Theta\) (un ``inverse à droite'') qui, pour
    une pose désirée \(p \in \mathcal{P}\), fournit un ensemble d'angles
    articulaires \(\theta = g(p)\) tel que \(f(\theta) = p\).
    C'est-à-dire \(f \circ g = \text{Id}_{\mathcal{P}}\). L'existence
    d'une telle fonction \(g\) est cruciale pour la planification de
    mouvement et le contrôle des robots, permettant de traduire une
    tâche dans l'espace opérationnel en commandes articulaires.
  \end{itemize}
\item
  \textbf{Apprentissage par Renforcement (Reinforcement Learning) :}

  \begin{itemize}
  \tightlist
  \item
    Dans certains contextes, on peut modéliser la politique d'un agent
    comme une fonction \(f: S \to A\) (état vers action) ou
    \(f: S \to \mathcal{D}(A)\) (état vers distribution de probabilité
    sur les actions).
  \item
    La surjectivité de \(f\) pourrait signifier que l'agent peut
    potentiellement prendre n'importe quelle action dans n'importe quel
    état (ou générer n'importe quelle distribution d'actions), ce qui
    est souvent une hypothèse de base pour l'exploration. L'existence
    d'un inverse à droite pourrait être liée à la capacité de
    ``reconstruire'' l'état à partir de l'action prise, bien que ce soit
    une analogie plus lâche et contextuelle.
  \end{itemize}
\item
  \textbf{Compression de Données et Hachage :}

  \begin{itemize}
  \tightlist
  \item
    Les fonctions de hachage \(h: D \to H\) (données vers hachage) ne
    sont généralement pas surjectives sur l'espace de hachage complet,
    ni injectives. Cependant, l'idée d'un ``inverse à droite'' peut être
    conceptualisée dans des systèmes de compression réversibles où l'on
    peut décompresser parfaitement. Si une fonction de compression
    \(C: X \to Y\) est surjective sur l'ensemble des données compressées
    valides \(Y\), alors il existe une fonction de décompression
    \(D: Y \to X\) telle que \(C \circ D = \text{Id}_Y\). Cela signifie
    que toute donnée compressée valide peut être décompressée pour
    retrouver une donnée originale, et que la compression de cette
    donnée originale donne la donnée compressée de départ.
  \end{itemize}
\end{enumerate}

En résumé, la compréhension de la surjectivité et de l'existence
d'inverses à droite est fondamentale pour analyser les capacités et les
limites des modèles d'IA, en particulier ceux qui impliquent des
transformations d'espaces (comme les auto-encodeurs, les générateurs) ou
des mappings entre différentes représentations. Elle permet de raisonner
sur la complétude d'un modèle à couvrir son espace cible ou sa capacité
à être ``inversé'' pour certaines opérations.
\newpage
\section{Exercice 10}
\section{Exercice 10/10 : Le Théorème de
Cantor}\label{exercice-1010-le-thuxe9oruxe8me-de-cantor}

\textbf{Jalon 5 : Applications, injections, surjections, bijections et
composition de fonctions}

\textbf{Niveau de difficulté :}
\(\star\)\(\star\)\(\star\)\(\star\)\(\star\)

\begin{center}\rule{0.5\linewidth}{0.5pt}\end{center}

\subsection{Énoncé de l'Exercice}\label{uxe9noncuxe9-de-lexercice}

Soit \(E\) un ensemble non vide. On note \(\mathcal{P}(E)\) l'ensemble
de tous les sous-ensembles de \(E\), appelé l'ensemble des parties de
\(E\).

Démontrer, en utilisant l'argument diagonal de Cantor, qu'il n'existe
aucune fonction surjective de \(E\) vers \(\mathcal{P}(E)\).

Autrement dit, pour toute fonction \(f: E \to \mathcal{P}(E)\), il
existe au moins un sous-ensemble \(Y \in \mathcal{P}(E)\) tel que pour
tout \(x \in E\), \(f(x) \neq Y\).

\begin{center}\rule{0.5\linewidth}{0.5pt}\end{center}

\subsection{Analyse de l'Énoncé}\label{analyse-de-luxe9noncuxe9}

Cet exercice aborde un résultat fondamental de la théorie des ensembles,
le théorème de Cantor, qui a des implications profondes sur la notion de
taille (cardinalité) des ensembles, en particulier pour les ensembles
infinis.

\begin{enumerate}
\def\labelenumi{\arabic{enumi}.}
\tightlist
\item
  \textbf{Ensemble \(E\) :} Il s'agit d'un ensemble arbitraire,
  potentiellement fini ou infini. La démonstration est valable dans tous
  les cas.
\item
  \textbf{Ensemble \(\mathcal{P}(E)\) :} C'est l'ensemble de tous les
  sous-ensembles de \(E\).

  \begin{itemize}
  \tightlist
  \item
    Si \(E = \{1, 2\}\), alors
    \(\mathcal{P}(E) = \{\emptyset, \{1\}, \{2\}, \{1, 2\}\}\).
  \item
    Si \(E\) a \(n\) éléments, \(\mathcal{P}(E)\) a \(2^n\) éléments.
  \item
    Pour \(n \ge 1\), on a \(2^n > n\). Cela suggère déjà que
    \(\mathcal{P}(E)\) est ``plus grand'' que \(E\) pour les ensembles
    finis. Le théorème de Cantor généralise cette intuition aux
    ensembles infinis.
  \end{itemize}
\item
  \textbf{Fonction surjective \(f: E \to \mathcal{P}(E)\) :} Une
  fonction \(f\) est surjective si chaque élément de l'ensemble
  d'arrivée (\(\mathcal{P}(E)\) ici) est l'image d'au moins un élément
  de l'ensemble de départ (\(E\) ici).

  \begin{itemize}
  \tightlist
  \item
    L'énoncé nous demande de prouver qu'une telle surjection
    \emph{n'existe pas}.
  \item
    Cela signifie que, quelle que soit la fonction \(f\) que l'on
    choisit de \(E\) vers \(\mathcal{P}(E)\), on doit pouvoir trouver un
    sous-ensemble \(Y \in \mathcal{P}(E)\) qui n'est l'image d'aucun
    élément de \(E\) par \(f\). C'est-à-dire, \(Y \notin \text{Im}(f)\).
  \end{itemize}
\item
  \textbf{L'argument diagonal de Cantor :} C'est la technique de preuve
  clé. Elle consiste à construire un élément ``spécial''
  \(Y \in \mathcal{P}(E)\) qui est défini de manière à différer de
  \emph{toutes} les images \(f(x)\) pour \emph{tous} les \(x \in E\).
  L'idée est de créer un ``point de désaccord'' avec chaque \(f(x)\).
\end{enumerate}

Ce théorème est d'une importance capitale car il montre qu'il existe
différents ``niveaux'' d'infini. Par exemple, l'ensemble des nombres
réels \(\mathbb{R}\) a la même cardinalité que
\(\mathcal{P}(\mathbb{N})\) (l'ensemble des parties des nombres
naturels), et le théorème de Cantor implique donc que
\(|\mathbb{N}| < |\mathbb{R}|\).

\begin{center}\rule{0.5\linewidth}{0.5pt}\end{center}

\subsection{Correction Exhaustive
Pas-à-Pas}\label{correction-exhaustive-pas-uxe0-pas}

Nous allons procéder par l'absurde.

\textbf{Étape 1 : Supposition par l'absurde}

Supposons qu'il existe une fonction \(f: E \to \mathcal{P}(E)\) qui est
surjective. Par définition de la surjectivité, cela signifie que pour
tout \(Y' \in \mathcal{P}(E)\), il existe au moins un \(x' \in E\) tel
que \(f(x') = Y'\).

\textbf{Étape 2 : Construction de l'ensemble diagonal de Cantor}

Nous allons maintenant construire un sous-ensemble
\(Y \in \mathcal{P}(E)\) qui sera la contradiction recherchée.
Définissons l'ensemble \(Y\) comme suit :
\[ Y = \{x \in E \mid x \notin f(x)\} \]

Analysons cette définition : * Pour chaque élément \(x \in E\), \(f(x)\)
est un sous-ensemble de \(E\), c'est-à-dire \(f(x) \in \mathcal{P}(E)\).
* La condition ``\(x \notin f(x)\)'' est une proposition logique qui est
soit vraie, soit fausse pour chaque \(x \in E\). * L'ensemble \(Y\) est
donc l'ensemble de tous les éléments \(x\) de \(E\) qui n'appartiennent
\emph{pas} à l'image d'eux-mêmes par \(f\). * Par l'axiome de
compréhension, l'ensemble \(Y\) est formé d'éléments \(x\) qui
appartiennent préalablement à \(E\). Ainsi, pour tout \(x \in Y\), nous
avons logiquement \(x \in E\), ce qui caractérise l'inclusion
\(Y \subseteq E\). Par conséquent, par définition de l'ensemble des
parties, \(Y \in \mathcal{P}(E)\).

\textbf{Étape 3 : Utilisation de l'hypothèse de surjectivité}

Puisque nous avons supposé que \(f\) est surjective et que
\(Y \in \mathcal{P}(E)\), il doit exister un élément \(a \in E\) tel que
\(f(a) = Y\). Ceci découle directement de la définition de la
surjectivité : chaque élément de l'ensemble d'arrivée
(\(\mathcal{P}(E)\)) doit être l'image d'au moins un élément de
l'ensemble de départ (\(E\)). Puisque \(Y\) est un élément de
\(\mathcal{P}(E)\), il doit avoir un antécédent \(a\) dans \(E\).

\textbf{Étape 4 : Dérivation de la contradiction}

Considérons l'élément \(a \in E\) que nous avons trouvé à l'Étape 3, tel
que \(f(a) = Y\). Nous allons maintenant nous poser la question cruciale
: l'élément \(a\) appartient-il à l'ensemble \(Y\) ?

Nous avons deux cas possibles, qui sont mutuellement exclusifs et
exhaustifs :

\textbf{Cas 1 : Supposons que \(a \in Y\).} 1. Si \(a \in Y\), alors,
par la définition de l'ensemble \(Y\) (Étape 2), tout élément
\(z \in Y\) satisfait la condition \(z \notin f(z)\). 2. En appliquant
cette définition à \(a\), nous obtenons \(a \notin f(a)\). 3. Cependant,
nous savons de l'Étape 3 que \(f(a) = Y\). 4. En substituant \(f(a)\)
par \(Y\) dans l'expression \(a \notin f(a)\), nous obtenons
\(a \notin Y\). 5. Nous sommes arrivés à une contradiction : nous avons
supposé \(a \in Y\) et nous avons déduit \(a \notin Y\). Ceci est
logiquement impossible.

\textbf{Cas 2 : Supposons que \(a \notin Y\).} 1. Si \(a \notin Y\),
alors, par la définition de l'ensemble \(Y\) (Étape 2), tout élément
\(z \in E\) qui n'est pas dans \(Y\) doit satisfaire la négation de la
condition de \(Y\), c'est-à-dire \(z \in f(z)\). 2. En appliquant cette
logique à \(a\), nous obtenons \(a \in f(a)\). 3. Cependant, nous savons
de l'Étape 3 que \(f(a) = Y\). 4. En substituant \(f(a)\) par \(Y\) dans
l'expression \(a \in f(a)\), nous obtenons \(a \in Y\). 5. Nous sommes
arrivés à une contradiction : nous avons supposé \(a \notin Y\) et nous
avons déduit \(a \in Y\). Ceci est également logiquement impossible.

\textbf{Étape 5 : Conclusion}

Dans les deux cas possibles (que \(a \in Y\) ou \(a \notin Y\)), nous
aboutissons à une contradiction logique. Puisque notre raisonnement est
correct et que les deux cas couvrent toutes les possibilités pour \(a\),
la seule conclusion possible est que notre supposition initiale (Étape
1) doit être fausse.

Par conséquent, il n'existe aucune fonction surjective
\(f: E \to \mathcal{P}(E)\). Ceci démontre le théorème de Cantor.

\begin{center}\rule{0.5\linewidth}{0.5pt}\end{center}

\subsection{Liens avec l'Intelligence
Artificielle}\label{liens-avec-lintelligence-artificielle}

Le théorème de Cantor, bien que fondamentalement mathématique, a des
résonances profondes et des analogies frappantes dans le domaine de
l'Intelligence Artificielle et de l'informatique théorique :

\begin{enumerate}
\def\labelenumi{\arabic{enumi}.}
\item
  \textbf{Le Problème de l'Arrêt (Halting Problem) :} C'est l'analogie
  la plus directe et la plus célèbre. Le problème de l'arrêt, qui
  demande s'il est possible de déterminer pour un programme arbitraire
  et une entrée arbitraire si le programme finira par s'arrêter ou
  s'exécutera indéfiniment, a été prouvé indécidable par Alan Turing en
  utilisant un argument diagonal très similaire à celui de Cantor.

  \begin{itemize}
  \tightlist
  \item
    Imaginez une fonction \(f\) qui, pour chaque programme \(P\)
    (élément de \(E\)), renvoie un ensemble de programmes \(f(P)\) avec
    lesquels \(P\) interagit d'une certaine manière.
  \item
    L'argument diagonal construit un programme ``paradoxal'' qui, si on
    lui applique la logique de \(f\), mène à une contradiction, prouvant
    ainsi qu'une telle \(f\) (qui résoudrait le problème de l'arrêt) ne
    peut exister.
  \item
    Cela établit une limite fondamentale à ce que les ordinateurs
    peuvent calculer, même en théorie.
  \end{itemize}
\item
  \textbf{Limites de la Représentation et de l'Énumération :}

  \begin{itemize}
  \tightlist
  \item
    Le théorème de Cantor montre que l'ensemble des sous-ensembles d'un
    ensemble est ``plus grand'' que l'ensemble lui-même. En IA, cela
    peut être interprété comme une limite à la capacité de représenter
    toutes les ``idées'' ou ``concepts'' (qui pourraient être vus comme
    des sous-ensembles d'un espace de base) en utilisant un système de
    symboles ou un langage formel qui a la même ``taille'' que l'espace
    de base.
  \item
    Par exemple, l'ensemble des fonctions de \(\mathbb{N}\) vers
    \(\mathbb{N}\) est non-dénombrable (a la même cardinalité que
    \(\mathcal{P}(\mathbb{N})\)), tandis que l'ensemble des programmes
    informatiques est dénombrable. Cela signifie qu'il existe infiniment
    plus de fonctions que de programmes pour les calculer. La plupart
    des fonctions ne sont donc pas calculables par un algorithme. C'est
    une limite fondamentale pour l'IA symbolique et algorithmique.
  \end{itemize}
\item
  \textbf{Théorèmes d'Incomplétude de Gödel :} Bien que plus complexes,
  les théorèmes d'incomplétude de Gödel, qui démontrent les limites des
  systèmes axiomatiques formels, utilisent également des techniques de
  ``référence à soi-même'' et des constructions de type diagonal. Ils
  montrent qu'il y aura toujours des énoncés vrais qui ne peuvent pas
  être prouvés au sein d'un système formel suffisamment puissant pour
  contenir l'arithmétique. Ces limites sont cruciales pour comprendre
  les fondements de la logique et de la connaissance, des domaines
  centraux pour l'IA.
\item
  \textbf{Complexité et Expressivité des Modèles :} Dans l'apprentissage
  automatique, nous construisons des modèles (par exemple, des réseaux
  de neurones) pour apprendre des fonctions. Le théorème de Cantor, à un
  niveau très abstrait, nous rappelle que l'espace des fonctions
  possibles est immensément vaste. Même si nos modèles sont très
  expressifs, ils ne peuvent ``représenter'' qu'une infime fraction de
  toutes les fonctions possibles. Cela souligne l'importance de la
  régularisation, des biais inductifs et de la recherche d'architectures
  qui capturent les propriétés pertinentes des fonctions que nous
  voulons apprendre, plutôt que de tenter de couvrir l'espace entier.
\end{enumerate}

En somme, le théorème de Cantor est un rappel puissant des limites
fondamentales de la logique, de la calculabilité et de la
représentation, des concepts qui sont au cœur des défis théoriques et
pratiques de l'Intelligence Artificielle.
\newpage
\chapter{Travaux Pratiques}
\section{TP 1}
\section{TP 1/5 - Jalon 5 : Injectivité et Surjectivité sur Ensembles
Finis}\label{tp-15---jalon-5-injectivituxe9-et-surjectivituxe9-sur-ensembles-finis}

\subsection{Introduction}\label{introduction}

Ce premier Travail Pratique du Jalon 5 se concentre sur la compréhension
et l'implémentation des concepts fondamentaux d'applications
(fonctions), d'injectivité et de surjectivité dans le contexte des
ensembles finis. En mathématiques, une application \(f\) d'un ensemble
\(A\) (domaine) vers un ensemble \(B\) (codomaine) associe à chaque
élément de \(A\) un unique élément de \(B\).

Nous allons représenter ces applications à l'aide de dictionnaires
Python, où les clés représenteront les éléments du domaine \(A\) et les
valeurs les éléments correspondants du codomaine \(B\). Cette approche
nous permettra de manipuler des ensembles finis et de vérifier les
propriétés d'injectivité et de surjectivité de manière algorithmique.

\subsubsection{Rappels Mathématiques}\label{rappels-mathuxe9matiques}

Soit une application \(f: A \to B\).

\begin{itemize}
\item
  \textbf{Injectivité} : L'application \(f\) est dite \textbf{injective}
  si et seulement si deux éléments distincts du domaine \(A\) ont
  toujours des images distinctes dans le codomaine \(B\). Formellement :
  \(\forall x_1, x_2 \in A, (x_1 \neq x_2 \implies f(x_1) \neq f(x_2))\).
  Une formulation équivalente est :
  \(\forall x_1, x_2 \in A, (f(x_1) = f(x_2) \implies x_1 = x_2)\). En
  d'autres termes, chaque élément du codomaine est l'image d'au plus un
  élément du domaine.
\item
  \textbf{Surjectivité} : L'application \(f\) est dite
  \textbf{surjective} si et seulement si chaque élément du codomaine
  \(B\) est l'image d'au moins un élément du domaine \(A\). Formellement
  : \(\forall y \in B, \exists x \in A \text{ tel que } f(x) = y\). Cela
  signifie que l'image de \(f\), notée
  \(Im(f) = \{f(x) \mid x \in A\}\), est égale au codomaine \(B\).
\item
  \textbf{Bijectivité} : L'application \(f\) est dite \textbf{bijective}
  si et seulement si elle est à la fois injective et surjective.
\end{itemize}

\subsection{Objectifs du TP}\label{objectifs-du-tp}

\begin{enumerate}
\def\labelenumi{\arabic{enumi}.}
\tightlist
\item
  Représenter une application \(f: A \to B\) à l'aide d'un dictionnaire
  Python.
\item
  Implémenter une fonction Python
  \passthrough{\lstinline!is\_injective(f\_dict)!} qui vérifie si une
  application donnée est injective.
\item
  Implémenter une fonction Python
  \passthrough{\lstinline!is\_surjective(f\_dict, codomain\_set)!} qui
  vérifie si une application donnée est surjective.
\item
  (Bonus) Implémenter une fonction Python
  \passthrough{\lstinline!is\_bijective(f\_dict, codomain\_set)!} qui
  vérifie si une application donnée est bijective.
\item
  Utiliser des assertions pour valider la justesse mathématique des
  implémentations.
\end{enumerate}

\subsection{Prérequis}\label{pruxe9requis}

\begin{itemize}
\tightlist
\item
  Connaissances de base en Python (dictionnaires, ensembles, boucles,
  conditions).
\item
  Compréhension des concepts d'ensembles, de fonctions, d'injectivité et
  de surjectivité.
\end{itemize}

\begin{center}\rule{0.5\linewidth}{0.5pt}\end{center}

\subsection{Partie 1 : Représentation d'une
fonction}\label{partie-1-repruxe9sentation-dune-fonction}

Une fonction \(f: A \to B\) peut être représentée par un dictionnaire où
les clés sont les éléments du domaine \(A\) et les valeurs sont leurs
images dans \(B\). Il est crucial de noter que le dictionnaire ne
contient que les paires \((x, f(x))\) pour \(x \in A\). Le codomaine
\(B\) doit être fourni séparément pour vérifier la surjectivité, car le
dictionnaire seul ne permet de déduire que l'image de la fonction, pas
l'ensemble \(B\) complet.

\textbf{Exemple de représentation :} Si \(A = \{1, 2, 3\}\) et
\(B = \{a, b, c, d\}\), et \(f(1)=a, f(2)=b, f(3)=c\), alors la fonction
peut être représentée par
\passthrough{\lstinline!f\_dict = \{1: 'a', 2: 'b', 3: 'c'\}!}. Le
codomaine serait
\passthrough{\lstinline!codomain\_set = \{'a', 'b', 'c', 'd'\}!}.

\begin{lstlisting}[language=Python]
# --- Partie 1 : Representation d'une fonction ---

def get_domain(f_dict: dict) -> set:
    """
    Extrait le domaine (ensemble de depart) d'une fonction representee par un dictionnaire.
    Le domaine est l'ensemble de toutes les cles du dictionnaire.

    Args:
        f_dict (dict): Dictionnaire representant la fonction {x: f(x)}.

    Returns:
        set: L'ensemble des elements du domaine.
    """
    return set(f_dict.keys())

def get_image(f_dict: dict) -> set:
    """
    Extrait l'image (ensemble d'arrivee des valeurs reellement atteintes) d'une fonction.
    L'image est l'ensemble de toutes les valeurs du dictionnaire.

    Args:
        f_dict (dict): Dictionnaire representant la fonction {x: f(x)}.

    Returns:
        set: L'ensemble des elements de l'image.
    """
    return set(f_dict.values())

# Exemple de fonction
f_example = {
    'pomme': 'fruit',
    'carotte': 'legume',
    'banane': 'fruit',
    'brocoli': 'legume'
}

# Definition explicite du codomaine pour les tests de surjectivite
codomain_example = {'fruit', 'legume', 'cereale'}

print(f"Fonction f_example: {f_example}")
print(f"Domaine de f_example: {get_domain(f_example)}")
print(f"Image de f_example: {get_image(f_example)}")
print(f"Codomaine explicite pour f_example: {codomain_example}")

# Assertions pour valider les fonctions utilitaires
assert get_domain(f_example) == {'pomme', 'carotte', 'banane', 'brocoli'}
assert get_image(f_example) == {'fruit', 'legume'}
print("\nAssertions pour les fonctions utilitaires passees avec succes.")
\end{lstlisting}

\subsection{Partie 2 : Vérification de
l'injectivité}\label{partie-2-vuxe9rification-de-linjectivituxe9}

Pour vérifier l'injectivité d'une fonction \(f: A \to B\), nous devons
nous assurer qu'aucun élément du codomaine n'est l'image de plus d'un
élément du domaine. En d'autres termes, si nous collectons toutes les
images \(f(x)\) pour \(x \in A\), il ne doit pas y avoir de doublons
parmi ces images.

\textbf{Algorithme :} 1. Extraire toutes les valeurs (images) de la
fonction représentée par le dictionnaire. 2. Convertir cette collection
de valeurs en un ensemble. Un ensemble ne contient pas de doublons. 3.
Si la taille de la collection originale de valeurs est égale à la taille
de l'ensemble des valeurs (c'est-à-dire qu'il n'y a pas eu de
suppression de doublons), alors la fonction est injective. Sinon, elle
ne l'est pas.

\begin{lstlisting}[language=Python]
# --- Partie 2 : Verification de l'injectivite ---

def is_injective(f_dict: dict) -> bool:
    """
    Verifie si une fonction representee par un dictionnaire est injective.
    Une fonction est injective si chaque element du codomaine est l'image
    d'au plus un element du domaine.
    Algorithmiquement, cela signifie que toutes les valeurs (images) dans le dictionnaire
    doivent etre uniques.

    Args:
        f_dict (dict): Dictionnaire representant la fonction {x: f(x)}.

    Returns:
        bool: True si la fonction est injective, False sinon.
    """
    # Etape 1: Obtenir toutes les valeurs (images) de la fonction.
    # Les valeurs peuvent contenir des doublons si la fonction n'est pas injective.
    values = list(f_dict.values())

    # Etape 2: Creer un ensemble a partir de ces valeurs.
    # Un ensemble elimine automatiquement les doublons.
    unique_values = set(values)

    # Etape 3: Comparer la taille de la liste originale des valeurs avec la taille de l'ensemble.
    # Si les tailles sont egales, cela signifie qu'il n'y avait pas de doublons,
    # donc chaque element du domaine a une image unique.
    return len(values) == len(unique_values)

# --- Tests et Assertions pour l'injectivite ---

print("\n--- Tests d'injectivite ---")

# Fonction injective (chaque element du domaine a une image unique)
f_injective = {'a': 1, 'b': 2, 'c': 3}
# Mathematiquement: f(a)=1, f(b)=2, f(c)=3. Toutes les images sont distinctes.
print(f"f_injective: {f_injective} est injective? {is_injective(f_injective)}")
assert is_injective(f_injective) is True, "f_injective devrait etre injective"

# Fonction non injective (deux elements du domaine ont la meme image)
f_not_injective = {'a': 1, 'b': 2, 'c': 1}
# Mathematiquement: f(a)=1 et f(c)=1, mais a != c. Donc non injective.
print(f"f_not_injective: {f_not_injective} est injective? {is_injective(f_not_injective)}")
assert is_injective(f_not_injective) is False, "f_not_injective ne devrait PAS etre injective"

# Fonction avec un seul element (toujours injective)
f_single_element = {'x': 10}
print(f"f_single_element: {f_single_element} est injective? {is_injective(f_single_element)}")
assert is_injective(f_single_element) is True, "f_single_element devrait etre injective"

# Fonction vide (toujours injective par definition, car la condition n'est jamais violee)
f_empty = {}
print(f"f_empty: {f_empty} est injective? {is_injective(f_empty)}")
assert is_injective(f_empty) is True, "f_empty devrait etre injective"

print("Assertions d'injectivite passees avec succes.")
\end{lstlisting}

\subsection{Partie 3 : Vérification de la
surjectivité}\label{partie-3-vuxe9rification-de-la-surjectivituxe9}

Pour vérifier la surjectivité d'une fonction \(f: A \to B\), nous devons
nous assurer que chaque élément du codomaine \(B\) est atteint par au
moins un élément du domaine \(A\). Cela signifie que l'ensemble des
images de la fonction (l'image de \(f\)) doit être exactement égal au
codomaine \(B\).

\textbf{Algorithme :} 1. Extraire l'ensemble des images de la fonction
(les valeurs du dictionnaire). 2. Comparer cet ensemble d'images avec
l'ensemble du codomaine \(B\) fourni explicitement. 3. Si les deux
ensembles sont égaux, la fonction est surjective. Sinon, elle ne l'est
pas.

\begin{lstlisting}[language=Python]
# --- Partie 3 : Verification de la surjectivite ---

def is_surjective(f_dict: dict, codomain_set: set) -> bool:
    """
    Verifie si une fonction representee par un dictionnaire est surjective par rapport
    a un codomaine donne.
    Une fonction est surjective si chaque element du codomaine est l'image
    d'au moins un element du domaine.
    Algorithmiquement, cela signifie que l'ensemble des images de la fonction
    doit etre egal a l'ensemble du codomaine.

    Args:
        f_dict (dict): Dictionnaire representant la fonction {x: f(x)}.
        codomain_set (set): L'ensemble explicite du codomaine B.

    Returns:
        bool: True si la fonction est surjective, False sinon.
    """
    # Etape 1: Obtenir l'ensemble des images de la fonction.
    # C'est l'ensemble de toutes les valeurs presentes dans le dictionnaire.
    image_set = set(f_dict.values())

    # Etape 2: Comparer l'ensemble des images avec l'ensemble du codomaine.
    # Si l'image est egale au codomaine, alors tous les elements du codomaine
    # sont atteints par la fonction.
    return image_set == codomain_set

# --- Tests et Assertions pour la surjectivite ---

print("\n--- Tests de surjectivite ---")

# Fonction surjective (l'image couvre tout le codomaine)
f_surjective = {'a': 1, 'b': 2, 'c': 3}
codomain_123 = {1, 2, 3}
# Mathematiquement: Im(f) = {1,2,3}. Codomaine = {1,2,3}. Im(f) == Codomaine. Surjective.
print(f"f_surjective: {f_surjective}, Codomaine: {codomain_123} est surjective? {is_surjective(f_surjective, codomain_123)}")
assert is_surjective(f_surjective, codomain_123) is True, "f_surjective devrait etre surjective"

# Fonction non surjective (le codomaine contient des elements non atteints)
f_not_surjective = {'a': 1, 'b': 2, 'c': 2}
codomain_123 = {1, 2, 3}
# Mathematiquement: Im(f) = {1,2}. Codomaine = {1,2,3}. Im(f) != Codomaine (3 n'est pas atteint). Non surjective.
print(f"f_not_surjective: {f_not_surjective}, Codomaine: {codomain_123} est surjective? {is_surjective(f_not_surjective, codomain_123)}")
assert is_surjective(f_not_surjective, codomain_123) is False, "f_not_surjective ne devrait PAS etre surjective"

# Fonction surjective (meme si non injective)
f_surjective_not_injective = {'a': 1, 'b': 2, 'c': 2, 'd': 1}
codomain_12 = {1, 2}
# Mathematiquement: Im(f) = {1,2}. Codomaine = {1,2}. Im(f) == Codomaine. Surjective.
print(f"f_surjective_not_injective: {f_surjective_not_injective}, Codomaine: {codomain_12} est surjective? {is_surjective(f_surjective_not_injective, codomain_12)}")
assert is_surjective(f_surjective_not_injective, codomain_12) is True, "f_surjective_not_injective devrait etre surjective"

# Fonction vide (non surjective si le codomaine n'est pas vide)
f_empty = {}
codomain_non_empty = {1}
# Mathematiquement: Im(f) = {}. Codomaine = {1}. Im(f) != Codomaine. Non surjective.
print(f"f_empty: {f_empty}, Codomaine: {codomain_non_empty} est surjective? {is_surjective(f_empty, codomain_non_empty)}")
assert is_surjective(f_empty, codomain_non_empty) is False, "f_empty ne devrait PAS etre surjective si codomaine non vide"

# Fonction vide (surjective si le codomaine est vide)
f_empty = {}
codomain_empty = set()
# Mathematiquement: Im(f) = {}. Codomaine = {}. Im(f) == Codomaine. Surjective.
print(f"f_empty: {f_empty}, Codomaine: {codomain_empty} est surjective? {is_surjective(f_empty, codomain_empty)}")
assert is_surjective(f_empty, codomain_empty) is True, "f_empty devrait etre surjective si codomaine vide"

print("Assertions de surjectivite passees avec succes.")
\end{lstlisting}

\subsection{Partie 4 : Vérification de la bijectivité
(Bonus)}\label{partie-4-vuxe9rification-de-la-bijectivituxe9-bonus}

Une fonction est bijective si et seulement si elle est à la fois
injective et surjective. L'implémentation est donc directe en utilisant
les fonctions précédemment définies.

\begin{lstlisting}[language=Python]
# --- Partie 4 : Verification de la bijectivite (Bonus) ---

def is_bijective(f_dict: dict, codomain_set: set) -> bool:
    """
    Verifie si une fonction est bijective.
    Une fonction est bijective si elle est a la fois injective et surjective.

    Args:
        f_dict (dict): Dictionnaire representant la fonction {x: f(x)}.
        codomain_set (set): L'ensemble explicite du codomaine B.

    Returns:
        bool: True si la fonction est bijective, False sinon.
    """
    # Une fonction est bijective si elle est injective ET surjective.
    return is_injective(f_dict) and is_surjective(f_dict, codomain_set)

# --- Tests et Assertions pour la bijectivite ---

print("\n--- Tests de bijectivite ---")

# Fonction bijective (injective et surjective)
f_bijective = {'a': 1, 'b': 2, 'c': 3}
codomain_123 = {1, 2, 3}
# Mathematiquement: Injective (images distinctes), Surjective (Im(f) == Codomaine). Bijective.
print(f"f_bijective: {f_bijective}, Codomaine: {codomain_123} est bijective? {is_bijective(f_bijective, codomain_123)}")
assert is_bijective(f_bijective, codomain_123) is True, "f_bijective devrait etre bijective"

# Fonction non bijective (non injective mais surjective)
f_not_bijective_not_injective = {'a': 1, 'b': 2, 'c': 2}
codomain_12 = {1, 2}
# Mathematiquement: Non injective (f(b)=f(c)=2), Surjective (Im(f) == Codomaine). Non bijective.
print(f"f_not_bijective_not_injective: {f_not_bijective_not_injective}, Codomaine: {codomain_12} est bijective? {is_bijective(f_not_bijective_not_injective, codomain_12)}")
assert is_bijective(f_not_bijective_not_injective, codomain_12) is False, "f_not_bijective_not_injective ne devrait PAS etre bijective"

# Fonction non bijective (injective mais non surjective)
f_not_bijective_not_surjective = {'a': 1, 'b': 2}
codomain_123 = {1, 2, 3}
# Mathematiquement: Injective (images distinctes), Non surjective (3 non atteint). Non bijective.
print(f"f_not_bijective_not_surjective: {f_not_bijective_not_surjective}, Codomaine: {codomain_123} est bijective? {is_bijective(f_not_bijective_not_surjective, codomain_123)}")
assert is_bijective(f_not_bijective_not_surjective, codomain_123) is False, "f_not_bijective_not_surjective ne devrait PAS etre bijective"

# Fonction non bijective (ni injective ni surjective)
f_neither = {'a': 1, 'b': 1, 'c': 2}
codomain_123 = {1, 2, 3}
# Mathematiquement: Non injective (f(a)=f(b)=1), Non surjective (3 non atteint). Non bijective.
print(f"f_neither: {f_neither}, Codomaine: {codomain_123} est bijective? {is_bijective(f_neither, codomain_123)}")
assert is_bijective(f_neither, codomain_123) is False, "f_neither ne devrait PAS etre bijective"

# Fonction vide (bijective si et seulement si le codomaine est vide)
f_empty = {}
codomain_empty = set()
print(f"f_empty: {f_empty}, Codomaine: {codomain_empty} est bijective? {is_bijective(f_empty, codomain_empty)}")
assert is_bijective(f_empty, codomain_empty) is True, "f_empty devrait etre bijective si codomaine vide"

f_empty_codomain_non_empty = {}
codomain_non_empty = {1}
print(f"f_empty: {f_empty_codomain_non_empty}, Codomaine: {codomain_non_empty} est bijective? {is_bijective(f_empty_codomain_non_empty, codomain_non_empty)}")
assert is_bijective(f_empty_codomain_non_empty, codomain_non_empty) is False, "f_empty ne devrait PAS etre bijective si codomaine non vide"

print("Assertions de bijectivite passees avec succes.")
\end{lstlisting}

\subsection{Conclusion}\label{conclusion}

Ce TP vous a permis d'implémenter les vérifications d'injectivité et de
surjectivité pour des fonctions sur des ensembles finis, représentées
par des dictionnaires Python. Vous avez vu comment traduire des
définitions mathématiques rigoureuses en algorithmes concrets et comment
valider ces implémentations à l'aide d'assertions.

La distinction entre l'image d'une fonction (ce que le dictionnaire
contient comme valeurs) et le codomaine (l'ensemble de toutes les
valeurs possibles à l'arrivée) est fondamentale pour la surjectivité.

Pour aller plus loin, vous pourriez explorer : * Comment représenter des
fonctions avec des domaines et codomaines infinis (nécessiterait une
approche différente, non basée sur des énumérations exhaustives). *
L'implémentation de la composition de fonctions, qui sera le sujet du
prochain TP. * La recherche de la fonction inverse pour une fonction
bijective.
\newpage
\section{TP 2}
\section{TP 2/5 : Composition de Fonctions - Implémentation et
Vérification
Empirique}\label{tp-25-composition-de-fonctions---impluxe9mentation-et-vuxe9rification-empirique}

\textbf{Jalon 5 : Applications, injections, surjections, bijections et
composition de fonctions}

\begin{center}\rule{0.5\linewidth}{0.5pt}\end{center}

\subsection{Introduction}\label{introduction}

En tant que développeurs et mathématiciens, la capacité à manipuler et
composer des fonctions est fondamentale. La composition de fonctions est
une opération qui prend deux fonctions et en produit une troisième, qui
représente l'application successive des deux fonctions originales. Ce
concept est au cœur de nombreuses constructions mathématiques et
informatiques, de la programmation fonctionnelle aux architectures de
réseaux neuronaux.

Ce Travail Pratique (TP) vise à solidifier votre compréhension de la
composition de fonctions en vous demandant de l'implémenter en Python de
manière générique, puis de valider empiriquement votre implémentation à
l'aide de tests rigoureux.

\begin{center}\rule{0.5\linewidth}{0.5pt}\end{center}

\subsection{Objectifs du TP}\label{objectifs-du-tp}

\begin{enumerate}
\def\labelenumi{\arabic{enumi}.}
\tightlist
\item
  \textbf{Comprendre la définition mathématique} de la composition de
  fonctions.
\item
  \textbf{Implémenter une fonction générique
  \passthrough{\lstinline!compose(f, g)!}} en Python pur, qui retourne
  une nouvelle fonction représentant \passthrough{\lstinline!g o f!}.
\item
  \textbf{Utiliser des assertions (\passthrough{\lstinline!assert!})}
  pour valider mathématiquement la justesse de l'implémentation avec
  divers exemples.
\item
  \textbf{Développer une approche rigoureuse} de test et de vérification
  de code.
\end{enumerate}

\begin{center}\rule{0.5\linewidth}{0.5pt}\end{center}

\subsection{Partie 1 : Rappels Mathématiques sur la Composition de
Fonctions}\label{partie-1-rappels-mathuxe9matiques-sur-la-composition-de-fonctions}

Soient deux fonctions \(f: A \to B\) et \(g: B \to C\). La
\textbf{composition de \(f\) par \(g\)}, notée \(g \circ f\) (lire ``g
rond f''), est une nouvelle fonction définie par :
\[ (g \circ f)(x) = g(f(x)) \] pour tout \(x\) dans le domaine de \(f\).

\textbf{Points clés :} * L'ordre est crucial : \(g \circ f\) signifie
que \(f\) est appliquée en premier, puis \(g\) est appliquée au résultat
de \(f\). * Le codomaine de la première fonction (\(f\)) doit être
compatible avec le domaine de la seconde fonction (\(g\)). Plus
précisément, l'image de \(f\) doit être un sous-ensemble du domaine de
\(g\). * La fonction résultante \(g \circ f\) a pour domaine \(A\) et
pour codomaine \(C\).

\textbf{Exemple :} Si \(f(x) = x + 1\) et \(g(x) = 2x\). Alors
\((g \circ f)(x) = g(f(x)) = g(x+1) = 2(x+1) = 2x + 2\). Et
\((f \circ g)(x) = f(g(x)) = f(2x) = 2x + 1\). On observe que
\(g \circ f \neq f \circ g\), la composition n'est pas commutative en
général.

\begin{center}\rule{0.5\linewidth}{0.5pt}\end{center}

\subsection{\texorpdfstring{Partie 2 : Implémentation de la Fonction
\texttt{compose(f,\ g)}}{Partie 2 : Implémentation de la Fonction compose(f, g)}}\label{partie-2-impluxe9mentation-de-la-fonction-composef-g}

Votre tâche est d'écrire une fonction Python
\passthrough{\lstinline!compose(f, g)!} qui prend deux fonctions
\passthrough{\lstinline!f!} et \passthrough{\lstinline!g!} en arguments
et retourne une nouvelle fonction. Cette nouvelle fonction doit
représenter la composition \(g \circ f\).

\textbf{Contraintes :} * Python PUR : Aucune bibliothèque externe (comme
NumPy, SciPy, etc.) n'est autorisée. * Le code doit être profondément
commenté pour expliquer chaque étape et la logique mathématique
sous-jacente.

\begin{lstlisting}[language=Python]
def compose(f, g):
    """
    Implemente la composition de deux fonctions f et g.
    Mathematiquement, cela correspond a (g o f)(x) = g(f(x)).

    Args:
        f (callable): La premiere fonction a appliquer. Elle prend un argument
                      et retourne une valeur.
        g (callable): La seconde fonction a appliquer. Elle prend le resultat
                      de f comme argument et retourne une valeur.

    Returns:
        callable: Une nouvelle fonction qui represente la composition (g o f).
                  Cette fonction prend un argument 'x' et retourne g(f(x)).

    Preconditions (implicites en Python, mais importantes mathematiquement) :
        - Le codomaine de f doit etre compatible avec le domaine de g.
          C'est-a-dire que le type et la nature de la valeur retournee par f(x)
          doivent etre acceptables comme argument pour g.
    """
    # La fonction 'compose' doit retourner une nouvelle fonction.
    # Nous utilisons une fonction interne (closure) pour capturer 'f' et 'g'
    # de l'environnement englobant.
    def composed_function(x):
        """
        La fonction resultante de la composition (g o f).
        Elle applique f a x en premier, puis applique g au resultat de f(x).
        """
        # Etape 1 : Appliquer la premiere fonction 'f' a l'argument 'x'.
        # Le resultat de cette operation est la valeur intermediaire.
        result_f = f(x)

        # Etape 2 : Appliquer la seconde fonction 'g' au resultat de 'f(x)'.
        # C'est l'essence meme de la definition (g o f)(x) = g(f(x)).
        result_g_of_f = g(result_f)

        # Retourner le resultat final de la composition.
        return result_g_of_f

    # Retourner la fonction composee, pas son execution.
    return composed_function
\end{lstlisting}

\begin{center}\rule{0.5\linewidth}{0.5pt}\end{center}

\subsection{\texorpdfstring{Partie 3 : Vérification Empirique et Tests
Unitaires (avec
\texttt{assert})}{Partie 3 : Vérification Empirique et Tests Unitaires (avec assert)}}\label{partie-3-vuxe9rification-empirique-et-tests-unitaires-avec-assert}

Pour valider votre implémentation, vous allez créer plusieurs fonctions
simples et les composer. Ensuite, vous utiliserez des assertions pour
vérifier que le résultat de la fonction composée correspond bien au
calcul manuel attendu.

Pour chaque test, vous devrez : 1. Définir les fonctions
\passthrough{\lstinline!f!} et \passthrough{\lstinline!g!}. 2. Calculer
manuellement l'expression de \((g \circ f)(x)\). 3. Appliquer votre
fonction \passthrough{\lstinline!compose!} pour obtenir la fonction
\passthrough{\lstinline!h = compose(f, g)!}. 4. Tester
\passthrough{\lstinline!h!} avec plusieurs valeurs d'entrée
\passthrough{\lstinline!x!} et utiliser \passthrough{\lstinline!assert!}
pour comparer le résultat de \passthrough{\lstinline!h(x)!} avec le
calcul manuel.

\subsubsection{Exemple de Structure de
Test}\label{exemple-de-structure-de-test}

\begin{lstlisting}[language=Python]
# --- Definition des fonctions de base ---
def add_one(x):
    return x + 1

def multiply_by_two(x):
    return x * 2

# --- Test 1 : (multiply_by_two o add_one)(x) ---
# Calcul manuel : (multiply_by_two o add_one)(x) = multiply_by_two(add_one(x))
#                                                = multiply_by_two(x + 1)
#                                                = 2 * (x + 1)
#                                                = 2x + 2

# Creation de la fonction composee
h1 = compose(add_one, multiply_by_two)

# Verifications avec assert
print(f"Test 1: (multiply_by_two o add_one)(x) = 2x + 2")
assert h1(0) == 2 * 0 + 2, f"Erreur pour x=0: Attendu {2*0+2}, Obtenu {h1(0)}"
assert h1(1) == 2 * 1 + 2, f"Erreur pour x=1: Attendu {2*1+2}, Obtenu {h1(1)}"
assert h1(5) == 2 * 5 + 2, f"Erreur pour x=5: Attendu {2*5+2}, Obtenu {h1(5)}"
assert h1(-3) == 2 * (-3) + 2, f"Erreur pour x=-3: Attendu {2*(-3)+2}, Obtenu {h1(-3)}"
print("  -> Test 1 reussi.")

# --- Test 2 : (add_one o multiply_by_two)(x) ---
# Calcul manuel : (add_one o multiply_by_two)(x) = add_one(multiply_by_two(x))
#                                                = add_one(2x)
#                                                = 2x + 1

# Creation de la fonction composee
h2 = compose(multiply_by_two, add_one)

# Verifications avec assert
print(f"\nTest 2: (add_one o multiply_by_two)(x) = 2x + 1")
assert h2(0) == 2 * 0 + 1, f"Erreur pour x=0: Attendu {2*0+1}, Obtenu {h2(0)}"
assert h2(1) == 2 * 1 + 1, f"Erreur pour x=1: Attendu {2*1+1}, Obtenu {h2(1)}"
assert h2(5) == 2 * 5 + 1, f"Erreur pour x=5: Attendu {2*5+1}, Obtenu {h2(5)}"
assert h2(-3) == 2 * (-3) + 1, f"Erreur pour x=-3: Attendu {2*(-3)+1}, Obtenu {h2(-3)}"
print("  -> Test 2 reussi.")
\end{lstlisting}

\subsubsection{Exercices de
Vérification}\label{exercices-de-vuxe9rification}

Implémentez les tests suivants en utilisant la même structure.

\paragraph{Exercice 3.1 : Fonctions
polynomiales}\label{exercice-3.1-fonctions-polynomiales}

Définissez les fonctions suivantes : *
\passthrough{\lstinline!square(x)!} qui retourne \(x^2\). *
\passthrough{\lstinline!subtract\_five(x)!} qui retourne \(x - 5\).

Testez les compositions suivantes : 1.
\passthrough{\lstinline!(subtract\_five o square)(x)!} 2.
\passthrough{\lstinline!(square o subtract\_five)(x)!}

\begin{lstlisting}[language=Python]
# --- Exercice 3.1 : Fonctions polynomiales ---

def square(x):
    """Retourne le carre de x."""
    return x ** 2

def subtract_five(x):
    """Retourne x moins 5."""
    return x - 5

# 1. Composition (subtract_five o square)(x)
# Calcul manuel : (subtract_five o square)(x) = subtract_five(square(x))
#                                              = subtract_five(x**2)
#                                              = x**2 - 5
h3_1 = compose(square, subtract_five)
print(f"\nExercice 3.1.1: (subtract_five o square)(x) = x**2 - 5")
assert h3_1(0) == 0**2 - 5, f"Erreur pour x=0: Attendu {0**2-5}, Obtenu {h3_1(0)}"
assert h3_1(2) == 2**2 - 5, f"Erreur pour x=2: Attendu {2**2-5}, Obtenu {h3_1(2)}"
assert h3_1(-2) == (-2)**2 - 5, f"Erreur pour x=-2: Attendu {(-2)**2-5}, Obtenu {h3_1(-2)}"
assert h3_1(3) == 3**2 - 5, f"Erreur pour x=3: Attendu {3**2-5}, Obtenu {h3_1(3)}"
print("  -> Exercice 3.1.1 reussi.")


# 2. Composition (square o subtract_five)(x)
# Calcul manuel : (square o subtract_five)(x) = square(subtract_five(x))
#                                              = square(x - 5)
#                                              = (x - 5)**2
h3_2 = compose(subtract_five, square)
print(f"\nExercice 3.1.2: (square o subtract_five)(x) = (x - 5)**2")
assert h3_2(0) == (0 - 5)**2, f"Erreur pour x=0: Attendu {(0-5)**2}, Obtenu {h3_2(0)}"
assert h3_2(5) == (5 - 5)**2, f"Erreur pour x=5: Attendu {(5-5)**2}, Obtenu {h3_2(5)}"
assert h3_2(7) == (7 - 5)**2, f"Erreur pour x=7: Attendu {(7-5)**2}, Obtenu {h3_2(7)}"
assert h3_2(-1) == (-1 - 5)**2, f"Erreur pour x=-1: Attendu {(-1-5)**2}, Obtenu {h3_2(-1)}"
print("  -> Exercice 3.1.2 reussi.")
\end{lstlisting}

\paragraph{Exercice 3.2 : Composition de trois fonctions
(associativité)}\label{exercice-3.2-composition-de-trois-fonctions-associativituxe9}

La composition de fonctions est associative, c'est-à-dire que
\((h \circ g) \circ f = h \circ (g \circ f)\). Nous allons vérifier cela
empiriquement.

Définissez les fonctions suivantes : * \passthrough{\lstinline!f1(x)!}
qui retourne \(x + 1\). * \passthrough{\lstinline!f2(x)!} qui retourne
\(x \times 2\). * \passthrough{\lstinline!f3(x)!} qui retourne
\(x - 3\).

Testez la composition \passthrough{\lstinline!(f3 o f2 o f1)(x)!} de
deux manières : 1.
\passthrough{\lstinline!h\_left = compose(f1, compose(f2, f3))!}
(équivalent à \passthrough{\lstinline!(f3 o f2) o f1!}) 2.
\passthrough{\lstinline!h\_right = compose(compose(f1, f2), f3)!}
(équivalent à \passthrough{\lstinline!f3 o (f2 o f1)!})

Vérifiez que \passthrough{\lstinline!h\_left(x)!} et
\passthrough{\lstinline!h\_right(x)!} donnent les mêmes résultats.

\begin{lstlisting}[language=Python]
# --- Exercice 3.2 : Composition de trois fonctions (associativite) ---

def f1(x):
    """Fonction f1(x) = x + 1."""
    return x + 1

def f2(x):
    """Fonction f2(x) = x * 2."""
    return x * 2

def f3(x):
    """Fonction f3(x) = x - 3."""
    return x - 3

# Calcul manuel de (f3 o f2 o f1)(x):
# f1(x) = x + 1
# f2(f1(x)) = f2(x + 1) = 2 * (x + 1) = 2x + 2
# f3(f2(f1(x))) = f3(2x + 2) = (2x + 2) - 3 = 2x - 1
expected_result_expr = lambda x: 2 * x - 1

# 1. Composition de gauche a droite : (f3 o f2) o f1
# C'est-a-dire : compose(f1, compose(f2, f3))
# compose(f2, f3) represente (f3 o f2)
# compose(f1, (f3 o f2)) represente ((f3 o f2) o f1)
h_left = compose(f1, compose(f2, f3))
print(f"\nExercice 3.2.1: (f3 o f2 o f1)(x) via (f3 o f2) o f1 = 2x - 1")
assert h_left(0) == expected_result_expr(0), f"Erreur pour x=0: Attendu {expected_result_expr(0)}, Obtenu {h_left(0)}"
assert h_left(1) == expected_result_expr(1), f"Erreur pour x=1: Attendu {expected_result_expr(1)}, Obtenu {h_left(1)}"
assert h_left(10) == expected_result_expr(10), f"Erreur pour x=10: Attendu {expected_result_expr(10)}, Obtenu {h_left(10)}"
assert h_left(-5) == expected_result_expr(-5), f"Erreur pour x=-5: Attendu {expected_result_expr(-5)}, Obtenu {h_left(-5)}"
print("  -> Exercice 3.2.1 reussi.")


# 2. Composition de droite a gauche : f3 o (f2 o f1)
# C'est-a-dire : compose(compose(f1, f2), f3)
# compose(f1, f2) represente (f2 o f1)
# compose((f2 o f1), f3) represente (f3 o (f2 o f1))
h_right = compose(compose(f1, f2), f3)
print(f"\nExercice 3.2.2: (f3 o f2 o f1)(x) via f3 o (f2 o f1) = 2x - 1")
assert h_right(0) == expected_result_expr(0), f"Erreur pour x=0: Attendu {expected_result_expr(0)}, Obtenu {h_right(0)}"
assert h_right(1) == expected_result_expr(1), f"Erreur pour x=1: Attendu {expected_result_expr(1)}, Obtenu {h_right(1)}"
assert h_right(10) == expected_result_expr(10), f"Erreur pour x=10: Attendu {expected_result_expr(10)}, Obtenu {h_right(10)}"
assert h_right(-5) == expected_result_expr(-5), f"Erreur pour x=-5: Attendu {expected_result_expr(-5)}, Obtenu {h_right(-5)}"
print("  -> Exercice 3.2.2 reussi.")

# Verification de l'associativite : les deux fonctions composees doivent donner le meme resultat
print(f"\nExercice 3.2.3: Verification de l'associativite (h_left == h_right)")
for x_val in [-10, -1, 0, 1, 10]:
    assert h_left(x_val) == h_right(x_val), \
        f"Erreur d'associativite pour x={x_val}: h_left={h_left(x_val)}, h_right={h_right(x_val)}"
print("  -> Exercice 3.2.3 (associativite) reussi : h_left et h_right donnent les memes resultats.")
\end{lstlisting}

\paragraph{Exercice 3.3 : Fonctions avec des types de données différents
(si
applicable)}\label{exercice-3.3-fonctions-avec-des-types-de-donnuxe9es-diffuxe9rents-si-applicable}

Bien que Python soit dynamiquement typé, il est bon de réfléchir aux
implications mathématiques. Considérons une fonction qui prend un nombre
et retourne une chaîne de caractères, puis une autre qui prend une
chaîne et retourne un nombre.

\begin{itemize}
\tightlist
\item
  \passthrough{\lstinline!num\_to\_str(x)!} qui retourne
  \passthrough{\lstinline'str(x) + "!"'}
\item
  \passthrough{\lstinline!str\_len(s)!} qui retourne
  \passthrough{\lstinline!len(s)!}
\end{itemize}

Testez la composition
\passthrough{\lstinline!(str\_len o num\_to\_str)(x)!}.

\begin{lstlisting}[language=Python]
# --- Exercice 3.3 : Fonctions avec des types de donnees differents ---

def num_to_str(x):
    """Convertit un nombre en chaine de caracteres et ajoute un '!'."""
    return str(x) + "!"

def str_len(s):
    """Retourne la longueur d'une chaine de caracteres."""
    return len(s)

# Composition (str_len o num_to_str)(x)
# Calcul manuel : (str_len o num_to_str)(x) = str_len(num_to_str(x))
#                                            = str_len(str(x) + "!")
# Pour x=10, str(10) + "!" = "10!", len("10!") = 3
# Pour x=123, str(123) + "!" = "123!", len("123!") = 4
h4 = compose(num_to_str, str_len)
print(f"\nExercice 3.3: (str_len o num_to_str)(x)")
assert h4(0) == len(str(0) + "!"), f"Erreur pour x=0: Attendu {len(str(0)+'!')}, Obtenu {h4(0)}"
assert h4(10) == len(str(10) + "!"), f"Erreur pour x=10: Attendu {len(str(10)+'!')}, Obtenu {h4(10)}"
assert h4(123) == len(str(123) + "!"), f"Erreur pour x=123: Attendu {len(str(123)+'!')}, Obtenu {h4(123)}"
assert h4(-5) == len(str(-5) + "!"), f"Erreur pour x=-5: Attendu {len(str(-5)+'!')}, Obtenu {h4(-5)}"
print("  -> Exercice 3.3 reussi.")
\end{lstlisting}

\begin{center}\rule{0.5\linewidth}{0.5pt}\end{center}

\subsection{Consignes de Rendu}\label{consignes-de-rendu}

\begin{itemize}
\tightlist
\item
  Le code Python doit être contenu dans un seul fichier Markdown, comme
  demandé.
\item
  Assurez-vous que tous les \passthrough{\lstinline!assert!} passent
  sans erreur.
\item
  Les commentaires doivent être clairs, concis et pertinents, expliquant
  à la fois la logique de programmation et les liens avec les concepts
  mathématiques.
\item
  Le code doit être ``pur Python'', sans aucune importation de
  bibliothèques non standard.
\end{itemize}

\begin{center}\rule{0.5\linewidth}{0.5pt}\end{center}

\subsection{Critères d'Évaluation}\label{crituxe8res-duxe9valuation}

\begin{itemize}
\tightlist
\item
  \textbf{Fonctionnalité (40\%)} : La fonction
  \passthrough{\lstinline!compose!} fonctionne-t-elle correctement pour
  tous les cas de test ?
\item
  \textbf{Conformité aux règles (20\%)} : Respect des contraintes
  ``Python pur'', utilisation des \passthrough{\lstinline!assert!},
  format Markdown.
\item
  \textbf{Rigueur Mathématique (20\%)} : Clarté des explications
  mathématiques pour chaque test et dans les commentaires du code.
\item
  \textbf{Qualité du Code (20\%)} : Lisibilité, commentaires pertinents,
  respect des bonnes pratiques de codage (PEP 8 si applicable, noms de
  variables clairs).
\end{itemize}

\begin{center}\rule{0.5\linewidth}{0.5pt}\end{center}
\newpage
\section{TP 3}
\section{Travail Pratique (TP) 3/5 : Calcul de l'Image Réciproque d'un
Sous-Ensemble}\label{travail-pratique-tp-35-calcul-de-limage-ruxe9ciproque-dun-sous-ensemble}

\textbf{Jalon 5 : Applications, injections, surjections, bijections et
composition de fonctions}

\begin{center}\rule{0.5\linewidth}{0.5pt}\end{center}

\subsection{Introduction}\label{introduction}

Ce troisième Travail Pratique de notre jalon sur les fonctions se
concentre sur un concept fondamental en théorie des ensembles et des
fonctions : l'image réciproque (ou pré-image) d'un sous-ensemble.
Comprendre et savoir implémenter ce concept est crucial pour manipuler
des fonctions de manière rigoureuse en programmation, notamment dans des
domaines comme l'analyse de données, la logique ou l'intelligence
artificielle.

En tant que Développeur Senior, vous serez souvent amené à traduire des
définitions mathématiques abstraites en code fonctionnel et efficace. En
tant que Professeur de Mathématiques, je m'assurerai que les fondements
théoriques soient solides et que l'implémentation reflète fidèlement la
définition mathématique.

\begin{center}\rule{0.5\linewidth}{0.5pt}\end{center}

\subsection{Rappel Mathématique : L'Image
Réciproque}\label{rappel-mathuxe9matique-limage-ruxe9ciproque}

Soit une fonction \(f: A \to B\), où \(A\) est l'ensemble de départ
(domaine) et \(B\) est l'ensemble d'arrivée (codomaine). Soit \(Y\) un
sous-ensemble de \(B\) (c'est-à-dire \(Y \subseteq B\)).

L'\textbf{image réciproque} de \(Y\) par \(f\), notée \(f^{-1}(Y)\), est
l'ensemble de tous les éléments \(x\) du domaine \(A\) tels que leur
image par \(f\), \(f(x)\), appartient à \(Y\).

Formellement, on définit \(f^{-1}(Y)\) comme suit :
\[ f^{-1}(Y) = \{x \in A \mid f(x) \in Y\} \]

Il est important de noter que \(f^{-1}(Y)\) est toujours un
sous-ensemble de \(A\). La notation \(f^{-1}\) ne signifie pas
nécessairement que la fonction \(f\) est inversible ; elle désigne
simplement l'opération de ``remonter'' un sous-ensemble du codomaine
vers le domaine.

\begin{center}\rule{0.5\linewidth}{0.5pt}\end{center}

\subsection{Énoncé du Problème}\label{uxe9noncuxe9-du-probluxe8me}

Votre tâche est d'implémenter une fonction Python pure, nommée
\passthrough{\lstinline!preimage!}, qui calcule l'image réciproque d'un
sous-ensemble donné par une fonction donnée.

\textbf{Contraintes et Spécifications :} 1. La fonction
\passthrough{\lstinline!preimage!} doit prendre trois arguments : *
\passthrough{\lstinline!f!}: La fonction mathématique (représentée par
un callable Python). * \passthrough{\lstinline!domain\_A!}: L'ensemble
de départ \(A\) (représenté par un \passthrough{\lstinline!set!}
Python). * \passthrough{\lstinline!subset\_Y!}: Le sous-ensemble \(Y\)
du codomaine (représenté par un \passthrough{\lstinline!set!} Python).
2. La fonction doit retourner un \passthrough{\lstinline!set!} Python
représentant \(f^{-1}(Y)\). 3. Le code doit être entièrement ``from
scratch'', sans utiliser de bibliothèques de calcul numérique avancées
(comme NumPy, SciPy, etc.). Les types de données Python natifs
(\passthrough{\lstinline!set!}, \passthrough{\lstinline!list!},
\passthrough{\lstinline!dict!}, fonctions, etc.) sont autorisés. 4. Le
code doit être profondément commenté pour expliquer chaque étape. 5. Des
assertions (\passthrough{\lstinline!assert!}) doivent être utilisées
pour valider la justesse mathématique de l'implémentation avec des
exemples concrets. 6. Des explications mathématiques rigoureuses doivent
accompagner le code pour justifier les choix d'implémentation.

\begin{center}\rule{0.5\linewidth}{0.5pt}\end{center}

\subsection{Implémentation en Python}\label{impluxe9mentation-en-python}

Nous allons représenter les ensembles \(A\) et \(Y\) par des objets
\passthrough{\lstinline!set!} de Python. La fonction \(f\) sera une
fonction Python standard (ou une lambda).

\begin{lstlisting}[language=Python]
def preimage(f, domain_A: set, subset_Y: set) -> set:
    """
    Calcule l'image reciproque (pre-image) d'un sous-ensemble Y par une fonction f.

    Soit f: A -> B une fonction, et Y un sous-ensemble de B.
    L'image reciproque de Y par f est l'ensemble {x appartenant a A | f(x) appartient a Y}.

    Args:
        f (callable): La fonction mathematique f. Elle doit prendre un element de domain_A
                      et retourner un element du codomaine B.
        domain_A (set): L'ensemble de depart A de la fonction f. C'est un ensemble
                        de tous les elements possibles pour x.
        subset_Y (set): Le sous-ensemble Y du codomaine B dont on veut calculer
                        l'image reciproque.

    Returns:
        set: L'ensemble des elements x de domain_A tels que f(x) est dans subset_Y.
             C'est f^-1(Y).
    """
    # Initialisation d'un ensemble vide pour stocker les elements de l'image reciproque.
    # Cet ensemble sera notre f^-1(Y).
    preimage_set = set()

    # Parcourir chaque element 'x' de l'ensemble de depart 'domain_A'.
    # C'est la premiere partie de la definition mathematique : {x appartenant a A ...}
    for x in domain_A:
        # Calculer l'image de 'x' par la fonction 'f'.
        # C'est la partie f(x) de la definition.
        image_of_x = f(x)

        # Verifier si l'image de 'x' appartient au sous-ensemble 'subset_Y'.
        # C'est la deuxieme partie de la definition : ... | f(x) appartient a Y}
        if image_of_x in subset_Y:
            # Si f(x) est dans Y, alors 'x' fait partie de l'image reciproque.
            # Ajouter 'x' a notre ensemble resultat.
            preimage_set.add(x)

    # Retourner l'ensemble construit, qui est l'image reciproque f^-1(Y).
    return preimage_set
\end{lstlisting}

\subsubsection{Exemples et Assertions pour la Validation
Mathématique}\label{exemples-et-assertions-pour-la-validation-mathuxe9matique}

Nous allons maintenant tester notre fonction
\passthrough{\lstinline!preimage!} avec plusieurs scénarios pour valider
sa justesse mathématique à l'aide d'assertions.

\paragraph{Exemple 1 : Fonction simple et sous-ensemble
clair}\label{exemple-1-fonction-simple-et-sous-ensemble-clair}

Soit \(f: \mathbb{Z} \to \mathbb{Z}\) définie par \(f(x) = x^2\).
Considérons un domaine \(A = \{-3, -2, -1, 0, 1, 2, 3\}\). Et un
sous-ensemble \(Y = \{0, 1, 4\}\).

\(f^{-1}(Y) = \{x \in A \mid x^2 \in \{0, 1, 4\}\}\) Pour
\(x=-3, f(-3)=9 \notin Y\) Pour \(x=-2, f(-2)=4 \in Y\) Pour
\(x=-1, f(-1)=1 \in Y\) Pour \(x=0, f(0)=0 \in Y\) Pour
\(x=1, f(1)=1 \in Y\) Pour \(x=2, f(2)=4 \in Y\) Pour
\(x=3, f(3)=9 \notin Y\) Donc, \(f^{-1}(Y) = \{-2, -1, 0, 1, 2\}\).

\begin{lstlisting}[language=Python]
# --- Exemple 1 ---
print("--- Exemple 1 : Fonction carree ---")
# Definition de la fonction f(x) = x^2
def square_function(x):
    return x * x

# Definition du domaine A
domain_A_ex1 = {-3, -2, -1, 0, 1, 2, 3}

# Definition du sous-ensemble Y
subset_Y_ex1 = {0, 1, 4}

# Calcul de l'image reciproque
result_ex1 = preimage(square_function, domain_A_ex1, subset_Y_ex1)

# Resultat attendu : {-2, -1, 0, 1, 2}
expected_ex1 = {-2, -1, 0, 1, 2}

print(f"Fonction f(x) = x^2")
print(f"Domaine A : {domain_A_ex1}")
print(f"Sous-ensemble Y : {subset_Y_ex1}")
print(f"Image reciproque calculee : {result_ex1}")
print(f"Image reciproque attendue : {expected_ex1}")

# Assertion pour valider le resultat
assert result_ex1 == expected_ex1, f"Erreur dans l'Exemple 1: Attendu {expected_ex1}, obtenu {result_ex1}"
print("Assertion Exemple 1 passee avec succes.\n")
\end{lstlisting}

\paragraph{Exemple 2 : Fonction linéaire et sous-ensemble
vide}\label{exemple-2-fonction-linuxe9aire-et-sous-ensemble-vide}

Soit \(f: \mathbb{R} \to \mathbb{R}\) définie par \(f(x) = 2x + 1\).
Considérons un domaine \(A = \{1, 2, 3, 4, 5\}\). Et un sous-ensemble
\(Y = \{10, 12\}\).

\(f^{-1}(Y) = \{x \in A \mid 2x+1 \in \{10, 12\}\}\) Pour
\(x=1, f(1)=3 \notin Y\) Pour \(x=2, f(2)=5 \notin Y\) Pour
\(x=3, f(3)=7 \notin Y\) Pour \(x=4, f(4)=9 \notin Y\) Pour
\(x=5, f(5)=11 \notin Y\) Aucun élément de \(A\) n'a son image dans
\(Y\). Donc, \(f^{-1}(Y) = \emptyset\).

\begin{lstlisting}[language=Python]
# --- Exemple 2 ---
print("--- Exemple 2 : Fonction lineaire et sous-ensemble vide ---")
# Definition de la fonction f(x) = 2x + 1
def linear_function(x):
    return 2 * x + 1

# Definition du domaine A
domain_A_ex2 = {1, 2, 3, 4, 5}

# Definition du sous-ensemble Y (aucun element de A ne devrait y etre mappe)
subset_Y_ex2 = {10, 12}

# Calcul de l'image reciproque
result_ex2 = preimage(linear_function, domain_A_ex2, subset_Y_ex2)

# Resultat attendu : ensemble vide
expected_ex2 = set()

print(f"Fonction f(x) = 2x + 1")
print(f"Domaine A : {domain_A_ex2}")
print(f"Sous-ensemble Y : {subset_Y_ex2}")
print(f"Image reciproque calculee : {result_ex2}")
print(f"Image reciproque attendue : {expected_ex2}")

# Assertion pour valider le resultat
assert result_ex2 == expected_ex2, f"Erreur dans l'Exemple 2: Attendu {expected_ex2}, obtenu {result_ex2}"
print("Assertion Exemple 2 passee avec succes.\n")
\end{lstlisting}

\paragraph{Exemple 3 : Fonction
constante}\label{exemple-3-fonction-constante}

Soit \(f: \{a, b, c, d\} \to \{1, 2, 3\}\) définie par \(f(x) = 2\) pour
tout \(x\). Considérons un domaine \(A = \{a, b, c, d\}\). Et un
sous-ensemble \(Y = \{1, 2\}\).

\(f^{-1}(Y) = \{x \in A \mid f(x) \in \{1, 2\}\}\) Pour
\(x=a, f(a)=2 \in Y\) Pour \(x=b, f(b)=2 \in Y\) Pour
\(x=c, f(c)=2 \in Y\) Pour \(x=d, f(d)=2 \in Y\) Donc,
\(f^{-1}(Y) = \{a, b, c, d\}\).

\begin{lstlisting}[language=Python]
# --- Exemple 3 ---
print("--- Exemple 3 : Fonction constante ---")
# Definition de la fonction f(x) = 2
def constant_function(x):
    return 2

# Definition du domaine A
domain_A_ex3 = {'a', 'b', 'c', 'd'}

# Definition du sous-ensemble Y
subset_Y_ex3 = {1, 2}

# Calcul de l'image reciproque
result_ex3 = preimage(constant_function, domain_A_ex3, subset_Y_ex3)

# Resultat attendu : {'a', 'b', 'c', 'd'}
expected_ex3 = {'a', 'b', 'c', 'd'}

print(f"Fonction f(x) = 2")
print(f"Domaine A : {domain_A_ex3}")
print(f"Sous-ensemble Y : {subset_Y_ex3}")
print(f"Image reciproque calculee : {result_ex3}")
print(f"Image reciproque attendue : {expected_ex3}")

# Assertion pour valider le resultat
assert result_ex3 == expected_ex3, f"Erreur dans l'Exemple 3: Attendu {expected_ex3}, obtenu {result_ex3}"
print("Assertion Exemple 3 passee avec succes.\n")
\end{lstlisting}

\paragraph{Exemple 4 : Sous-ensemble Y ne contenant aucune
image}\label{exemple-4-sous-ensemble-y-ne-contenant-aucune-image}

Soit \(f: \{1, 2, 3\} \to \{'A', 'B', 'C'\}\) définie par
\(f(1)='A', f(2)='B', f(3)='C'\). Considérons un domaine
\(A = \{1, 2, 3\}\). Et un sous-ensemble \(Y = \{'X', 'Y'\}\).

\(f^{-1}(Y) = \{x \in A \mid f(x) \in \{'X', 'Y'\}\}\) Aucun élément de
\(A\) n'a son image dans \(Y\). Donc, \(f^{-1}(Y) = \emptyset\).

\begin{lstlisting}[language=Python]
# --- Exemple 4 ---
print("--- Exemple 4 : Sous-ensemble Y sans intersection avec l'image de f ---")
# Definition de la fonction f
def map_letters(x):
    if x == 1: return 'A'
    if x == 2: return 'B'
    if x == 3: return 'C'
    return None # Pour les cas non definis dans le domaine

# Definition du domaine A
domain_A_ex4 = {1, 2, 3}

# Definition du sous-ensemble Y
subset_Y_ex4 = {'X', 'Y'}

# Calcul de l'image reciproque
result_ex4 = preimage(map_letters, domain_A_ex4, subset_Y_ex4)

# Resultat attendu : ensemble vide
expected_ex4 = set()

print(f"Fonction f(1)='A', f(2)='B', f(3)='C'")
print(f"Domaine A : {domain_A_ex4}")
print(f"Sous-ensemble Y : {subset_Y_ex4}")
print(f"Image reciproque calculee : {result_ex4}")
print(f"Image reciproque attendue : {expected_ex4}")

# Assertion pour valider le resultat
assert result_ex4 == expected_ex4, f"Erreur dans l'Exemple 4: Attendu {expected_ex4}, obtenu {result_ex4}"
print("Assertion Exemple 4 passee avec succes.\n")
\end{lstlisting}

\begin{center}\rule{0.5\linewidth}{0.5pt}\end{center}

\subsection{Justification Mathématique de
l'Implémentation}\label{justification-mathuxe9matique-de-limpluxe9mentation}

L'implémentation Python de la fonction
\passthrough{\lstinline!preimage!} adhère strictement à la définition
mathématique de l'image réciproque :
\(f^{-1}(Y) = \{x \in A \mid f(x) \in Y\}\).

\begin{enumerate}
\def\labelenumi{\arabic{enumi}.}
\item
  \textbf{\passthrough{\lstinline!preimage\_set = set()!}}: Nous
  commençons par un ensemble vide. Ceci représente l'initialisation de
  notre ensemble \(f^{-1}(Y)\) avant d'y ajouter les éléments qui
  satisfont la condition. C'est l'équivalent de commencer la
  construction de l'ensemble par l'ensemble vide.
\item
  \textbf{\passthrough{\lstinline!for x in domain\_A:!}}: Cette boucle
  itère sur chaque élément \passthrough{\lstinline!x!} de
  \passthrough{\lstinline!domain\_A!}. Ceci correspond à la clause
  ``\(x \in A\)'' de la définition. Pour trouver tous les éléments du
  domaine qui satisfont la condition, nous devons nécessairement les
  examiner un par un.
\item
  \textbf{\passthrough{\lstinline!image\_of\_x = f(x)!}}: Pour chaque
  \passthrough{\lstinline!x!} du domaine, nous calculons son image par
  la fonction \passthrough{\lstinline!f!}. Ceci est la partie
  ``\(f(x)\)'' de la condition ``\(f(x) \in Y\)''.
\item
  \textbf{\passthrough{\lstinline!if image\_of\_x in subset\_Y:!}}:
  Cette condition vérifie si l'image calculée
  \passthrough{\lstinline!f(x)!} appartient au sous-ensemble
  \passthrough{\lstinline!subset\_Y!}. Ceci est la partie
  ``\(f(x) \in Y\)'' de la définition. C'est le critère de sélection
  pour qu'un élément \passthrough{\lstinline!x!} fasse partie de l'image
  réciproque.
\item
  \textbf{\passthrough{\lstinline!preimage\_set.add(x)!}}: Si la
  condition est vraie (c'est-à-dire si \(f(x) \in Y\)), alors l'élément
  \passthrough{\lstinline!x!} est ajouté à notre ensemble
  \passthrough{\lstinline!preimage\_set!}. C'est ainsi que nous
  construisons l'ensemble des éléments qui satisfont la propriété.
  L'utilisation d'un \passthrough{\lstinline!set!} Python garantit que
  chaque élément n'est ajouté qu'une seule fois, ce qui est conforme à
  la nature des ensembles mathématiques (pas de doublons).
\end{enumerate}

En résumé, le code parcourt systématiquement le domaine de la fonction,
applique la fonction à chaque élément, puis vérifie si le résultat tombe
dans le sous-ensemble cible. Si c'est le cas, l'élément original du
domaine est inclus dans l'ensemble résultant. Cette approche est une
traduction directe et fidèle de la définition mathématique.

\begin{center}\rule{0.5\linewidth}{0.5pt}\end{center}

\subsection{Conclusion}\label{conclusion}

Ce TP a permis d'implémenter un concept fondamental de la théorie des
fonctions : l'image réciproque d'un sous-ensemble. L'exercice a mis en
lumière l'importance de traduire des définitions mathématiques précises
en algorithmes clairs et corrects. L'utilisation de types de données
Python natifs comme les \passthrough{\lstinline!set!} et les
\passthrough{\lstinline!callable!} pour représenter les ensembles et les
fonctions, combinée à une logique itérative simple, démontre qu'il est
possible de construire des outils mathématiques robustes sans dépendre
de bibliothèques complexes. Les assertions ont joué un rôle clé pour
garantir la justesse mathématique de notre implémentation.

Ce concept est une brique essentielle pour aborder des notions plus
avancées en mathématiques discrètes et en informatique théorique.
\newpage
\section{TP 4}
\section{TP 4/5 - Jalon 5: Dictionnaire Inversible (Bijection) avec
Gestion Stricte des
Erreurs}\label{tp-45---jalon-5-dictionnaire-inversible-bijection-avec-gestion-stricte-des-erreurs}

\subsection{Introduction}\label{introduction}

Bienvenue à ce quatrième Travail Pratique du Jalon 5, axé sur les
applications, injections, surjections et bijections. Dans ce TP, nous
allons mettre en pratique ces concepts fondamentaux de la théorie des
fonctions en développant une structure de données en Python : un
``dictionnaire inversible''.

Un dictionnaire Python standard est une \emph{application} (ou fonction)
qui associe chaque clé unique à une valeur. Cependant, il ne garantit
pas que chaque valeur soit unique (plusieurs clés peuvent pointer vers
la même valeur), ni qu'il soit inversible directement, car si plusieurs
clés pointent vers la même valeur, la fonction réciproque associerait
plusieurs antécédents à une même image, violant la définition même d'une
application. Pour qu'une fonction soit inversible, elle doit être une
\textbf{bijection}.

L'objectif de ce TP est de créer une classe
\passthrough{\lstinline!InvertibleDict!} qui modélise une bijection
parfaite entre un ensemble de clés et un ensemble de valeurs. Cela
signifie que chaque clé doit être unique, chaque valeur doit être
unique, et il doit exister une correspondance biunivoque entre elles.
Nous mettrons un accent particulier sur la gestion stricte des erreurs
pour garantir l'intégrité de cette bijection.

\subsection{Prérequis}\label{pruxe9requis}

\begin{itemize}
\tightlist
\item
  Connaissances solides en programmation orientée objet en Python
  (classes, méthodes, attributs, exceptions).
\item
  Compréhension approfondie des concepts mathématiques d'application,
  injection, surjection et bijection.
\item
  Maîtrise de l'utilisation des \passthrough{\lstinline!assert!} pour la
  validation de code.
\end{itemize}

\subsection{Objectifs du TP}\label{objectifs-du-tp}

\begin{enumerate}
\def\labelenumi{\arabic{enumi}.}
\tightlist
\item
  \textbf{Implémenter une classe
  \passthrough{\lstinline!InvertibleDict!}} : Cette classe doit se
  comporter comme un dictionnaire standard pour l'accès direct (clé
  -\textgreater{} valeur) mais doit également permettre un accès inverse
  (valeur -\textgreater{} clé).
\item
  \textbf{Garantir la Bijection} : Assurer que chaque clé est unique et
  que chaque valeur est unique au sein du dictionnaire.
\item
  \textbf{Gestion Stricte des Erreurs} : Lever des exceptions
  appropriées (\passthrough{\lstinline!KeyError!},
  \passthrough{\lstinline!ValueError!}) en cas de tentative d'ajout de
  clés ou de valeurs dupliquées, ou d'accès à des éléments inexistants.
\item
  \textbf{Code Commenté et Validé} : Le code doit être profondément
  commenté et inclure des \passthrough{\lstinline!assert!} pour valider
  la justesse mathématique des opérations.
\end{enumerate}

\subsection{Rappel Mathématique: La
Bijection}\label{rappel-mathuxe9matique-la-bijection}

Une fonction \(f: A \to B\) est une \textbf{bijection} si et seulement
si elle est à la fois :

\begin{enumerate}
\def\labelenumi{\arabic{enumi}.}
\tightlist
\item
  \textbf{Injective (ou injection)} : Chaque élément de l'ensemble
  d'arrivée \(B\) est atteint par au plus un élément de l'ensemble de
  départ \(A\). Autrement dit, si \(f(x_1) = f(x_2)\), alors
  \(x_1 = x_2\). Il n'y a pas de valeurs dupliquées pour des clés
  différentes.
\item
  \textbf{Surjective (ou surjection)} : Chaque élément de l'ensemble
  d'arrivée \(B\) est atteint par au moins un élément de l'ensemble de
  départ \(A\). Autrement dit, pour tout \(y \in B\), il existe au moins
  un \(x \in A\) tel que \(f(x) = y\). Toutes les valeurs ``possibles''
  sont utilisées.
\end{enumerate}

Lorsqu'une fonction est une bijection, elle possède une \textbf{fonction
inverse} unique \(f^{-1}: B \to A\) telle que \(f^{-1}(f(x)) = x\) pour
tout \(x \in A\) et \(f(f^{-1}(y)) = y\) pour tout \(y \in B\).

Dans le contexte de notre \passthrough{\lstinline!InvertibleDict!}: *
L'ensemble des clés \(K\) est l'ensemble de départ \(A\). * L'ensemble
des valeurs \(V\) est l'ensemble d'arrivée \(B\). * La relation clé
\(\to\) valeur est la fonction \(f\). * La relation valeur \(\to\) clé
est la fonction inverse \(f^{-1}\).

Pour garantir la bijection, lors de l'ajout d'une paire (clé, valeur) :
* La clé ne doit pas déjà exister (garantie par les dictionnaires
standards pour l'unicité des clés). * La valeur ne doit pas déjà exister
(c'est ce que nous devons implémenter pour l'injection).

\subsection{\texorpdfstring{Exercice 1: Conception et Implémentation de
la classe
\texttt{InvertibleDict}}{Exercice 1: Conception et Implémentation de la classe InvertibleDict}}\label{exercice-1-conception-et-impluxe9mentation-de-la-classe-invertibledict}

Vous devez implémenter la classe
\passthrough{\lstinline!InvertibleDict!} en respectant les
spécifications suivantes.

\subsubsection{Structure de la classe}\label{structure-de-la-classe}

La classe \passthrough{\lstinline!InvertibleDict!} doit maintenir deux
dictionnaires internes pour assurer une recherche rapide dans les deux
sens : * \passthrough{\lstinline!\_key\_to\_value!}: Le dictionnaire
principal (clé -\textgreater{} valeur). *
\passthrough{\lstinline!\_value\_to\_key!}: Le dictionnaire inverse
(valeur -\textgreater{} clé).

\subsubsection{Méthodes à
implémenter}\label{muxe9thodes-uxe0-impluxe9menter}

\begin{enumerate}
\def\labelenumi{\arabic{enumi}.}
\tightlist
\item
  \textbf{\passthrough{\lstinline!\_\_init\_\_(self, initial\_data=None)!}}:

  \begin{itemize}
  \tightlist
  \item
    Initialise les deux dictionnaires internes.
  \item
    Si \passthrough{\lstinline!initial\_data!} est fourni (un
    dictionnaire ou une liste de tuples), il doit tenter d'ajouter
    toutes les paires en utilisant la logique d'ajout stricte (lever des
    erreurs si des doublons sont trouvés dans les données initiales).
  \end{itemize}
\item
  \textbf{\passthrough{\lstinline!\_\_setitem\_\_(self, key, value)!}}:

  \begin{itemize}
  \tightlist
  \item
    Permet d'ajouter ou de modifier une paire
    \passthrough{\lstinline!(key, value)!} en utilisant la syntaxe
    \passthrough{\lstinline!idict[key] = value!}.
  \item
    \textbf{Gestion des erreurs cruciale} :

    \begin{itemize}
    \tightlist
    \item
      Si \passthrough{\lstinline!key!} existe déjà et est associée à une
      \passthrough{\lstinline!value!} différente, ou si
      \passthrough{\lstinline!value!} existe déjà et est associée à une
      \passthrough{\lstinline!key!} différente, cela signifie une
      tentative de briser la bijection. Une
      \passthrough{\lstinline!KeyError!} ou
      \passthrough{\lstinline!ValueError!} appropriée doit être levée.
    \item
      Si \passthrough{\lstinline!key!} existe déjà et est associée à la
      \emph{même} \passthrough{\lstinline!value!}, l'opération est
      idempotente et ne fait rien.
    \item
      Si \passthrough{\lstinline!key!} n'existe pas mais
      \passthrough{\lstinline!value!} existe déjà, lever une
      \passthrough{\lstinline!ValueError!} (la valeur n'est pas unique).
    \item
      Si \passthrough{\lstinline!key!} existe mais
      \passthrough{\lstinline!value!} n'existe pas, et que la
      \passthrough{\lstinline!key!} est réassignée à une nouvelle
      \passthrough{\lstinline!value!}, l'ancienne paire doit être
      supprimée et la nouvelle ajoutée.
    \end{itemize}
  \item
    \textbf{Logique d'ajout/modification} :

    \begin{itemize}
    \tightlist
    \item
      Si la clé existe déjà, et la valeur est la même, ne rien faire.
    \item
      Si la clé existe déjà, et la valeur est différente :

      \begin{itemize}
      \tightlist
      \item
        Vérifier si la \emph{nouvelle} valeur existe déjà. Si oui, lever
        \passthrough{\lstinline!ValueError!}.
      \item
        Supprimer l'ancienne paire (clé, ancienne\_valeur) des deux
        dictionnaires.
      \item
        Ajouter la nouvelle paire (clé, nouvelle\_valeur) aux deux
        dictionnaires.
      \end{itemize}
    \item
      Si la clé n'existe pas :

      \begin{itemize}
      \tightlist
      \item
        Vérifier si la valeur existe déjà. Si oui, lever
        \passthrough{\lstinline!ValueError!}.
      \item
        Ajouter la nouvelle paire (clé, valeur) aux deux dictionnaires.
      \end{itemize}
    \end{itemize}
  \end{itemize}
\item
  \textbf{\passthrough{\lstinline!\_\_getitem\_\_(self, key)!}}:

  \begin{itemize}
  \tightlist
  \item
    Permet de récupérer la valeur associée à une clé :
    \passthrough{\lstinline!value = idict[key]!}.
  \item
    Lève une \passthrough{\lstinline!KeyError!} si la clé n'existe pas.
  \end{itemize}
\item
  \textbf{\passthrough{\lstinline!get\_inverse(self, value)!}}:

  \begin{itemize}
  \tightlist
  \item
    Permet de récupérer la clé associée à une valeur.
  \item
    Lève une \passthrough{\lstinline!ValueError!} si la valeur n'existe
    pas. (Utilisation de \passthrough{\lstinline!ValueError!} pour les
    recherches par valeur est une convention raisonnable ici).
  \end{itemize}
\item
  \textbf{\passthrough{\lstinline!\_\_delitem\_\_(self, key)!}}:

  \begin{itemize}
  \tightlist
  \item
    Permet de supprimer une paire par sa clé :
    \passthrough{\lstinline!del idict[key]!}.
  \item
    Lève une \passthrough{\lstinline!KeyError!} si la clé n'existe pas.
  \item
    Doit supprimer l'entrée des deux dictionnaires internes.
  \end{itemize}
\item
  \textbf{\passthrough{\lstinline!\_\_len\_\_(self)!}}:

  \begin{itemize}
  \tightlist
  \item
    Retourne le nombre de paires dans le dictionnaire.
  \end{itemize}
\item
  \textbf{\passthrough{\lstinline!\_\_contains\_\_(self, key)!}}:

  \begin{itemize}
  \tightlist
  \item
    Vérifie si une clé existe dans le dictionnaire.
  \end{itemize}
\item
  \textbf{\passthrough{\lstinline!keys(self)!}}:

  \begin{itemize}
  \tightlist
  \item
    Retourne un itérateur sur les clés.
  \end{itemize}
\item
  \textbf{\passthrough{\lstinline!values(self)!}}:

  \begin{itemize}
  \tightlist
  \item
    Retourne un itérateur sur les valeurs.
  \end{itemize}
\item
  \textbf{\passthrough{\lstinline!items(self)!}}:

  \begin{itemize}
  \tightlist
  \item
    Retourne un itérateur sur les paires (clé, valeur).
  \end{itemize}
\item
  \textbf{\passthrough{\lstinline!\_\_repr\_\_(self)!} et
  \passthrough{\lstinline!\_\_str\_\_(self)!}}:

  \begin{itemize}
  \tightlist
  \item
    Pour une représentation lisible de l'objet.
  \end{itemize}
\end{enumerate}

\subsubsection{Implémentation du code}\label{impluxe9mentation-du-code}

\begin{lstlisting}[language=Python]
class InvertibleDict:
    """
    Implemente un dictionnaire inversible qui garantit une bijection
    entre les cles et les valeurs.

    Chaque cle est unique et chaque valeur est unique.
    Les tentatives d'ajout de cles ou de valeurs dupliquees levent des erreurs.
    """

    def __init__(self, initial_data=None):
        """
        Initialise le dictionnaire inversible.

        Args:
            initial_data (dict ou list de tuple, optionnel): Donnees initiales
                pour peupler le dictionnaire. Chaque paire (cle, valeur) sera
                ajoutee avec la logique de verification de bijection.
        """
        # Dictionnaire principal: cle -> valeur
        self._key_to_value = {}
        # Dictionnaire inverse: valeur -> cle
        self._value_to_key = {}

        if initial_data:
            # Si des donnees initiales sont fournies, les ajouter une par une.
            # Cela permet de valider la bijection des l'initialisation.
            if isinstance(initial_data, dict):
                items = initial_data.items()
            elif isinstance(initial_data, (list, tuple)):
                items = initial_data
            else:
                raise TypeError("initial_data doit etre un dict ou une liste/tuple de paires.")

            for key, value in items:
                # Utilise la methode __setitem__ pour garantir la logique de bijection
                # et la gestion des erreurs lors de l'initialisation.
                self[key] = value

    def __setitem__(self, key, value):
        """
        Ajoute ou modifie une paire (cle, valeur) dans le dictionnaire.
        Garantit que la bijection est maintenue.

        Args:
            key: La cle a ajouter ou modifier.
            value: La valeur associee a la cle.

        Raises:
            KeyError: Si la cle existe deja et est associee a une valeur differente,
                      ou si la cle est tentee d'etre reassignee a une valeur deja existante
                      par une autre cle.
            ValueError: Si la valeur existe deja et est associee a une cle differente.
        """
        # Cas 1: La cle existe deja dans le dictionnaire
        if key in self._key_to_value:
            current_value = self._key_to_value[key]

            # Si la cle est deja associee a la meme valeur, l'operation est idempotente.
            if current_value == value:
                return

            # Si la cle existe mais est associee a une valeur differente,
            # nous devons verifier la nouvelle valeur et potentiellement mettre a jour.

            # Verifier si la nouvelle valeur existe deja, mais est associee a une autre cle.
            if value in self._value_to_key and self._value_to_key[value] != key:
                raise ValueError(f"La valeur '{value}' est deja associee a la cle '{self._value_to_key[value]}'. "
                                 f"Impossible de l'associer a la cle '{key}' car cela briserait l'injection.")

            # Supprimer l'ancienne paire des deux dictionnaires
            del self._key_to_value[key]
            del self._value_to_key[current_value]

            # Ajouter la nouvelle paire
            self._key_to_value[key] = value
            self._value_to_key[value] = key

        # Cas 2: La cle n'existe pas encore
        else:
            # Verifier si la valeur existe deja (pour garantir l'injection)
            if value in self._value_to_key:
                raise ValueError(f"La valeur '{value}' est deja associee a la cle '{self._value_to_key[value]}'. "
                                 f"Impossible de l'associer a la nouvelle cle '{key}' car cela briserait l'injection.")

            # Si tout est bon, ajouter la nouvelle paire aux deux dictionnaires
            self._key_to_value[key] = value
            self._value_to_key[value] = key

    def __getitem__(self, key):
        """
        Recupere la valeur associee a une cle.

        Args:
            key: La cle dont on veut recuperer la valeur.

        Returns:
            La valeur associee a la cle.

        Raises:
            KeyError: Si la cle n'est pas trouvee dans le dictionnaire.
        """
        if key not in self._key_to_value:
            raise KeyError(f"La cle '{key}' n'existe pas dans le dictionnaire.")
        return self._key_to_value[key]

    def get_inverse(self, value):
        """
        Recupere la cle associee a une valeur (fonction inverse).

        Args:
            value: La valeur dont on veut recuperer la cle.

        Returns:
            La cle associee a la valeur.

        Raises:
            ValueError: Si la valeur n'est pas trouvee dans le dictionnaire.
        """
        if value not in self._value_to_key:
            # Utilisation de ValueError pour les recherches par valeur est une convention raisonnable.
            raise ValueError(f"La valeur '{value}' n'existe pas dans le dictionnaire.")
        return self._value_to_key[value]

    def __delitem__(self, key):
        """
        Supprime une paire (cle, valeur) par sa cle.

        Args:
            key: La cle de la paire a supprimer.

        Raises:
            KeyError: Si la cle n'est pas trouvee dans le dictionnaire.
        """
        if key not in self._key_to_value:
            raise KeyError(f"La cle '{key}' n'existe pas dans le dictionnaire.")

        value_to_delete = self._key_to_value[key]

        # Supprimer des deux dictionnaires pour maintenir la coherence
        del self._key_to_value[key]
        del self._value_to_key[value_to_delete]

    def __len__(self):
        """
        Retourne le nombre de paires (cle, valeur) dans le dictionnaire.
        """
        # Les deux dictionnaires doivent toujours avoir la meme taille
        assert len(self._key_to_value) == len(self._value_to_key), \
            "Incoherence interne: les dictionnaires direct et inverse n'ont pas la meme taille."
        return len(self._key_to_value)

    def __contains__(self, key):
        """
        Verifie si une cle existe dans le dictionnaire.
        """
        return key in self._key_to_value

    def keys(self):
        """
        Retourne un iterateur sur les cles du dictionnaire.
        """
        return self._key_to_value.keys()

    def values(self):
        """
        Retourne un iterateur sur les valeurs du dictionnaire.
        """
        return self._key_to_value.values()

    def items(self):
        """
        Retourne un iterateur sur les paires (cle, valeur) du dictionnaire.
        """
        return self._key_to_value.items()

    def __repr__(self):
        """
        Representation officielle de l'objet.
        """
        return f"InvertibleDict({self._key_to_value})"

    def __str__(self):
        """
        Representation lisible de l'objet.
        """
        return str(self._key_to_value)
\end{lstlisting}

\subsubsection{Tests et Assertions
Mathématiques}\label{tests-et-assertions-mathuxe9matiques}

Voici un bloc de code pour tester la classe
\passthrough{\lstinline!InvertibleDict!} et valider son comportement
avec des \passthrough{\lstinline!assert!}.

```python \# --- Bloc de tests pour InvertibleDict ---

print(``--- Démarrage des tests pour InvertibleDict ---'')

\section{Test 1: Initialisation et ajout
basique}\label{test-1-initialisation-et-ajout-basique}

print(``\nTest 1: Initialisation et ajout basique'') idict =
InvertibleDict() idict{[}`a'{]} = 1 idict{[}`b'{]} = 2 idict{[}`c'{]} =
3

assert len(idict) == 3, f''Erreur: Taille attendue 3, obtenue
\{len(idict)\}'' assert idict{[}`a'{]} == 1, ``Erreur: idict{[}`a'{]}
devrait être 1'' assert idict.get\_inverse(1) == `a', ``Erreur:
get\_inverse(1) devrait être `a'\,'' assert `b' in idict, ``Erreur: `b'
devrait être dans le dictionnaire'' assert 2 in idict.values(),
``Erreur: 2 devrait être dans les valeurs'' print(``Test 1 réussi:
Initialisation et ajout basique OK.'')

\section{Test 2: Initialisation avec des
données}\label{test-2-initialisation-avec-des-donnuxe9es}

print(``\nTest 2: Initialisation avec des données'') initial\_data =
\{`x': 10, `y': 20\} idict2 = InvertibleDict(initial\_data) assert
len(idict2) == 2, f''Erreur: Taille attendue 2, obtenue
\{len(idict2)\}'' assert idict2{[}`x'{]} == 10 assert
idict2.get\_inverse(20) == `y' print(``Test 2 réussi: Initialisation
avec des données OK.'')

\section{Test 3: Tentative d'ajout de clé dupliquée (avec valeur
différente)}\label{test-3-tentative-dajout-de-cluxe9-dupliquuxe9e-avec-valeur-diffuxe9rente}

print(``\nTest 3: Tentative d'ajout de clé dupliquée (avec valeur
différente)'') try: idict{[}`a'{]} = 4 \# `a' est déjà associé à 1
assert False, ``Erreur: KeyError attendue pour clé dupliquée avec valeur
différente'' except ValueError as e: \# Note: Notre implémentation lève
ValueError si la nouvelle valeur existe déjà assert ``valeur `4' est
déjà associée'' not in str(e) \# 4 n'existe pas encore print(f''Test 3
réussi: Exception levée comme attendu: \{e\}``) except Exception as e:
assert False, f''Erreur: Type d'exception inattendu:
\{type(e).\_\_name\_\_\}''

\section{Test 4: Tentative d'ajout de valeur dupliquée (avec clé
différente)}\label{test-4-tentative-dajout-de-valeur-dupliquuxe9e-avec-cluxe9-diffuxe9rente}

print(``\nTest 4: Tentative d'ajout de valeur dupliquée (avec clé
différente)'') try: idict{[}`d'{]} = 1 \# 1 est déjà associé à `a'
assert False, ``Erreur: ValueError attendue pour valeur dupliquée''
except ValueError as e: assert ``valeur `1' est déjà associée à la clé
`a'\,'' in str(e) print(f''Test 4 réussi: Exception levée comme attendu:
\{e\}``) except Exception as e: assert False, f''Erreur: Type
d'exception inattendu: \{type(e).\_\_name\_\_\}''

\section{Test 5: Modification d'une paire existante (clé -\textgreater{}
nouvelle valeur
unique)}\label{test-5-modification-dune-paire-existante-cluxe9---nouvelle-valeur-unique}

print(``\nTest 5: Modification d'une paire existante (clé
-\textgreater{} nouvelle valeur unique)'') idict{[}`a'{]} = 100 \# `a'
était 1, maintenant 100. 100 n'existe pas. assert len(idict) == 3,
f''Erreur: Taille attendue 3, obtenue \{len(idict)\}'' assert
idict{[}`a'{]} == 100, ``Erreur: idict{[}`a'{]} devrait être 100''
assert idict.get\_inverse(100) == `a', ``Erreur: get\_inverse(100)
devrait être `a'\,'' try: idict.get\_inverse(1) \# L'ancienne valeur 1
ne devrait plus exister assert False, ``Erreur: L'ancienne valeur 1
devrait avoir été supprimée'' except ValueError: print(``Test 5 réussi:
Modification de valeur OK, ancienne valeur supprimée.'')

\section{Test 6: Tentative de modification d'une clé vers une valeur
déjà
existante}\label{test-6-tentative-de-modification-dune-cluxe9-vers-une-valeur-duxe9juxe0-existante}

print(``\nTest 6: Tentative de modification d'une clé vers une valeur
déjà existante'') try: idict{[}`b'{]} = 100 \# 100 est déjà associé à
`a' assert False, ``Erreur: ValueError attendue pour modification vers
valeur dupliquée'' except ValueError as e: assert ``valeur `100' est
déjà associée à la clé `a'\,'' in str(e) print(f''Test 6 réussi:
Exception levée comme attendu: \{e\}``)

\section{Test 7: Suppression d'une
paire}\label{test-7-suppression-dune-paire}

print(``\nTest 7: Suppression d'une paire'') del idict{[}`b'{]} assert
len(idict) == 2, f''Erreur: Taille attendue 2 après suppression, obtenue
\{len(idict)\}'' assert `b' not in idict, ``Erreur: `b' ne devrait plus
être dans le dictionnaire'' try: idict{[}`b'{]} assert False, ``Erreur:
KeyError attendue après suppression de `b'\,'' except KeyError: pass
try: idict.get\_inverse(2) \# 2 était la valeur de `b' assert False,
``Erreur: ValueError attendue après suppression de la valeur 2'' except
ValueError: pass print(``Test 7 réussi: Suppression de paire OK.'')

\section{Test 8: Tentative de suppression d'une clé
inexistante}\label{test-8-tentative-de-suppression-dune-cluxe9-inexistante}

print(``\nTest 8: Tentative de suppression d'une clé inexistante'') try:
del idict{[}`z'{]} assert False, ``Erreur: KeyError attendue pour
suppression de clé inexistante'' except KeyError as e: assert ``clé `z'
n'existe pas'' in str(e) print(f''Test 8 réussi: Exception levée comme
attendu: \{e\}``)

\section{Test 9: Méthodes keys(), values(),
items()}\label{test-9-muxe9thodes-keys-values-items}

print(``\nTest 9: Méthodes keys(), values(), items()'') idict{[}`e'{]} =
5 idict{[}`f'{]} = 6 assert sorted(list(idict.keys())) == {[}`a', `c',
`e', `f'{]}, ``Erreur: keys() incorrect'' assert
sorted(list(idict.values())) == {[}3, 5, 6, 100{]}, ``Erreur: values()
incorrect'' assert sorted(list(idict.items())) == {[}(`a', 100), (`c',
3), (`e', 5), (`f', 6){]}, ``Erreur: items() incorrect'' print(``Test 9
réussi: keys(), values(), items() OK.'')

\section{Test 10: Représentation
string}\label{test-10-repruxe9sentation-string}

print(``\nTest 10: Représentation string'') assert str(idict) ==
``\{`a': 100, `c': 3, `e': 5, `f': 6\}'', ``Erreur: \textbf{str}
incorrect'' assert repr(idict).startswith(``InvertibleDict(''),
``Erreur: \textbf{repr} incorrect'' print(``Test 10 réussi:
Représentation string OK.'')

\section{Test 11: Initialisation avec des données dupliquées (doit
échouer)}\label{test-11-initialisation-avec-des-donnuxe9es-dupliquuxe9es-doit-uxe9chouer}

print(``\nTest 11: Initialisation avec des données dupliquées (doit
échouer)'') try: InvertibleDict(\{`k1': `v1', `k2': `v1'\}) \# v1 est
dupliquée assert False, ``Erreur: Initialisation avec valeur dupliquée
devrait échouer'' except ValueError as e: assert ``valeur `v1' est déjà
associée'' in str(e) print(f''Test 11 réussi: Initialisation avec valeur
dupliquée échoue comme attendu: \{e\}``)

try: InvertibleDict({[}(`k1', `v1'), (`k1', `v2'){]}) \# k1 est
dupliquée assert False, ``Erreur: Initialisation avec clé dupliquée
devrait échouer'' except ValueError as e: \# Notre implémentation de
\textbf{setitem} lève ValueError si la nouvelle valeur existe déjà
assert ``valeur `v2' est déjà associée'' not in str(e) \# v2 n'existe
pas encore print(f''Test 11 réussi: Initialisation avec clé dupliquée
échoue comme attendu: \{e\}``)

print(``\n--- Tous les tests pour InvertibleDict sont passés avec succès
! ---'')
\newpage
\section{TP 5}
\section{TP 5/5 - Mini-projet IA : Implémentation d'une transformation
(fonction) et de son inverse pour normaliser/dénormaliser des données en
1D}\label{tp-55---mini-projet-ia-impluxe9mentation-dune-transformation-fonction-et-de-son-inverse-pour-normaliserduxe9normaliser-des-donnuxe9es-en-1d}

\textbf{Jalon 5 : Applications, injections, surjections, bijections et
composition de fonctions}

\subsection{Objectifs pédagogiques}\label{objectifs-puxe9dagogiques}

Ce Travail Pratique vise à vous faire explorer les concepts fondamentaux
des fonctions (applications, injections, surjections, bijections) à
travers un cas d'usage concret et essentiel en Intelligence Artificielle
: la normalisation des données. Vous devrez implémenter une fonction de
normalisation et son inverse en Python pur, sans l'aide de bibliothèques
de haut niveau, et valider mathématiquement leur justesse.

À la fin de ce TP, vous devriez être capable de : * Comprendre la
nécessité et le principe de la normalisation Min-Max en IA. *
Implémenter une fonction de transformation (normalisation) et son
inverse (dénormalisation) en Python pur. * Appliquer les concepts
d'application, d'injection, de surjection et de bijection à ces
fonctions. * Valider mathématiquement la justesse de vos implémentations
à l'aide d'\passthrough{\lstinline!assert!}. * Expliquer rigoureusement
les propriétés mathématiques des fonctions implémentées.

\subsection{Introduction : Pourquoi normaliser les données en IA
?}\label{introduction-pourquoi-normaliser-les-donnuxe9es-en-ia}

Dans de nombreux algorithmes d'apprentissage automatique, la performance
et la convergence peuvent être fortement affectées par l'échelle des
caractéristiques (features). Si les caractéristiques ont des plages de
valeurs très différentes, celles avec de grandes valeurs peuvent dominer
celles avec de petites valeurs, même si ces dernières sont tout aussi
importantes.

La \textbf{normalisation} est une technique de prétraitement des données
qui consiste à redimensionner les valeurs des caractéristiques pour
qu'elles se situent dans une plage standard, souvent entre 0 et 1. Cela
permet de : 1. \textbf{Accélérer la convergence des algorithmes
d'optimisation} (comme la descente de gradient), car les pas de mise à
jour des poids seront plus équilibrés. 2. \textbf{Éviter la domination
de certaines caractéristiques} sur d'autres. 3. \textbf{Améliorer la
performance} de certains modèles (ex: SVM, K-Means, réseaux de
neurones).

Dans ce TP, nous allons nous concentrer sur la \textbf{normalisation
Min-Max}, une méthode simple mais efficace. Nous verrons comment cette
transformation est une fonction bijective et comment son inverse permet
de retrouver les données originales.

\subsection{Partie 1 : Rappels Mathématiques et Définition de la
Transformation
Min-Max}\label{partie-1-rappels-mathuxe9matiques-et-duxe9finition-de-la-transformation-min-max}

\subsubsection{1.1 La Normalisation
Min-Max}\label{la-normalisation-min-max}

La normalisation Min-Max transforme une valeur \(x\) d'une plage
originale \([\min_{original}, \max_{original}]\) vers une nouvelle plage
\([0, 1]\). La formule est la suivante :

\[ x' = f(x) = \frac{x - \min_{original}}{\max_{original} - \min_{original}} \]

Où : * \(x\) est la valeur originale. * \(\min_{original}\) est la
valeur minimale de l'ensemble de données original. * \(\max_{original}\)
est la valeur maximale de l'ensemble de données original. * \(x'\) est
la valeur normalisée, qui sera comprise entre 0 et 1.

\textbf{Conditions et Propriétés :} * Cette fonction est bien définie si
et seulement si \(\max_{original} - \min_{original} \neq 0\). Si
\(\min_{original} = \max_{original}\), toutes les valeurs sont
identiques, et la normalisation n'est pas nécessaire (ou le dénominateur
serait zéro). Dans ce cas, toutes les valeurs normalisées pourraient
être définies à 0.5 ou 0. * \textbf{Application :} Pour tout
\(x \in [\min_{original}, \max_{original}]\), \(f(x)\) est une valeur
unique dans \([0, 1]\). C'est donc bien une application (ou fonction). *
\textbf{Injectivité :} Si \(\min_{original} \neq \max_{original}\), la
fonction \(f(x)\) est strictement croissante (sa dérivée est
\(1 / (\max_{original} - \min_{original}) > 0\)). Une fonction
strictement monotone est toujours injective. Cela signifie que deux
valeurs originales différentes auront toujours deux valeurs normalisées
différentes. * \textbf{Surjectivité :} La fonction \(f(x)\) mappe
l'intervalle \([\min_{original}, \max_{original}]\) sur l'intervalle
\([0, 1]\). Pour toute valeur \(y \in [0, 1]\), il existe un
\(x \in [\min_{original}, \max_{original}]\) tel que \(f(x) = y\). Elle
est donc surjective sur son codomaine \([0, 1]\). * \textbf{Bijectivité
:} Puisqu'elle est à la fois injective et surjective (sur son codomaine
restreint), la fonction de normalisation Min-Max est une
\textbf{bijection} de \([\min_{original}, \max_{original}]\) vers
\([0, 1]\) (si \(\min_{original} \neq \max_{original}\)).

\subsubsection{1.2 La Dénormalisation Min-Max (Fonction
Inverse)}\label{la-duxe9normalisation-min-max-fonction-inverse}

Puisque la normalisation Min-Max est une bijection, elle possède une
fonction inverse unique. Cette fonction inverse est appelée
\textbf{dénormalisation}. Elle permet de retrouver la valeur originale
\(x\) à partir de sa valeur normalisée \(x'\).

Pour trouver l'inverse, nous résolvons l'équation de normalisation pour
\(x\):
\[ x' = \frac{x - \min_{original}}{\max_{original} - \min_{original}} \]
\[ x' (\max_{original} - \min_{original}) = x - \min_{original} \]
\[ x = x' (\max_{original} - \min_{original}) + \min_{original} \]

Donc, la fonction de dénormalisation est :

\[ x = f^{-1}(x') = x' (\max_{original} - \min_{original}) + \min_{original} \]

Où : * \(x'\) est la valeur normalisée. * \(\min_{original}\) et
\(\max_{original}\) sont les valeurs min et max de l'ensemble de données
original. * \(x\) est la valeur dénormalisée, qui devrait correspondre à
la valeur originale.

\textbf{Propriétés de l'inverse :} * La fonction \(f^{-1}(x')\) est
également une bijection de \([0, 1]\) vers
\([\min_{original}, \max_{original}]\). * La composition des fonctions
doit donner l'identité : \(f^{-1}(f(x)) = x\) et \(f(f^{-1}(x')) = x'\).
C'est cette propriété que nous validerons avec des
\passthrough{\lstinline!assert!}.

\subsection{Partie 2 : Implémentation en Python
Pur}\label{partie-2-impluxe9mentation-en-python-pur}

Nous allons implémenter les fonctions nécessaires en Python pur, sans
utiliser de bibliothèques comme \passthrough{\lstinline!numpy!} ou
\passthrough{\lstinline!scikit-learn!}.

\begin{lstlisting}[language=Python]
import math # Pour la tolerance des flottants, si necessaire, mais on peut faire sans pour ce TP.

def find_min_max(data):
    """
    Trouve les valeurs minimale et maximale dans une liste de donnees.
    Implementation manuelle pour respecter la contrainte "Python PUR from scratch".

    Args:
        data (list): Une liste de nombres (int ou float).

    Returns:
        tuple: Un tuple (min_val, max_val) contenant la valeur minimale et maximale.
               Retourne (None, None) si la liste est vide.
    """
    if not data:
        return None, None

    min_val = data[0]
    max_val = data[0]

    for value in data:
        if value < min_val:
            min_val = value
        if value > max_val:
            max_val = value

    return min_val, max_val

def normalize_min_max(value, data_min, data_max):
    """
    Normalise une seule valeur en utilisant la methode Min-Max.
    La valeur est transformee pour etre dans l'intervalle [0, 1].

    Args:
        value (float or int): La valeur a normaliser.
        data_min (float or int): La valeur minimale de l'ensemble de donnees original.
        data_max (float or int): La valeur maximale de l'ensemble de donnees original.

    Returns:
        float: La valeur normalisee.
               Retourne 0.5 si data_min == data_max pour eviter la division par zero,
               indiquant que toutes les valeurs etaient identiques.
    """
    if data_max == data_min:
        # Si toutes les valeurs sont identiques, la plage est nulle.
        # Par convention, on peut retourner 0.5 (milieu de [0,1]) ou 0.
        # Ici, 0.5 est choisi pour representer une "valeur moyenne" dans la plage normalisee.
        return 0.5

    # Applique la formule de normalisation Min-Max
    normalized_value = (value - data_min) / (data_max - data_min)
    return normalized_value

def denormalize_min_max(normalized_value, data_min, data_max):
    """
    Denormalise une seule valeur (precedemment normalisee Min-Max) pour retrouver
    sa plage originale.

    Args:
        normalized_value (float): La valeur normalisee (attendue dans [0, 1]).
        data_min (float or int): La valeur minimale de l'ensemble de donnees original.
        data_max (float or int): La valeur maximale de l'ensemble de donnees original.

    Returns:
        float: La valeur denormalisee, correspondant a la valeur originale.
               Retourne data_min si data_min == data_max, car toutes les valeurs
               originales etaient egales a data_min.
    """
    if data_max == data_min:
        # Si toutes les valeurs etaient identiques, la denormalisation ramene a cette valeur unique.
        return data_min

    # Applique la formule de denormalisation (fonction inverse)
    original_value = normalized_value * (data_max - data_min) + data_min
    return original_value

def apply_normalization(data):
    """
    Applique la normalisation Min-Max a une liste entiere de donnees.

    Args:
        data (list): La liste des donnees a normaliser.

    Returns:
        tuple: Un tuple contenant (normalized_data, original_min, original_max).
               normalized_data est la liste des valeurs normalisees.
               original_min et original_max sont les min/max de la liste originale,
               necessaires pour la denormalisation.
    """
    if not data:
        return [], None, None

    original_min, original_max = find_min_max(data)

    # Gerer le cas ou toutes les valeurs sont identiques
    if original_min == original_max:
        # Si toutes les valeurs sont les memes, toutes les valeurs normalisees seront 0.5
        # (selon notre convention dans normalize_min_max).
        # On retourne une liste de 0.5 de la meme taille que les donnees originales.
        return [0.5] * len(data), original_min, original_max

    normalized_data = []
    for value in data:
        normalized_data.append(normalize_min_max(value, original_min, original_max))

    return normalized_data, original_min, original_max

def apply_denormalization(normalized_data, original_min, original_max):
    """
    Applique la denormalisation Min-Max a une liste de donnees normalisees.

    Args:
        normalized_data (list): La liste des valeurs normalisees.
        original_min (float or int): La valeur minimale de l'ensemble de donnees original.
        original_max (float or int): La valeur maximale de l'ensemble de donnees original.

    Returns:
        list: La liste des valeurs denormalisees.
    """
    if not normalized_data:
        return []

    denormalized_data = []
    for value in normalized_data:
        denormalized_data.append(denormalize_min_max(value, original_min, original_max))

    return denormalized_data
\end{lstlisting}

\subsection{\texorpdfstring{Partie 3 : Validation Mathématique et Tests
avec
\texttt{assert}}{Partie 3 : Validation Mathématique et Tests avec assert}}\label{partie-3-validation-mathuxe9matique-et-tests-avec-assert}

Nous allons maintenant tester nos fonctions et utiliser des
\passthrough{\lstinline!assert!} pour valider leur comportement et leurs
propriétés mathématiques, notamment la bijectivité et la relation
inverse.

\begin{lstlisting}[language=Python]
# --- Jeu de donnees de test ---
data_set_1 = [10, 20, 30, 40, 50]
data_set_2 = [-5, 0, 5, 10, 15.5, 20]
data_set_3 = [7.2, 1.1, 9.8, 3.5, 6.0]
data_set_4_single_value = [100] # Cas ou min == max
data_set_5_empty = [] # Cas de liste vide
data_set_6_negative = [-10, -5, -1, -0.5, 0]

# --- Tolerance pour la comparaison des flottants ---
# Les operations en virgule flottante peuvent introduire de petites erreurs.
# Nous utilisons une petite tolerance (epsilon) pour comparer les resultats.
EPSILON = 1e-9

print("--- Debut des tests de validation ---")

# Test 1: find_min_max
print("\nTest 1: find_min_max")
min1, max1 = find_min_max(data_set_1)
assert min1 == 10 and max1 == 50, f"Test 1.1 failed: Expected (10, 50), got ({min1}, {max1})"
print(f"  Min/Max pour {data_set_1}: ({min1}, {max1})")

min2, max2 = find_min_max(data_set_2)
assert min2 == -5 and max2 == 20, f"Test 1.2 failed: Expected (-5, 20), got ({min2}, {max2})"
print(f"  Min/Max pour {data_set_2}: ({min2}, {max2})")

min4, max4 = find_min_max(data_set_4_single_value)
assert min4 == 100 and max4 == 100, f"Test 1.3 failed: Expected (100, 100), got ({min4}, {max4})"
print(f"  Min/Max pour {data_set_4_single_value}: ({min4}, {max4})")

min5, max5 = find_min_max(data_set_5_empty)
assert min5 is None and max5 is None, f"Test 1.4 failed: Expected (None, None), got ({min5}, {max5})"
print(f"  Min/Max pour {data_set_5_empty}: ({min5}, {max5})")
print("  find_min_max: OK")


# Test 2: Normalisation et Denormalisation d'une seule valeur
print("\nTest 2: Normalisation/Denormalisation d'une seule valeur")
original_min_ds1, original_max_ds1 = find_min_max(data_set_1) # (10, 50)

# Normalisation d'une valeur au milieu
val_to_normalize = 30
normalized_val = normalize_min_max(val_to_normalize, original_min_ds1, original_max_ds1)
assert abs(normalized_val - 0.5) < EPSILON, f"Test 2.1 failed: Expected 0.5, got {normalized_val}"
print(f"  {val_to_normalize} normalise: {normalized_val}")

# Normalisation de la valeur min
normalized_min = normalize_min_max(original_min_ds1, original_min_ds1, original_max_ds1)
assert abs(normalized_min - 0.0) < EPSILON, f"Test 2.2 failed: Expected 0.0, got {normalized_min}"
print(f"  {original_min_ds1} normalise: {normalized_min}")

# Normalisation de la valeur max
normalized_max = normalize_min_max(original_max_ds1, original_min_ds1, original_max_ds1)
assert abs(normalized_max - 1.0) < EPSILON, f"Test 2.3 failed: Expected 1.0, got {normalized_max}"
print(f"  {original_max_ds1} normalise: {normalized_max}")

# Denormalisation de la valeur normalisee
denormalized_val = denormalize_min_max(normalized_val, original_min_ds1, original_max_ds1)
assert abs(denormalized_val - val_to_normalize) < EPSILON, f"Test 2.4 failed: Expected {val_to_normalize}, got {denormalized_val}"
print(f"  {normalized_val} denormalise: {denormalized_val}")

# Test de la propriete inverse f^-1(f(x)) = x
# On prend une valeur arbitraire de data_set_1
test_val_inverse = 25
normalized_test_val = normalize_min_max(test_val_inverse, original_min_ds1, original_max_ds1)
denormalized_test_val = denormalize_min_max(normalized_test_val, original_min_ds1, original_max_ds1)
assert abs(denormalized_test_val - test_val_inverse) < EPSILON, \
    f"Test 2.5 failed (f^-1(f(x)) = x): Expected {test_val_inverse}, got {denormalized_test_val}"
print(f"  f^-1(f({test_val_inverse})) = {denormalized_test_val} (attendu {test_val_inverse})")

# Test du cas min == max pour une seule valeur
original_min_ds4, original_max_ds4 = find_min_max(data_set_4_single_value) # (100, 100)
normalized_val_ds4 = normalize_min_max(100, original_min_ds4, original_max_ds4)
assert abs(normalized_val_ds4 - 0.5) < EPSILON, f"Test 2.6 failed: Expected 0.5 for single value, got {normalized_val_ds4}"
print(f"  {data_set_4_single_value[0]} normalise (min=max): {normalized_val_ds4}")
denormalized_val_ds4 = denormalize_min_max(normalized_val_ds4, original_min_ds4, original_max_ds4)
assert abs(denormalized_val_ds4 - 100) < EPSILON, f"Test 2.7 failed: Expected 100 for single value denormalization, got {denormalized_val_ds4}"
print(f"  {normalized_val_ds4} denormalise (min=max): {denormalized_val_ds4}")
print("  Normalisation/Denormalisation d'une seule valeur: OK")


# Test 3: Normalisation et Denormalisation de listes completes
print("\nTest 3: Normalisation/Denormalisation de listes completes")

# Test avec data_set_1
normalized_data_1, min_orig_1, max_orig_1 = apply_normalization(data_set_1)
print(f"  Original: {data_set_1}")
print(f"  Normalise: {normalized_data_1}")
# Verifier que toutes les valeurs normalisees sont dans [0, 1]
for val in normalized_data_1:
    assert 0.0 - EPSILON <= val <= 1.0 + EPSILON, f"Test 3.1 failed: Normalized value {val} out of [0, 1] range."
# Verifier les valeurs specifiques
assert abs(normalized_data_1[0] - 0.0) < EPSILON, "Test 3.2 failed: First value not 0.0"
assert abs(normalized_data_1[2] - 0.5) < EPSILON, "Test 3.3 failed: Middle value not 0.5"
assert abs(normalized_data_1[-1] - 1.0) < EPSILON, "Test 3.4 failed: Last value not 1.0"

denormalized_data_1 = apply_denormalization(normalized_data_1, min_orig_1, max_orig_1)
print(f"  Denormalise: {denormalized_data_1}")
# Verifier que les donnees denormalisees sont (presque) identiques aux originales
for i in range(len(data_set_1)):
    assert abs(denormalized_data_1[i] - data_set_1[i]) < EPSILON, \
        f"Test 3.5 failed: Denormalisation incorrecte pour l'element {i}. Attendu {data_set_1[i]}, obtenu {denormalized_data_1[i]}"
print("  Normalisation/Denormalisation de data_set_1: OK")

# Test avec data_set_3 (flottants)
normalized_data_3, min_orig_3, max_orig_3 = apply_normalization(data_set_3)
print(f"\n  Original: {data_set_3}")
print(f"  Normalise: {normalized_data_3}")
for val in normalized_data_3:
    assert 0.0 - EPSILON <= val <= 1.0 + EPSILON, f"Test 3.6 failed: Normalized value {val} out of [0, 1] range."

denormalized_data_3 = apply_denormalization(normalized_data_3, min_orig_3, max_orig_3)
print(f"  Denormalise: {denormalized_data_3}")
for i in range(len(data_set_3)):
    assert abs(denormalized_data_3[i] - data_set_3[i]) < EPSILON, \
        f"Test 3.7 failed: Denormalisation incorrecte pour l'element {i}. Attendu {data_set_3[i]}, obtenu {denormalized_data_3[i]}"
print("  Normalisation/Denormalisation de data_set_3: OK")

# Test avec data_set_4_single_value (min == max)
normalized_data_4, min_orig_4, max_orig_4 = apply_normalization(data_set_4_single_value)
print(f"\n  Original: {data_set_4_single_value}")
print(f"  Normalise: {normalized_data_4}")
assert len(normalized_data_4) == 1 and abs(normalized_data_4[0] - 0.5) < EPSILON, \
    f"Test 3.8 failed: Expected [0.5] for single value list, got {normalized_data_4}"

denormalized_data_4 = apply_denormalization(normalized_data_4, min_orig_4, max_orig_4)
print(f"  Denormalise: {denormalized_data_4}")
assert len(denormalized_data_4) == 1 and abs(denormalized_data_4[0] - data_set_4_single_value[0]) < EPSILON, \
    f"Test 3.9 failed: Expected {data_set_4_single_value} for single value list denormalization, got {denormalized_data_4}"
print("  Normalisation/Denormalisation de data_set_4_single_value: OK")

# Test avec data_set_5_empty (liste vide)
normalized_data_5, min_orig_5, max_orig_5 = apply_normalization(data_set_5_empty)
print(f"\n  Original: {data_set_5_empty}")
print(f"  Normalise: {normalized_data_5}")
assert normalized_data_5 == [] and min_orig_5 is None and max_orig_5 is None, \
    f"Test 3.10 failed: Expected ([], None, None) for empty list, got ({normalized_data_5}, {min_orig_5}, {max_orig_5})"
denormalized_data_5 = apply_denormalization([], None, None) # Simuler la denormalisation d'une liste vide
assert denormalized_data_5 == [], f"Test 3.11 failed: Expected [] for empty list denormalization, got {denormalized_data_5}"
print("  Normalisation/Denormalisation de data_set_5_empty: OK")

print("\n--- Tous les tests de validation ont reussi ! ---")
\end{lstlisting}

\subsection{Conclusion et Réflexions}\label{conclusion-et-ruxe9flexions}

Ce TP vous a permis d'implémenter une transformation de données
fondamentale en Intelligence Artificielle : la normalisation Min-Max. En
le faisant en Python pur, vous avez dû comprendre et gérer les détails
de l'algorithme, y compris les cas limites (comme
\passthrough{\lstinline!min == max!}).

Nous avons rigoureusement démontré et validé que la fonction de
normalisation Min-Max, sous la condition
\(\min_{original} \neq \max_{original}\), est une \textbf{bijection} de
l'intervalle \([\min_{original}, \max_{original}]\) vers l'intervalle
\([0, 1]\). Sa fonction inverse, la dénormalisation, est également une
bijection et permet de retrouver les données originales sans perte
d'information. Cette propriété de bijectivité est cruciale pour toute
transformation réversible en traitement de données.

\textbf{Points à retenir :} * La normalisation est essentielle pour de
nombreux algorithmes d'IA. * La normalisation Min-Max est une bijection,
ce qui garantit l'existence et l'unicité de son inverse. * La capacité à
dénormaliser est vitale pour interpréter les résultats des modèles d'IA
dans l'échelle originale des données. * Les
\passthrough{\lstinline!assert!} sont des outils puissants pour valider
mathématiquement et logiquement le comportement de votre code.

\textbf{Pour aller plus loin :} * Explorez d'autres méthodes de
normalisation, comme la \textbf{standardisation (Z-score)} :
\(x' = (x - \mu) / \sigma\), où \(\mu\) est la moyenne et \(\sigma\)
l'écart-type. Discutez de ses propriétés (injection, surjection,
bijection) et de son inverse. * Considérez l'impact des valeurs
aberrantes (outliers) sur la normalisation Min-Max. Comment une seule
valeur extrême peut-elle ``écraser'' toutes les autres valeurs dans une
petite plage ? * Implémentez une version de normalisation robuste qui
utilise les quartiles (IQR) au lieu du min/max pour être moins sensible
aux outliers.
\newpage

\chapter{Annexes: Représentations Visuelles (TikZ)}
\section{Schéma d'une Application}
\begin{center}
\begin{tikzpicture}
    \draw (0,0) ellipse (1cm and 2cm);
    \draw (4,0) ellipse (1cm and 2cm);
    \node at (0,2.3) {$E$};
    \node at (4,2.3) {$F$};
    \node (x1) at (0,1) {$\bullet$}; \node at (-0.3,1) {$x_1$};
    \node (x2) at (0,0) {$\bullet$}; \node at (-0.3,0) {$x_2$};
    \node (x3) at (0,-1) {$\bullet$}; \node at (-0.3,-1) {$x_3$};
    \node (y1) at (4,1.5) {$\bullet$}; \node at (4.3,1.5) {$y_1$};
    \node (y2) at (4,0.5) {$\bullet$}; \node at (4.3,0.5) {$y_2$};
    \node (y3) at (4,-0.5) {$\bullet$}; \node at (4.3,-0.5) {$y_3$};
    \node (y4) at (4,-1.5) {$\bullet$}; \node at (4.3,-1.5) {$y_4$};
    \draw[->] (x1) -- (y1);
    \draw[->] (x2) -- (y3);
    \draw[->] (x3) -- (y3);
\end{tikzpicture}
\end{center}

\section{Schéma d'une Injection}
\begin{center}
\begin{tikzpicture}
    \draw (0,0) ellipse (1cm and 2cm);
    \draw (4,0) ellipse (1cm and 2cm);
    \node at (0,2.3) {$E$};
    \node at (4,2.3) {$F$};
    \node (x1) at (0,1) {$\bullet$}; \node at (-0.3,1) {$x_1$};
    \node (x2) at (0,0) {$\bullet$}; \node at (-0.3,0) {$x_2$};
    \node (x3) at (0,-1) {$\bullet$}; \node at (-0.3,-1) {$x_3$};
    \node (y1) at (4,1.5) {$\bullet$}; \node at (4.3,1.5) {$y_1$};
    \node (y2) at (4,0.5) {$\bullet$}; \node at (4.3,0.5) {$y_2$};
    \node (y3) at (4,-0.5) {$\bullet$}; \node at (4.3,-0.5) {$y_3$};
    \node (y4) at (4,-1.5) {$\bullet$}; \node at (4.3,-1.5) {$y_4$};
    \draw[->] (x1) -- (y1);
    \draw[->] (x2) -- (y3);
    \draw[->] (x3) -- (y2);
\end{tikzpicture}
\end{center}
\section{Schéma d'une Surjection}
\begin{center}
\begin{tikzpicture}
    \draw (0,0) ellipse (1cm and 2cm);
    \draw (4,0) ellipse (1cm and 1.5cm);
    \node at (0,2.3) {$E$};
    \node at (4,1.8) {$F$};
    \node (x1) at (0,1.5) {$\bullet$}; \node at (-0.3,1.5) {$x_1$};
    \node (x2) at (0,0.5) {$\bullet$}; \node at (-0.3,0.5) {$x_2$};
    \node (x3) at (0,-0.5) {$\bullet$}; \node at (-0.3,-0.5) {$x_3$};
    \node (x4) at (0,-1.5) {$\bullet$}; \node at (-0.3,-1.5) {$x_4$};
    \node (y1) at (4,1) {$\bullet$}; \node at (4.3,1) {$y_1$};
    \node (y2) at (4,0) {$\bullet$}; \node at (4.3,0) {$y_2$};
    \node (y3) at (4,-1) {$\bullet$}; \node at (4.3,-1) {$y_3$};
    \draw[->] (x1) -- (y1);
    \draw[->] (x2) -- (y2);
    \draw[->] (x3) -- (y2);
    \draw[->] (x4) -- (y3);
\end{tikzpicture}
\end{center}
\section{Schéma d'une Bijection}
\begin{center}
\begin{tikzpicture}
    \draw (0,0) ellipse (1cm and 1.5cm);
    \draw (4,0) ellipse (1cm and 1.5cm);
    \node at (0,1.8) {$E$};
    \node at (4,1.8) {$F$};
    \node (x1) at (0,1) {$\bullet$}; \node at (-0.3,1) {$x_1$};
    \node (x2) at (0,0) {$\bullet$}; \node at (-0.3,0) {$x_2$};
    \node (x3) at (0,-1) {$\bullet$}; \node at (-0.3,-1) {$x_3$};
    \node (y1) at (4,1) {$\bullet$}; \node at (4.3,1) {$y_1$};
    \node (y2) at (4,0) {$\bullet$}; \node at (4.3,0) {$y_2$};
    \node (y3) at (4,-1) {$\bullet$}; \node at (4.3,-1) {$y_3$};
    \draw[->] (x1) -- (y2);
    \draw[->] (x2) -- (y1);
    \draw[->] (x3) -- (y3);
\end{tikzpicture}
\end{center}
\end{document}
